%=====================================================================================
%  CHAPITRE 4 -- Cadre de modelisation et strategie de resolution numerique
%  Version complete et redigee, avec integration des notes
%  "Derivation of the Navier-Stokes equations, and implications for groundwater flow"
%
%  FICHIERS GRAPHIQUES ATTENDUS dans le dossier du chapitre :
%    figure1.png (schema du conduit)  figure2.png (boucle de couplage)
%    figure3.png (profils spatiaux)   figure4.png (convergence de maillage)
%  Retirer les deux environnements "vigilance" avant remise du manuscrit.
%
%  -----------------------------------------------------------------------------------
%  A. INTEGRATION DES NOTES (5 insertions)
%   I1  Bilan de masse 3D (div v = 0) -> contrainte globale d(Au)/dx = 0    [section 4.2.2]
%   I2  Hierarchie des echelles et justification du quasi-statique            [section 4.3.4]
%   I3  Reduction dimensionnelle du modele de Navier-Stokes                   [section 4.3]
%   I4  Poiseuille / forme darcienne comme archetype de fermeture locale      [section 4.3.3]
%   I5  Tableau des nombres sans dimension et criteres de validite            [section 4.3.5]
%   Annexe A : chaine d'etablissement Cauchy -> Navier-Stokes -> Poiseuille -> Darcy
%
%  B. CORRECTIONS APPORTEES AU CHAPITRE D'ORIGINE
%   T1  "equation de pression elliptique / methode de Poisson" -> quadrature + contrainte
%       globale + solveur scalaire (4.3.2, 4.6.3, 4.7.1). Aucun Laplacien n'est resolu.
%   T2  Ancienne section 2.4 (redondante) condensee : section 4.6 (organisation de
%       l'algorithme) ; l'argumentaire de robustesse constitue la section 4.7.5.
%   T3  "la dichotomie conserve l'encadrement comme Newton" -> reformule (4.7.5).
%   T4  coefficient beta_u (vitesse de transport des particules) defini explicitement.
%   T5  "globalement implicite" -> "hydraulique implicite + couplage sequentiel decale".
%   T6  Nombre de Reynolds : rho du melange, coherent avec la fermeture de frottement.
%   T7  phi_soil et parametres : tableau 4.5 par campagne (plus de valeur unique ambigue).
%   T8  Limiteur : seuil de stabilite explicite (CFL <= 1/2), choix documente.
%   T9  "1D inconditionnellement stable" -> stabilite conditionnelle, seuil quantifie.
%   T10 Rapport d'injection : ecriture dimensionnellement correcte (section 4.2.2).
%
%  C. RESULTATS DE VERIFICATION AJOUTES AU CHAPITRE (tous reproductibles)
%   V1  tau_b = R dP/(2L) : verifie a 1e-6 pres pour R/R0 = 1 a 10 (section 4.7.6).
%   V2  Q proportionnel a R^(19/7) : pente log 2,714 (mesuree 2,714) (section 4.7.6).
%   V3  t2 = t_er ln 2 = 658 s contre 666 s simule -> accord a 1 % (section 4.7.6).
%   V4  Seuil du schema d'advection : CFL <= 1/2 sans predicteur, CFL <= 1 avec
%       predicteur MUSCL-Hancock (section 4.8.4) -> justifie le reglage CFL = 0,35.
%   V5  Loi asymptotique des erreurs sur un front capture : L1 ~ dx,
%       L2 ~ dx^(1/2), L infini = constante = saut/8 (section 4.10.6).
%   V6  Sensibilite a K_out : tableau 4.11 (t2 = 666 s pour K=0, 2,1e4 s pour K=10).
%   V7  Loi f_m : le plafond est applique DANS le solveur. Non-regression de la campagne A
%       (trois traitements du plafond -> t2 = 668,5 s identiques, f_m = 1,0001).
%   V8  Sans plafond, la campagne B atteint f_m = 1,46e12, R_max = 179 R0 pour R_out = 3,5 R0
%       et s'arrete sur R_break (artefact). Avec plafond (fm_max = 2000) : runs reguliers.
%   V9  Validation croisee campagne B (Tab. 5 publie) : t_fail = 128,60 / 162,13 / 55,26 / 63,22 s
%       contre 128,1 / 161,8 / 55,0 / 62,9 s publies -> ecart < 0,5 % sur 4 configurations ;
%       facteurs de retard par K_out = 500 : x1,261 et x1,144 contre x1,263 et x1,144 publies.
%   V10 Campagne A : 668 s contre 666 s publies (+0,4 %) ; 877,8 s contre 872 s (+0,7 %) pour la
%       configuration avec resistance aval, qui correspond a K_out ~ 0,15 et NON a K_out = 10.
%   V11 Perte d'entree K_in : t2 = 668 / 1937 / 2522 / 3096 / 4235 s pour K_in,tot = 0 / 1 / 1,5 / 2 / 3 ;
%       tau_b0 = R0[Pin - K_in,tot rho u0^2/2]/(2L) verifie a 0,3 % (section 4.10, tab. 4.13).
%
%  D. POINTS DE VIGILANCE (encadres "vigilance" dans le texte)
%   0. RESOLUS (ex-alertes 1 et 2) :
%      - l'alerte sur t2 = 872 s pour K_out = 10 venait d'une confusion entre campagnes :
%        les K_out eleves (10, 500) decrivent la campagne B (terrain), pas la campagne A.
%        Campagne A avec resistance aval <=> K_out ~ 0,15 (mesure : 877,8 s contre 872 publies).
%      - l'alerte sur le "maintien de l'emballement" malgre K_out = 500 venait de la loi f_m
%        non plafonnee dans le solveur. Avec le plafond, le calcul atteint R_break en 162,1 s
%        et la perte aval est auto-limitante (elle decroit en R^-4) : elle ne peut que retarder.
%   1. Reste ouvert : la valeur de fm_max gouverne qualitativement la campagne B (2000 -> accord
%      avec la publication ; 5 -> arret de l'erosion). A presenter comme parametre identifie.
%   2. Autres points (parametres, unites de k_er, delta_t_max de la campagne B,
%      d_p de la campagne B, criteres d'arret, limiteur unique) : voir
%      docs/integration_notes_NS_chapitre.md, section 7.
%=====================================================================================

\documentclass[11pt,a4paper]{report}
\usepackage[utf8]{inputenc}
\usepackage[T1]{fontenc}
\usepackage{lmodern}
\usepackage[french]{babel}
\usepackage{amsmath,amssymb}
\usepackage{graphicx}
\usepackage{booktabs}
\usepackage{geometry}
\usepackage{siunitx}
\usepackage{hyperref}
\usepackage{caption}
\usepackage{float}
\usepackage{array}
\usepackage{enumitem}
\usepackage{threeparttable}
\geometry{margin=1in}
\captionsetup{font=small}
\hypersetup{colorlinks=true,linkcolor=blue,citecolor=blue,urlcolor=blue}

\setlength{\parskip}{0.4em}
\setlength{\parindent}{0pt}

% ---- Macros de notation ----
\newcommand{\phisol}{\phi_{\mathrm{soil}}}
\newcommand{\kout}{K_{\mathrm{out}}}
\newcommand{\kin}{K_{\mathrm{in}}}
\newcommand{\mdot}{\dot{m}}
\newcommand{\ter}{{t_{er}}}
\newcommand{\dt}{\Delta t}
\newcommand{\Chat}{\hat{P}}
\newcommand{\divg}{\nabla\!\cdot\!}
\newcommand{\vx}{v_x}
\newcommand{\Disc}{D_{\mathrm{num}}}
% Encadre de vigilance : a retirer avant remise definitive du manuscrit
\newenvironment{vigilance}{%
  \begin{quote}\small\noindent\rule{\linewidth}{0.4pt}\par\vspace{3pt}%
  \textbf{Note de vigilance --- \`a v\'erifier avant publication.}\par\vspace{2pt}}%
  {\par\vspace{3pt}\noindent\rule{\linewidth}{0.4pt}\end{quote}}

\begin{document}

%=====================================================================================
\chapter{Cadre de modélisation et stratégie de résolution numérique}
\label{chap:cadre}
%=====================================================================================

\section{Positionnement du modèle d'ordre réduit}
\label{sec:positionnement}

L'érosion interne par renard hydraulique constitue un problème d'interaction complexe entre les
sollicitations hydrauliques, l'érosion progressive des parois du conduit, le transport des particules
détachées et l'évolution de la géométrie de la zone érodée. La représentation numérique de cette
dynamique peut être abordée à différents niveaux de fidélité, depuis les modèles analytiques ou réduits
jusqu'aux simulations multiphysiques tridimensionnelles à haute résolution. Dans ce contexte, le présent
travail s'inscrit dans une démarche de modélisation d'ordre réduit, dont l'objectif n'est pas de
reproduire l'ensemble des phénomènes tridimensionnels locaux, mais de représenter de manière cohérente
les mécanismes dominants qui gouvernent l'évolution temporelle d'un conduit de renard hydraulique. Le
modèle proposé vise ainsi à fournir un cadre de calcul suffisamment représentatif des interactions
hydrauliques et érosives tout en conservant une complexité numérique compatible avec des simulations
répétées.

Cette orientation répond directement à la finalité recherchée : disposer d'un outil permettant
d'explorer rapidement l'influence des paramètres du problème, d'examiner différents scénarios de
conditions aux limites et d'identifier les régimes d'évolution susceptibles de conduire à une
amplification de l'érosion. Les simulations tridimensionnelles de type CFD avec frontières mobiles,
ainsi que les approches sans maillage telles que la méthode SPH, permettent certes une description plus
détaillée des mécanismes locaux, mais leur coût de calcul limite leur utilisation dans des campagnes
comportant de nombreux cas paramétriques. Le modèle d'ordre réduit constitue ainsi un niveau
intermédiaire entre la description expérimentale du phénomène et les approches numériques de
haute-fidélité : il conserve les rétroactions essentielles du système tout en offrant l'efficacité
nécessaire aux études de criblage paramétrique, à l'analyse de sensibilité et à l'interprétation rapide
des régimes transitoires.

Dans cette perspective, le modèle développé s'inscrit dans la continuité des approches
unidimensionnelles antérieures consacrées à la modélisation du piping, notamment les travaux de
Lachouette et al.\ (2008, 2009), tout en introduisant une formulation spécifiquement orientée vers la
résolution numérique sous pressions d'entrée et de sortie imposées, la prise en compte de la résistance
du mélange dépendante de la concentration et une procédure explicite de vérification numérique.

Une remarque préliminaire est nécessaire sur la nature de la réduction effectuée. Le caractère
\og unidimensionnel \fg{} du modèle ne doit pas être confondu avec l'hypothèse d'un écoulement
stationnaire, ni avec celle d'un problème sans inertie. La réduction opérée consiste à \emph{moyenner}
les équations de la mécanique des fluides sur la section du conduit, puis à exploiter la séparation
entre l'échelle de temps hydraulique et l'échelle de temps d'érosion. Le présent chapitre établit
explicitement cette chaîne de réductions à la section~\ref{sec:reduction-ns} : bilan de masse
tridimensionnel, bilan de quantité de mouvement, fermeture constitutive, puis dégénérescence
quasi-statique. C'est à ce prix que les hypothèses du modèle deviennent vérifiables et que ses limites
peuvent être quantifiées plutôt que postulées.

%=====================================================================================
\section{Hypothèses de réduction unidimensionnelle}
\label{sec:hypotheses}
%=====================================================================================

La réduction du problème tridimensionnel à une formulation unidimensionnelle repose sur l'hypothèse
fondamentale selon laquelle la dynamique de l'écoulement et du transport est principalement gouvernée
par la direction longitudinale du conduit. Le domaine physique est ainsi représenté par un conduit
préférentiel de longueur \(L\), décrit par une coordonnée axiale \(x \in [0, L]\), tandis que sa
géométrie est caractérisée, à chaque position et à chaque instant, par un rayon effectif \(R(x,t)\). La
section transversale correspondante est définie par
\[
A(x,t) = \pi R^2(x,t)
\]
et le débit est relié à la vitesse moyenne sur la section par
\[
Q(t) = A(x,t)\,u(x,t).
\]

Les principales variables du modèle sont ainsi le rayon \(R(x,t)\), la pression \(p(x,t)\), la vitesse
moyenne \(u(x,t)\) et la fraction volumique de solides en suspension \(\phi(x,t)\). Cette représentation
suppose explicitement que les variables hydrauliques et de transport peuvent être décrites par des
grandeurs moyennes sur chaque section, sans résolution directe de leur distribution radiale. Cette
approximation implique plusieurs hypothèses structurantes. Premièrement, une seule valeur effective du
rayon est attribuée à chaque abscisse longitudinale. La section du conduit est donc ramenée à une
géométrie circulaire équivalente, ce qui permet d'exprimer directement les grandeurs hydrauliques en
fonction de \(R\). Deuxièmement, l'écoulement et la concentration sont considérés sous leur forme
moyennée sur la section, de sorte que les variations transversales ne sont pas résolues explicitement.
Troisièmement, la dynamique dominante est supposée longitudinale, les phénomènes bidimensionnels ou
tridimensionnels étant volontairement exclus de la formulation. Enfin, le mélange eau--particules est
traité dans le cadre d'une approximation de suspension diluée, avec une masse volumique effective
définie par
\begin{equation}
\rho(\phi) = \rho_w + (\rho_p - \rho_w)\phi,
\label{eq:rho-mix}
\end{equation}
où \(\rho_w\) et \(\rho_p\) désignent respectivement les masses volumiques de l'eau et des particules.

\subsection{Validité géométrique de la réduction}

La validité de cette réduction est liée à l'élancement du conduit. Dans le modèle, l'approximation
unidimensionnelle est considérée acceptable lorsque le diamètre demeure faible devant la longueur, avec
la condition
\begin{equation}
\frac{D}{L} < 0.1,
\qquad\text{soit de manière équivalente}\qquad
\frac{R}{L} < 0.05 .
\label{eq:validite-1D}
\end{equation}
Cette condition traduit le caractère quasi-filiforme du conduit et limite l'importance des effets
transversaux qui ne sont pas représentés dans le modèle. Lorsque le rayon augmente au point de dépasser
cette limite, les effets d'entrée, les zones de séparation et plus généralement les mécanismes
bidimensionnels deviennent susceptibles de modifier significativement la structure de l'écoulement. Dans
cette situation, les lois de fermeture hydrauliques adoptées dans la formulation 1D ne peuvent plus être
considérées comme quantitativement représentatives. Les résultats obtenus au-delà de cette limite
géométrique, notamment dans les régimes de croissance rapide du conduit, doivent par conséquent être
interprétés comme des indicateurs qualitatifs de l'instabilité plutôt que comme une description
quantitative de la phase terminale du phénomène.

Il est essentiel de souligner que cette limite s'exprime en temps de simulation par un \emph{horizon}
très concret. Pour la géométrie de référence de la campagne A (\(L = 0.117\) m, \(R_0 = 3\) mm), la
condition~\eqref{eq:validite-1D} est vérifiée pour \(R < 5.85\) mm, soit \(R/R_0 < 1.95\) : la limite de
validité quantitative est donc atteinte lorsque le rayon a simplement doublé. Toute grandeur extraite
d'un calcul au-delà de cet horizon (rapports d'agrandissement, profils spatiaux, temps d'échec) doit
être présentée comme un indicateur de régime et non comme une prédiction. Cette contrainte est reprise
quantitativement à la section~\ref{sec:validite}.

\subsection{Continuité hydraulique et bilan de masse}
\label{sec:bilan-masse}

Une seconde hypothèse importante concerne la continuité hydraulique. Bien que l'érosion constitue une
source de matière dans l'écoulement, son effet direct sur le débit total est négligé dans l'équation de
continuité :
\begin{equation}
\frac{\partial (Au)}{\partial x} = 0 .
\label{eq:continuite}
\end{equation}
Cette simplification n'est pas une hypothèse indépendante : elle résulte de la moyenne de section de
l'équation d'incompressibilité, comme établi à la section~\ref{sec:reduction-ns}. Le débit \(Q(t)\) est
donc constant le long du conduit à chaque instant, tandis que l'érosion intervient essentiellement dans
l'évolution de la géométrie et dans l'équation de transport de la phase solide.

Sa justification quantitative repose sur la comparaison entre le débit volumique injecté par l'érosion
et le débit transporté. Le terme source volumique par unité de longueur s'écrit
\begin{equation}
S_A = \frac{2\pi R\,\mdot}{\rho_s}
\qquad \left[\mathrm{m^2\,s^{-1}}\right],
\label{eq:source-vol}
\end{equation}
où \(\rho_s\) est la masse volumique apparente saturée du sol intact. Sa comparaison directe avec le
débit \(Q\) (de dimension \(\mathrm{m^3\,s^{-1}}\)) n'est pas homogène : le rapport dépourvu de
dimension est obtenu en intégrant la source sur la longueur du conduit,
\begin{equation}
\frac{\displaystyle\int_0^L S_A\,\mathrm{d}x}{Q}
= \frac{2\pi R\,\mdot\,L}{\rho_s\,Q} .
\label{eq:rapport-injection}
\end{equation}
Pour le cas de référence de la campagne A, on obtient \(S_A \approx 6\times10^{-8}\ \mathrm{m^2\,s^{-1}}\),
\(Q \approx 1.3\times10^{-4}\ \mathrm{m^3\,s^{-1}}\), soit un rapport
\eqref{eq:rapport-injection} de l'ordre de \(6\times10^{-5}\). L'apport de matière érodée modifie donc le
bilan de volume de moins de \(10^{-4}\) : la simplification~\eqref{eq:continuite} est pleinement
justifiée, à condition de conserver la source dans l'équation de transport de la phase solide.

\begin{figure}[H]
\centering
\includegraphics[width=0.85\linewidth]{figure1.png}
\caption{Représentation schématique du conduit effectif unidimensionnel : pressions imposées à
l'entrée et à la sortie, vitesse moyenne \(u\), concentration en suspension \(\phi\), contrainte
pariétale \(\tau_b\) et élargissement induit par l'érosion.}
\label{fig:conduit-1D}
\end{figure}

\subsection{Hypothèses constitutives complémentaires}

Enfin, la réduction adoptée s'accompagne d'hypothèses constitutives supplémentaires qui délimitent le
domaine d'application du modèle :
\begin{itemize}
  \item \textbf{écoulement en conduite lisse} : la fermeture de frottement est de type conduite lisse,
        sans rugosité évolutive ;
  \item \textbf{loi d'érosion linéaire à seuil} : \(\mdot = k_{er}(\tau_b - \tau_c)\) pour
        \(\tau_b > \tau_c\), et \(\mdot = 0\) sinon, sans mécanisme de dépôt ni d'armurage ;
  \item \textbf{suspension diluée} : l'interaction entre particules est négligée dans la masse volumique
        \eqref{eq:rho-mix} ; cette hypothèse est discutée au paragraphe~\ref{par:fm-plafond} ;
  \item \textbf{absence de colmatage, d'armurage et de ségrégation granulométrique}.
\end{itemize}
Ces hypothèses ne constituent pas une description exhaustive de la physique du piping, mais les
simplifications nécessaires à la construction d'un modèle d'ordre réduit destiné à identifier les
mécanismes dominants et à explorer efficacement les régimes d'évolution du conduit.

Cette réduction définit donc le niveau de description à partir duquel est construit le système couplé
étudié dans la suite : l'hydraulique, l'érosion et le transport sont représentés au même niveau de
moyennage longitudinal, tout en conservant les dépendances fonctionnelles nécessaires à la restitution
de leur rétroaction mutuelle.

%=====================================================================================
\section{Réduction dimensionnelle du modèle de Navier--Stokes et nature des fermetures}
\label{sec:reduction-ns}
%=====================================================================================

Le modèle unidimensionnel utilisé dans ce travail peut être présenté de deux manières. La première
consiste à l'énoncer directement sous sa forme réduite, comme un système d'équations posé \emph{a
priori}. La seconde, retenue ici, consiste à l'obtenir par une suite de réductions explicites à partir
des équations de la mécanique des fluides. Cette seconde démarche présente un double intérêt : elle
montre quelles hypothèses sont réellement mobilisées à chaque étape, et elle rend ces hypothèses
quantifiables sous forme de nombres sans dimension. C'est cette quantification qui permet, à la
section~\ref{sec:validite}, de délimiter précisément le domaine dans lequel les résultats du modèle
conservent une valeur quantitative.

\subsection{Du bilan de masse tridimensionnel à la contrainte globale de débit}
\label{sec:div-v}

Pour un fluide incompressible, la conservation de la masse s'écrit
\begin{equation}
\divg\vec{v} = \frac{\partial v_x}{\partial x} + \frac{\partial v_y}{\partial y}
+ \frac{\partial v_z}{\partial z} = 0 .
\label{eq:div-v}
\end{equation}
Cette relation est purement cinématique : elle exprime que la masse d'un volume matériel est constante,
ou, de manière équivalente, que le taux de variation relatif du volume d'une particule fluide est nul.
Elle ne fait intervenir ni la rhéologie du fluide ni le détail de la géométrie interne du conduit. En
particulier, elle reste valable quelle que soit la forme locale du champ de vitesse, y compris dans les
zones de recirculation ou d'établissement.

Intégrons~\eqref{eq:div-v} sur le volume de contrôle compris entre deux sections droites \(\Sigma(x)\) et
\(\Sigma(x+\Delta x)\) normales à l'axe du conduit. En appliquant le théorème de la divergence, le
résultat se réduit aux flux à travers les deux sections :
\begin{equation}
\int_{\Sigma(x)} v_x\,\mathrm{d}A = \int_{\Sigma(x+\Delta x)} v_x\,\mathrm{d}A .
\label{eq:debit-conserve}
\end{equation}
En définissant la vitesse moyenne de section
\begin{equation}
u(x,t) = \frac{1}{A(x,t)}\int_{\Sigma(x)} v_x\,\mathrm{d}A,
\label{eq:def-u}
\end{equation}
la relation~\eqref{eq:debit-conserve} devient \(A(x,t)u(x,t) = \) constante, c'est-à-dire exactement la
forme~\eqref{eq:continuite} utilisée dans le modèle :
\begin{equation}
\frac{\partial (Au)}{\partial x} = 0
\qquad\Longleftrightarrow\qquad
Q(t) = A(x,t)\,u(x,t) .
\label{eq:continuite-reduite}
\end{equation}

Trois conséquences méritent d'être soulignées. Premièrement, l'équation~\eqref{eq:continuite} n'est pas
une hypothèse ajoutée au modèle : c'est la \emph{moyenne de section} de l'incompressibilité, la seule
information perdue dans l'opération étant la répartition radiale de \(v_x\). Deuxièmement, la contrainte
obtenue est \emph{globale} et non locale : le débit \(Q\) est une inconnue scalaire unique pour tout le
conduit, et non un champ. Cette structure est directement responsable de l'organisation numérique
retenue à la section~\ref{sec:hydraulique}, où \(Q\) est déterminé par une condition intégrale
(la pression aval) plutôt que par une équation locale. Troisièmement, le fait que l'érosion ajoute de la
matière aux parois ne contredit pas~\eqref{eq:div-v} : l'apport volumique est simplement négligeable
devant le débit transporté, comme établi par~\eqref{eq:rapport-injection}.

\subsection{Dégénérescence du bilan de quantité de mouvement}
\label{sec:degenerescence}

Le bilan de quantité de mouvement d'un fluide newtonien s'obtient en appliquant le principe fondamental
de la dynamique à un élément fluide, dont la forme finale est l'équation de Navier--Stokes
(dont l'établissement complet est rappelé en annexe) :
\begin{equation}
\rho\frac{D\vec{v}}{Dt} = -\rho g\,\vec{k} - \nabla P + \mu\,\nabla^2\vec{v},
\label{eq:ns}
\end{equation}
où \(\frac{D}{Dt}=\frac{\partial}{\partial t}+\vec{v}\cdot\nabla\) désigne la dérivée matérielle, \(P\) la
pression, \(\mu\) la viscosité dynamique et \(\vec{k}\) la verticale ascendante. Les quatre contributions
sont respectivement la force d'inertie, le poids, la force de pression et la force visqueuse. Le passage
à la formulation monodimensionnelle retenue dans ce travail s'effectue par trois réductions successives.

\paragraph{Gravité et pression réduite.} En retranchant l'équation hydrostatique de l'équation de
Navier--Stokes, le terme de gravité disparaît au profit de la pression réduite
\(\Chat = P + \rho g z\) : seul le gradient de pression \emph{en excès} de l'hydrostatique met le fluide
en mouvement. Le modèle de conduit horizontal n'a donc pas à traiter explicitement la pesanteur, ce qui
justifie l'écriture du bilan sous la forme \(\partial p/\partial x = 2\tau_b/R\).

\paragraph{Moyennage de section.} En supposant l'écoulement axisymétrique et en intégrant le bilan de
quantité de mouvement sur la section, on obtient pour la composante axiale
\begin{equation}
\underbrace{\rho\left(\frac{\partial u}{\partial t} + u\frac{\partial u}{\partial x}\right)}_{\text{inertie}}
+ \frac{\partial p}{\partial x}
= -\underbrace{\frac{2\tau_b}{R}}_{\text{frottement pariétal}},
\label{eq:momentum-1D}
\end{equation}
où \(\tau_b\) est la contrainte de cisaillement exercée par la paroi sur le fluide. Le terme visqueux
volumique a été converti en une contrainte pariétale : c'est cette opération qui introduit la nécessité
d'une \emph{fermeture}, c'est-à-dire d'une relation reliant \(\tau_b\) au champ de vitesse moyen.

\paragraph{Annulation du terme inertiel en géométrie uniforme.} Si la géométrie est uniforme le long du
conduit, alors \(A\) et donc \(u = Q/A\) sont indépendants de \(x\) : le terme convectif
\(u\,\partial u/\partial x\) est \emph{exactement} nul. Le caractère stationnaire en temps est une
conséquence distincte, analysée au paragraphe~\ref{sec:echelles}. Sous ces deux conditions,
l'équation~\eqref{eq:momentum-1D} dégénère en une relation locale exacte entre le gradient de pression
et la contrainte pariétale :
\begin{equation}
\frac{\partial p}{\partial x} = -\frac{2\tau_b}{R},
\qquad\text{soit, en valeur absolue,}\qquad
\left|\frac{\partial p}{\partial x}\right| = \frac{2\tau_b}{R}.
\label{eq:quasi-statique}
\end{equation}
Cette relation, qui constitue l'équation hydraulique fondamentale du modèle, a une conséquence numérique
majeure. Le problème hydraulique n'est plus un problème de champ : la pression s'obtient par une simple
\emph{quadrature} en espace,
\begin{equation}
p(x,t) = P_{\mathrm{in}} + \int_0^x \frac{2\tau_b(\xi,t)}{R(\xi,t)}\,\mathrm{d}\xi ,
\label{eq:quadrature-p}
\end{equation}
et la seule inconnue restante est le scalaire \(Q(t)\), déterminé par la condition intégrale
\(p(L,t) = p_{\mathrm{out}}(Q)\). Il n'y a donc, dans le modèle, ni inversion d'un opérateur
différentiel spatial, ni résolution d'un système linéaire creux, ni itération de type Poisson : le
problème de champ est remplacé par une quadrature et un problème scalaire de recherche de racine. Cette
structure est exploitée directement à la section~\ref{sec:hydraulique}.

\subsection{De la solution de Poiseuille à la forme darcienne : le rôle de la fermeture}
\label{sec:poiseuille}

L'équation~\eqref{eq:quasi-statique} ne constitue pas à elle seule une relation fermée : \(\tau_b\)
dépend de l'écoulement. La manière dont cette fermeture est construite détermine entièrement la nature
du modèle, et il est instructif de la construire d'abord dans le cas le plus simple, celui du régime
visqueux rampant.

\paragraph{Cas rampant : la solution de Poiseuille.} Pour un écoulement laminaire établi dans un tube
circulaire de rayon \(R\), en négligeant le terme inertiel (hypothèse justifiée lorsque le nombre de
Reynolds est très inférieur à l'unité), l'équation de Navier--Stokes se réduit à l'équilibre entre le
gradient de pression réduite et la diffusion visqueuse. La résolution avec la condition de non-glissement
à la paroi donne le profil parabolique
\begin{equation}
v_x(r) = \frac{1}{4\mu}\frac{\mathrm{d}\Chat}{\mathrm{d}x}\left(r^2 - R^2\right),
\label{eq:poiseuille}
\end{equation}
et, en intégrant sur la section, le débit volumique
\begin{equation}
Q = -\frac{\pi R^4}{8\mu}\frac{\mathrm{d}\Chat}{\mathrm{d}x}
\qquad\Longrightarrow\qquad
q = \frac{Q}{\pi R^2} = -\frac{R^2}{8\mu}\frac{\mathrm{d}\Chat}{\mathrm{d}x}.
\label{eq:poiseuille-q}
\end{equation}
En introduisant la charge hydraulique \(\Chat = \rho g h\), cette dernière relation prend exactement la
forme de la loi de Darcy,
\begin{equation}
q = -K\frac{\mathrm{d}h}{\mathrm{d}s},
\qquad
K = \frac{\rho g R^2}{8\mu} .
\label{eq:darcy}
\end{equation}
Autrement dit, la perméabilité d'un milieu poreux idéalisé par un faisceau de conduits est une
manifestation directe de la \emph{fermeture visqueuse} : la loi de Darcy est une loi de comportement,
pas une loi fondamentale. La même structure se retrouve pour l'écoulement dans une fracture
(\(q \propto -W^2\,\mathrm{d}\Chat/\mathrm{d}x\)) ou pour tout autre géométrie dont la solution locale
est linéaire en vitesse. Ce résultat est instructif pour le modèle de piping, mais il doit être manipulé
avec précaution.

\paragraph{Dégénérescence attendue et changement de régime.} Les relations~\eqref{eq:poiseuille-q}
et~\eqref{eq:darcy} montrent qu'une fermeture locale, résolue sur la géométrie du conduit, produit
toujours une relation du type \og débit proportionnel à un gradient \fg{}. Deux régimes doivent
néanmoins être distingués avec soin :
\begin{itemize}
  \item en régime \emph{laminaire} (\(Re \ll 1\)), la contrainte pariétale est linéaire en vitesse,
        \(\tau_b = 4\mu u/R\), et le débit varie comme \(R^4\) à gradient imposé ;
  \item en régime \emph{turbulent lisse} (\(Re \sim 10^3\)--\(10^5\)), la fermeture devient quadratique,
        \(\tau_b \propto u^2\), et le débit varie comme \(R^{19/7}\) à pression imposée pour une loi de
        Blasius. Cette puissance, qui peut paraître élevée, traduit le double effet de l'élargissement :
        la section croît comme \(R^2\) et la vitesse accessible comme \(R^{5/7}\).
\end{itemize}
Le conduit de renard hydraulique relève du second régime : pour la configuration de référence,
\(Re \approx 2.7\times10^4\). La fermeture retenue dans ce travail est donc de type Darcy--Weisbach,
\begin{equation}
\tau_b = -\frac{1}{8}f_D(Re)\,f_m(\phi)\,\rho(\phi)\,u|u|,
\label{eq:dw}
\end{equation}
et non de type Poiseuille. La solution~\eqref{eq:poiseuille-q} n'est pas utilisée pour décrire le
conduit : elle sert de \emph{gabarit structurel}, qui indique quelle opération est effectuée dans le
modèle (remplacer un champ visqueux par une relation locale entre vitesse et contrainte) et pourquoi
cette opération est légitime dès lors que le conduit est élancé. La justification empirique de la
fermeture~\eqref{eq:dw} et de son facteur correctif \(f_m(\phi)\) est donnée à la
section~\ref{sec:fermeture-hydraulique}.

\subsection{Hiérarchie des échelles et caractère quasi-statique}
\label{sec:echelles}

La suppression du terme inertiel et du caractère instationnaire de l'équation~\eqref{eq:momentum-1D}
repose sur une séparation d'échelles de temps qu'il convient de quantifier précisément. Trois échelles
doivent être comparées à l'échelle d'évolution géométrique
\begin{equation}
\ter = \frac{2L\rho_s}{k_{er}\,\Delta p},
\label{eq:ter}
\end{equation}
qui est le temps caractéristique d'érosion introduit par le modèle, et dont la
section~\ref{sec:instabilite} montre qu'il constitue également le taux de croissance exact du régime
d'emballement.

\paragraph{Échelle convective.} Le temps de transit d'une particule fluide dans le conduit vaut
\(\tau_h \sim L/u\). Pour le cas de référence (\(L = 0.117\) m, \(u \approx 4.5\) m/s), on trouve
\(\tau_h \approx 2.6\times10^{-2}\) s, soit
\begin{equation}
\frac{\tau_h}{\ter} \sim 3\times10^{-5} .
\label{eq:rapport-tau}
\end{equation}
La réponse hydraulique au sens \emph{advectif} est donc environ \(10^4\) fois plus rapide que
l'évolution géométrique : à l'échelle d'un pas d'érosion, l'écoulement a atteint son régime établi.

\paragraph{Échelle de propagation des ondes.} Une troisième échelle, souvent omise, est le temps de
propagation d'une onde de pression, \(\tau_a \sim L/c\) où \(c\) est la célérité des ondes dans le
fluide (\(\approx 1480\) m/s pour l'eau). Pour la même configuration, \(\tau_a \approx 8\times10^{-5}\) s,
soit
\begin{equation}
\frac{\tau_a}{\ter} \sim 10^{-7} .
\label{eq:rapport-ta}
\end{equation}
Cette seconde inégalité est décisive : elle garantit que la dynamique \emph{acoustique} est entièrement
relaxée entre deux états d'érosion, et donc qu'aucune onde de pression transitoire ne peut se
développer à l'échelle du calcul. Il est important de noter que c'est cette séparation acoustique, et
non la seule topologie unidimensionnelle du domaine, qui autorise à traiter le champ de pression comme
une quantité s'ajustant instantanément à l'état géométrique courant. La topologie 1D, quant à elle,
facilite la résolution, mais ne dispenserait pas d'une résolution instationnaire si les échelles de
temps étaient comparables.

\paragraph{Portée pratique.} Ces deux inégalités étant vérifiées sur toute la plage de paramètres
étudiée, la stratégie de résolution retenue est la suivante : à chaque état géométrique et de
concentration, le champ hydraulique est recalculé comme un état d'équilibre instantané, tandis que les
variables lentes \(R\) et \(\phi\) sont avancées en temps. La cohérence de ce traitement est contrôlée à
chaque pas par le diagnostic d'effectivité CFL présenté à la section~\ref{sec:cfl-diagnostic}, qui est
précisément le garde-fou permettant de détecter une éventuelle perte de la séparation d'échelles.

\subsection{Nombres sans dimension et critères de validité}
\label{sec:sansdim}

En rassemblant les résultats des paragraphes précédents, les hypothèses du modèle se traduisent en un
petit nombre de nombres sans dimension dont la valeur fixe le domaine de validité. Le
tableau~\ref{tab:sansdim} les récapitule pour la configuration de référence de la campagne A
(échelle laboratoire, \(L = 0.117\) m, \(R_0 = 3\) mm, \(\Delta p = 0.5\) m de colonne d'eau).

\begin{table}[H]
\centering
\caption{Nombres sans dimension et critères de validité de la réduction unidimensionnelle. Les valeurs
correspondent à l'état initial de la campagne A (\(u_0 \approx 4.5\) m/s, \(R_0 = 3\) mm).}
\label{tab:sansdim}
\small
\begin{tabular}{lllll}
\toprule
Grandeur & Définition & Valeur & Critère & Rôle \\
\midrule
Nombre de Reynolds & \(Re = 2\rho u R/\mu_w\) & \(2.7\times10^{4}\) & \(\gg 2300\) & régime turbulent \\
Élancement du conduit & \(R/L\) & \(0.026\) (initial) & \(< 0.05\) & validité 1D \\
Rapport d'échelles convectif & \(\tau_h/\ter = L/(u\,\ter)\) & \(3\times10^{-5}\) & \(\ll 1\) & quasi-statique \\
Rapport d'échelles acoustique & \(\tau_a/\ter = L/(c\,\ter)\) & \(10^{-7}\) & \(\ll 1\) & pression d'équilibre \\
Longueur d'établissement & \(L_e/D\) & \(10\)--\(40\) & \(\ll L/D = 19.5\) & écoulement établi \\
Inertie / frottement & \(8\,|\partial_x R|/f_D\) & \(0 \to \mathcal{O}(10)\) & \(\ll 1\) & géométrie graduelle \\
Charge cinétique / charge motrice & \(\frac{1}{2}\rho u^2/\Delta p\) & \(2.1\) & \(\ll 1\) & effets d'entrée négligeables \\
Supercriticité érosive & \(\tau_b/\tau_c\) & \(63\) & \(>1\) & érosion active \\
\bottomrule
\end{tabular}
\end{table}

Trois de ces critères ne sont pas satisfaits dans la configuration de référence et doivent donc être
commentés sans détour :
\begin{enumerate}
  \item \textbf{Inertie / frottement.} Le rapport \(8|\partial_x R|/f_D\) est nul tant que la géométrie
        reste uniforme, mais croît avec le gradient de rayon. Il quantifie la perte de validité de
        l'équation~\eqref{eq:quasi-statique} lorsque le profil se creuse : une fois le conduit élargi de
        manière non uniforme, le terme convectif n'est plus négligeable et la relation locale doit être
        considérée comme approchée.
  \item \textbf{Longueur d'établissement.} Pour un écoulement turbulent, la longueur nécessaire à
        l'établissement du profil de vitesse est de l'ordre de \(10\) à \(40\) diamètres. Avec
        \(L/D \approx 20\), le conduit de la campagne A est du même ordre que sa propre longueur
        d'établissement : l'application d'une fermeture de conduite \emph{établie} sur toute la
        longueur constitue une approximation, et non un résultat.
  \item \textbf{Charge cinétique d'entrée.} La condition d'entrée du modèle, \(p(0,t) = P_{\mathrm{in}}\)
        avec \(u(0,t) = Q/A_0\), suppose le fluide déjà en mouvement à la pression imposée. Or
        l'accélération depuis le réservoir amont exige au minimum une charge cinétique
        \(\frac{1}{2}\rho u^2\), dont la valeur dépasse ici la charge motrice disponible. Ce point est
        analysé en détail au paragraphe~\ref{par:entree}.
\end{enumerate}

\subsection{Effets non représentés par la réduction}
\label{sec:effets-non-representes}

\paragraph{Effets d'entrée et échelle de la configuration.}
\label{par:entree}
La relation entre l'état de référence et la condition d'entrée mérite un examen quantitatif, car il
conditionne l'interprétation des temps caractéristiques. Avec \(u_0 \approx 4.5\) m/s, la charge
cinétique vaut \(u_0^2/2g \approx 1.04\) m, à comparer à la charge motrice imposée \(\Delta p/\rho_w g =
0.5\) m. Autrement dit, l'énergie cinétique minimale que doit acquérir le fluide pour atteindre cette
vitesse dépasse à elle seule la charge disponible ; de manière équivalente, la borne de Bernoulli
frictionless \(u \le \sqrt{2\Delta p/\rho_w}\) donne \(3.13\) m/s, valeur inférieure à la vitesse
calculée. La vitesse rapportée par le modèle n'est donc réalisable que si le fluide entre dans le conduit
avec une charge supérieure à \(P_{\mathrm{in}}\), ce que la condition aux limites ne représente pas.

L'ordre de grandeur de l'effet est facile à évaluer. En introduisant un coefficient de perte singulière
d'entrée \(\kin\) dans le bilan de charge,
\begin{equation}
\Delta p = \underbrace{\frac{f_D\,\rho\,u^2\,L}{4R}}_{\text{frottement réparti}}
+ \underbrace{\kin\,\frac{1}{2}\rho u^2}_{\text{perte d'entrée}},
\label{eq:charge-entree}
\end{equation}
l'application numérique pour \(\kin = 0.5\) (entrée à arête vive) conduit à \(u = 3.1\) m/s au lieu de
\(4.5\) m/s, soit une réduction de \(31\%\) du débit et de \(49\%\) de la contrainte pariétale
\(\tau_b\). La simulation couplée confirme cet ordre de grandeur : \(u_0 = 3.17\) m/s
(\(-30\%\)) et \(\tau_{b,0} = 30.7\) Pa (\(-51\%\)), pour un temps de doublement porté de \(668\) à
\(1330\) s (\emph{cf.} section~\ref{sec:kin}). Comme \(\ter \propto 1/(\tau_b-\tau_c)\), les temps d'érosion varient dans le même rapport :
la représentation des effets d'entrée n'est donc pas un raffinement, c'est une question de premier
ordre. Deux choix sont possibles et doivent être explicités : soit introduire \(\kin\) dans le résidu
hydraulique, par symétrie exacte avec \(\kout\) (section~\ref{sec:residu}), soit interpréter
\(P_{\mathrm{in}}\) comme la pression \emph{interne} au conduit, mesurée après la zone d'entrée, et le
justifier expérimentalement.

\paragraph{Inertie convective et gradients de géométrie.} Le terme \(\rho u\,\partial u/\partial x\)
n'est rigoureusement nul que pour une géométrie uniforme. Sa contribution relative au frottement
s'écrit, en utilisant \(u = Q/(\pi R^2)\),
\begin{equation}
\frac{\left|\rho u\,\partial_x u\right|}{\left|2\tau_b/R\right|}
= \frac{8\,|\partial_x R|}{f_D}.
\label{eq:inertie-frottement}
\end{equation}
Cette relation fournit un critère simple : le modèle reste quasi-statique tant que
\(|\partial_x R| \ll f_D/8 \approx 3\times10^{-3}\). Comme le profil de rayon s'écarte rapidement de
l'uniformité dans les régimes d'emballement, ce critère constitue, avec \(R/L < 0.05\), la seconde
condition de validité à surveiller.

\paragraph{Retour sur la nature du problème hydraulique.} Les paragraphes précédents permettent de
préciser une question souvent formulée de manière ambiguë. Le fait que l'écoulement soit confiné et
quasi-1D n'implique pas que le champ de pression soit gouverné par une équation de Poisson. Une équation
de Poisson est une équation \emph{tridimensionnelle} qui exprime le bilan local de masse et requiert, à
chaque instant, une inversion globale du champ. Dans la réduction retenue ici, cette inversion n'a pas
lieu : la relation~\eqref{eq:div-v} a déjà été intégrée en une contrainte globale sur \(Q\), et le champ
de pression est reconstruit par la quadrature~\eqref{eq:quadrature-p}. Formellement :
\begin{itemize}
  \item ce qui subsiste de l'équation de continuité est une \emph{contrainte scalaire} (\eqref{eq:continuite}),
        et non une équation aux dérivées partielles ;
  \item le champ de pression n'est pas la solution d'un problème aux limites du second ordre, mais
        l'intégrale d'une relation locale~;
  \item l'information globale (la condition de pression aval) est portée par une \emph{unique} inconnue
        scalaire \(Q\), traitée par recherche de racine.
\end{itemize}
L'analogie avec la théorie des écoulements souterrains se situe donc à un autre niveau : les solutions de
Poiseuille et de Darcy rappellent que la \emph{fermeture locale} (relation vitesse--contrainte) est
l'élément qui remplace l'opérateur de diffusion dans la description réduite. C'est bien cette fermeture,
et non le caractère elliptique du problème 3D d'origine, qui subsiste dans le modèle 1D.


%=====================================================================================
\section{Structure du couplage hydraulique--érosion--transport}
\label{sec:couplage}
%=====================================================================================

La formulation adoptée repose sur un couplage bidirectionnel entre les processus hydrauliques, l'érosion
de la paroi du conduit et le transport des matériaux détachés. Les différentes équations ne constituent
pas des sous-modèles indépendants : les variables produites par chacune d'elles interviennent directement
dans les autres, de sorte que l'état du système à un instant donné résulte d'une succession de
rétroactions locales et globales. La figure~\ref{fig:boucle} représente cette organisation sous la forme
d'une boucle : la différence de pression imposée détermine le débit et la vitesse d'écoulement ; la
vitesse contrôle la contrainte de cisaillement exercée sur la paroi ; cette contrainte gouverne l'érosion
et l'agrandissement du conduit ; les matériaux érodés alimentent le champ de concentration en suspension ;
à son tour, cette concentration modifie les propriétés hydrauliques du mélange et réintroduit son effet
dans la détermination de l'écoulement.

\begin{figure}[H]
\centering
\includegraphics[width=\linewidth]{figure2.png}
\caption{Structure de rétroaction du modèle couplé : hydraulique, frottement du mélange, érosion,
évolution géométrique et transport en suspension.}
\label{fig:boucle}
\end{figure}

\subsection{Premier maillon : le problème hydraulique}

À chaque instant, l'état géométrique du conduit, défini par \(R(x,t)\), ainsi que la concentration locale
\(\phi(x,t)\), déterminent la section \(A = \pi R^2\), la masse volumique du mélange~\eqref{eq:rho-mix} et
le facteur de résistance du mélange \(f_m(\phi)\). La fermeture de type Darcy--Weisbach~\eqref{eq:dw}
relie alors la contrainte pariétale à la vitesse moyenne. Le nombre de Reynolds,
\begin{equation}
Re = \frac{2\,\rho(\phi)\,u\,R}{\mu_w},
\label{eq:re}
\end{equation}
fixe le facteur de frottement de Darcy \(f_D\), tandis que \(f_m\) traduit l'effet de la concentration des
particules sur la résistance à l'écoulement. Le modèle réalise ainsi un premier couplage entre le
transport des particules et la mécanique de l'écoulement. Le choix de la masse volumique du \emph{mélange}
(dans~\eqref{eq:re}, et non de l'eau seule) est cohérent avec la fermeture~\eqref{eq:dw}, où \(\rho(\phi)\)
apparaît également : le nombre de Reynolds et la contrainte pariétale font alors intervenir la même
densité effective.

\subsection{Deuxième maillon : la loi d'érosion}

La contrainte pariétale issue du premier maillon est comparée à une contrainte critique \(\tau_c\). Le flux
massique de matière détachée est donné par la loi à seuil
\begin{equation}
\mdot =
\begin{cases}
k_{er}\left(|\tau_b| - \tau_c\right), & |\tau_b| > \tau_c,\\[2mm]
0, & |\tau_b| \le \tau_c,
\end{cases}
\label{eq:erosion}
\end{equation}
où \(k_{er}\) est le coefficient d'érosion. Cette relation introduit une non-linéarité fondamentale :
tant que la contrainte hydraulique demeure inférieure à la valeur critique, aucune érosion n'est produite ;
lorsque le seuil est franchi, le taux d'érosion croît linéairement avec l'excès de contrainte. Le flux
érodé est ensuite converti en évolution géométrique par
\begin{equation}
\frac{\partial R}{\partial t} = \frac{\mdot}{\rho_s}\sqrt{1+\left(\frac{\partial R}{\partial x}\right)^2},
\label{eq:dRdt}
\end{equation}
où \(\rho_s\) représente la masse volumique apparente saturée du sol intact et où le facteur
\(\sqrt{1+(\partial_x R)^2}\) traduit la surface réelle de paroi par unité de longueur axiale. L'érosion
agit donc directement sur la géométrie, et cette modification de \(R\) rétroagit immédiatement sur le
problème hydraulique.

Cette rétroaction géométrique est particulièrement importante dans les régimes à pression imposée. Une
augmentation du rayon accroît la section disponible à l'écoulement et diminue la résistance hydraulique
du conduit ; le débit nécessaire pour satisfaire les pressions imposées augmente alors, ce qui accroît la
vitesse et la contrainte de cisaillement. Le mécanisme peut être schématisé par la chaîne
\begin{equation}
R \uparrow \;\Rightarrow\; \text{résistance} \downarrow \;\Rightarrow\; Q \uparrow \;\Rightarrow\;
u \uparrow \;\Rightarrow\; \tau_b \uparrow \;\Rightarrow\; \mdot \uparrow \;\Rightarrow\; R \uparrow .
\label{eq:retro-positive}
\end{equation}
Le modèle contient ainsi, de manière structurelle, une rétroaction positive susceptible de conduire à une
amplification rapide de l'érosion. La section~\ref{sec:instabilite} montre que cette boucle admet une
traduction analytique exacte en géométrie uniforme, et que le régime d'emballement est une propriété de la
\emph{condition aux limites} (pression imposée) et non de la loi d'érosion retenue.

\subsection{Troisième maillon : le transport des matériaux érodés}

La fraction volumique de solides en suspension \(\phi\) satisfait l'équation d'advection--relaxation
\begin{equation}
\frac{\partial\phi}{\partial t} + \frac{\partial(u\phi)}{\partial x} = k\left(\phisol - \phi\right),
\qquad
k = \frac{2\mdot}{R\,\rho_s},
\label{eq:transport}
\end{equation}
où \(k\) traduit directement l'intensité locale de l'érosion. Toute augmentation de \(\mdot\) produite par
le sous-modèle hydraulique se traduit par une augmentation de la source de particules dans l'équation de
transport. À l'inverse, le champ de concentration produit par cette équation intervient de nouveau dans
la fermeture hydraulique par l'intermédiaire de \(\rho(\phi)\) et de \(f_m(\phi)\).

Le choix d'une source de type relaxation,
\begin{equation}
S_\phi = k\left(\phisol - \phi\right),
\label{eq:source-phi}
\end{equation}
introduit un mécanisme de saturation : lorsque \(\phi \to \phisol\), la source tend vers zéro et
l'accumulation de solides en suspension est bornée par la fraction solide du matériau intact. Ce mécanisme
d'auto-limitation est absent d'une loi de source directe non saturante, qui autoriserait
\(\phi > \phisol\) en l'absence de contrôle explicite.

\subsection{Concurrence entre rétroaction géométrique et rétroaction résistive}

Le système possède donc simultanément deux rétroactions de sens opposé. La première, géométrique, suit la
chaîne~\eqref{eq:retro-positive} et amplifie l'érosion. La seconde, résistive, suit la chaîne
\begin{equation}
\mdot \uparrow \;\Rightarrow\; \phi \uparrow \;\Rightarrow\; \rho \uparrow,\; f_m \uparrow
\;\Rightarrow\; \text{résistance hydraulique} \uparrow,\; Q \downarrow,\; \tau_b \downarrow,
\label{eq:retro-negative}
\end{equation}
et tend à s'opposer à l'amplification. L'importance relative de ces deux mécanismes dépend de la
géométrie, de l'échelle du conduit et de la granulométrie ; elle est quantifiée à la
section~\ref{sec:fermeture-hydraulique} au moyen du paramètre \(f_m(\phi)\).

Cette compétition explique également pourquoi la condition d'aval peut modifier substantiellement le
comportement global : l'introduction d'une perte de charge supplémentaire à la sortie, représentée par
\(\kout\), réduit le débit pour une différence de pression donnée et constitue une rétroaction négative
\emph{externe}, qui s'ajoute à la rétroaction résistive interne. La section~\ref{sec:resultats} montre que
cette résistance aval retarde l'agrandissement du conduit et repousse le moment où le seuil géométrique de
validité est atteint, sans pour autant supprimer l'instabilité de fond.

Il convient enfin de distinguer le couplage physique de sa traduction numérique. Le caractère couplé
signifie ici que chaque état temporel est obtenu par une succession de mises à jour interdépendantes :
l'état géométrique et la concentration définissent les propriétés hydrauliques ; la solution hydraulique
fournit la contrainte pariétale ; celle-ci détermine l'érosion ; l'érosion modifie simultanément la
géométrie et la source du transport. Le modèle referme donc la boucle de rétroaction à chaque pas de
temps, selon la séquence détaillée à la section~\ref{sec:algo}.

%=====================================================================================
\section{Fermeture hydraulique et résistance du mélange}
\label{sec:fermeture-hydraulique}
%=====================================================================================

\subsection{Contrainte pariétale et nombre de Reynolds}

L'équation quasi-statique~\eqref{eq:quasi-statique} ne devient une relation fermée qu'une fois précisée la
loi reliant \(\tau_b\) à l'écoulement. La fermeture retenue est de type Darcy--Weisbach,
\begin{equation}
\tau_b = -\frac{1}{8}f_D(Re)\,f_m(\phi)\,\rho(\phi)\,u|u|,
\label{eq:dw2}
\end{equation}
où \(f_D\) est le facteur de frottement de Darcy de l'écoulement de référence, \(f_m(\phi)\) l'augmentation
relative de résistance due aux particules, \(\rho(\phi)\) la masse volumique du mélange~\eqref{eq:rho-mix}
et \(u\) la vitesse moyenne. Cette écriture distingue deux contributions : celle de la dynamique de
l'écoulement, portée par \(f_D\), et celle de la concentration des particules, portée par \(f_m\).

Le régime hydraulique est caractérisé par le nombre de Reynolds~\eqref{eq:re}, dans lequel le diamètre
hydraulique vaut \(2R\), cohérent avec la représentation du conduit par une section circulaire de rayon
effectif \(R\). Le facteur de frottement est calculé par une corrélation explicite de type conduite lisse.
Le choix d'une corrélation explicite, plutôt que d'une relation implicite de type Colebrook, répond
directement à la philosophie d'ordre réduit : dans une simulation transitoire comportant plusieurs milliers
de recalculs hydrauliques, l'évaluation du frottement doit rester peu coûteuse. Deux options ont été
utilisées dans ce travail :
\begin{itemize}
  \item une corrélation explicite de type Blasius, \(f_D = 0.3164\,Re^{-1/4}\), qui constitue
        l'approximation de référence pour un conduit lisse en régime turbulent ;
  \item une corrélation de type Barenblatt, fondée sur la loi de puissance en \(Re\) et valable sur une
        plage de Reynolds plus étendue.
\end{itemize}
L'écart entre les deux fermetures est faible sur la plage explorée : à \(Re = 2.7\times10^{4}\), les deux
corrélations donnent \(f_D = 0.0247\) et \(0.0246\), soit un écart inférieur à \(0.5\%\). Ce point doit être
mentionné explicitement dans la présentation des résultats, car il garantit que les conclusions ne
dépendent pas du choix de la corrélation. En revanche, la fermeture reste associée à l'hypothèse de
conduite lisse : elle n'intègre ni rugosité évolutive, ni ségrégation granulométrique, ni transition vers
un comportement viscoplastique.

\subsection{Masse volumique du mélange}

La concentration intervient une première fois dans la fermeture par l'intermédiaire de la masse volumique
effective~\eqref{eq:rho-mix}. Cette loi de mélange linéaire correspond à l'approximation de suspension
diluée : l'augmentation de la fraction volumique solide se traduit directement par une augmentation de la
masse volumique apparente du fluide transporté. Pour une vitesse donnée, l'augmentation de \(\phi\) modifie
donc directement le terme \(\rho(\phi)\,u|u|\) de la contrainte pariétale. Le champ de transport n'est pas
une variable purement descriptive : il participe à la détermination de la réponse hydraulique.

\subsection{Résistance supplémentaire du mélange}

L'effet de la suspension sur la résistance hydraulique est représenté par un facteur multiplicatif
\(f_m(\phi)\) qui croît avec la concentration et devient singulier lorsque \(\phi\) tend vers la fraction
solide du sol intact. Le modèle met en œuvre une forme de type Julien, construite sur la longueur de
mélange \(l_m\) et sur la concentration relative,
\begin{equation}
f_m(\phi) = \min\left[f_{m,\max},\;
1 + C_B\,\frac{\rho_p}{\rho(\phi)}\left(\frac{d_p}{l_m}\right)^{2}\lambda(\phi)^{n_\lambda}\right],
\qquad
\lambda(\phi) = \left[\left(\frac{\phisol}{\phi}\right)^{1/3}-1\right]^{-1},
\label{eq:fm}
\end{equation}
avec \(l_m = C_r R\) (ou, dans une variante simplifiée, \(l_m = C_r R_0\)), \(d_p\) le diamètre
caractéristique des particules, \(C_B\) le coefficient dispersif, \(C_r\) le coefficient de longueur de
mélange et \(n_\lambda\) l'exposant de concentration. La fonction \(\lambda(\phi)\) est celle qui décrit
l'accroissement de la dissipation collisionnelle lorsque l'empilement granulaire se densifie ; elle diverge
lorsque \(\phi \to \phisol\), ce qui exprime physiquement la transition vers un comportement de pâte.

Deux propriétés de cette fermeture doivent être soulignées avec précision, car elles conditionnent la
lisibilité des résultats.

\paragraph{Dépendance quadratique en diamètre.} Le terme \((d_p/l_m)^2\) rend la rétroaction rhéologique
extrêmement sensible à la granulométrie. À titre d'illustration, pour une même concentration relative
\(\phi/\phisol = 0.6\), le facteur \(f_m\) passe de \(1.0\)--\(1.1\) pour un silt fin (\(d_p \approx 20\)
\si{\micro m}) à plusieurs unités pour un sable fin (\(d_p \approx 200\ \si{\micro m}\)). Autrement dit,
la compétition décrite à la section~\ref{sec:couplage} n'est significative que pour des matériaux
suffisamment grossiers. Pour un sol silto-argileux fin, la rétroaction résistive interne est
essentiellement inactive et la dynamique est contrôlée par la seule rétroaction
géométrique~\eqref{eq:retro-positive}, modulée par la condition d'aval.

\paragraph{Plafonnement et limite de validité.}
\label{par:fm-plafond}
La présence du minimum dans~\eqref{eq:fm} n'est pas un simple artifice numérique. Sans limitation, la
fermeture devient singulière lorsque \(\phi \to \phisol\), puisque \(\lambda \to +\infty\). Le modèle adopte
donc
\begin{equation}
f_m \le f_{m,\max},
\label{eq:fmmax}
\end{equation}
valeur qui doit être présentée à la fois comme un dispositif de régularisation et comme une limite
physique de la fermeture : lorsque la concentration devient suffisamment importante, l'hypothèse d'un
mélange assimilable à un écoulement turbulent newtonien cesse d'être pertinente et un comportement de type
pâte viscoplastique peut apparaître.

Ce plafond est appliqué \emph{au sein du solveur}, et non en post-traitement des résultats : c'est lui qui
est consommé à chaque pas par le calcul de la contrainte pariétale~\eqref{eq:dw} et par le résidu
hydraulique~\eqref{eq:residu}. Ce point n'est pas un détail d'implémentation, car il change la nature du
problème. En l'absence de plafond, la campagne B (matériau grossier, \(d_p = 1\) mm) conduit à
\(f_m \sim 10^{12}\) dans une zone étroite du conduit : le gradient de résistance locale qui en résulte
concentre l'érosion en un pic isolé, le rayon maximal dépasse alors le rayon de sortie de plus d'un facteur
cinquante, et le calcul s'interrompt sur le critère de sécurité \(R_{\mathrm{break}}\). Le plafond supprime
ce pic, rétablit un champ de résistance régulier et rend au contraire la comparaison avec les temps publiés
possible (section~\ref{sec:verif-fm}).

Le caractère actif ou non de ce plafond doit être documenté pour chaque configuration étudiée, en
indiquant explicitement la valeur maximale de \(f_m\) atteinte au cours du calcul. Deux situations
distinctes peuvent se présenter : soit le plafond n'est jamais atteint, et la fermeture reste dans son
domaine de validité asymptotique ; soit il est atteint, et l'interprétation des résultats doit être
limitée au régime où la fermeture reste valable. Cette distinction est essentielle pour l'interprétation
des simulations de forte érosion : attribuer au modèle une capacité prédictive illimitée lorsque \(\phi\)
approche \(\phisol\) serait incorrect. Elle l'est d'autant plus que le plafond n'est pas une petite
correction : la réponse de la campagne B en dépend qualitativement (section~\ref{sec:kout}).

\paragraph{Cohérence entre les paramètres publiés et les paramètres simulés.} Le lecteur doit pouvoir
retrouver, sans ambiguïté, la relation exacte utilisée et le jeu de paramètres associé. Le
tableau~\ref{tab:param-campagnes} récapitule les deux campagnes de simulation exploitées dans ce travail,
avec les valeurs effectivement utilisées pour chaque grandeur. Les valeurs de \(d_p\), \(C_B\), \(C_r\),
\(n_\lambda\) et \(f_{m,\max}\) y figurent explicitement par campagne, et non comme des constantes
globales, précisément parce que la nature de la rétroaction rhéologique en dépend.

\begin{table}[H]
\centering
\caption{Paramètres des deux campagnes de simulation. Campagne A : échelle laboratoire, référence de
type essai au trou (HET). Campagne B : échelle de terrain, gradient hydraulique élevé.}
\label{tab:param-campagnes}
\small
\begin{tabular}{lccc}
\toprule
Paramètre & Symbole & Campagne A (laboratoire) & Campagne B (terrain) \\
\midrule
Longueur du conduit & \(L\) & \(0.117\) m & \(100\) m \\
Rayon initial & \(R_0\) & \(3\) mm & \(3\) mm \\
Pression d'entrée & \(P_{\mathrm{in}}\) & \(0.5\) m de colonne d'eau & \(3.33\) MPa \\
Pression de sortie & \(P_{\mathrm{out}}\) & \(0\) Pa & \(0\) Pa \\
Contrainte critique & \(\tau_c\) & \(1.0\) Pa & \(10\) Pa \\
Coefficient d'érosion & \(k_{er}\) & \(1.0\times10^{-4}\ \mathrm{s/m}\) & \(1.0\times10^{-2}\ \mathrm{s/m}\) \\
Masse volumique saturée & \(\rho_s\) & \(1990\ \mathrm{kg/m^3}\) & \(1850\ \mathrm{kg/m^3}\) \\
Fraction solide du sol & \(\phisol\) & \(0.60\) & \(0.50\) \\
Diamètre des particules & \(d_p\) & \(20\ \si{\micro m}\) & \(1.0\) mm \\
Coefficient dispersif & \(C_B\) & \(0.2\) & \(0.01\) \\
Coefficient de longueur de mélange & \(C_r\) & \(0.07\) & \(0.10\) \\
Exposant de concentration & \(n_\lambda\) & \(2\) & \(2\) \\
Plafond du facteur de résistance & \(f_{m,\max}\) & \(5\) & \(2000\) \\
Temps caractéristique & \(\ter\) & \(949\) s & \(11.1\) s \\
Nombre de cellules & \(N_x\) & \(300\), \(600\), \(1200\) & \(300\), \(600\) \\
Limiteur du schéma & --- & van Leer / MC & van Leer / MC \\
\bottomrule
\end{tabular}
\begin{tablenotes}
\small
\item Les valeurs de \(\ter\) sont calculées par~\eqref{eq:ter}. Le coefficient \(k_{er}\) est exprimé en
\(\mathrm{s/m}\), conformément à la loi~\eqref{eq:erosion} écrite en flux massique ; toute autre unité doit
être convertie explicitement.
\item Les paramètres \(d_p\), \(C_B\), \(C_r\), \(n_\lambda\) et \(f_{m,\max}\) sont ceux du
code de calcul et sont donnés par campagne. Le calcul de référence de la campagne A utilise
\(d_p = 75\ \si{\micro m}\) au lieu des \(20\ \si{\micro m}\) du sol étudié ; cette substitution est sans
conséquence \emph{mesurable} sur cette configuration, car le facteur de résistance reste égal à \(1.000\)
dans les deux cas (section~\ref{sec:verif-fm}) : la rétroaction rhéologique est inactive à l'échelle du
laboratoire, quel que soit le matériau fin considéré.
\end{tablenotes}
\end{table}

% TODO auteur : verifier que les valeurs de k_er, tau_c et rho_s des deux campagnes correspondent bien
% aux fichiers de configuration effectivement utilises, et completer la ligne "limiteur" avec la valeur
% unique retenue pour les resultats publies.


%=====================================================================================
\section{Organisation générale de l'algorithme de résolution}
\label{sec:algo}
%=====================================================================================

La résolution numérique du modèle repose sur une succession d'étapes interdépendantes destinées à
assurer, à chaque instant, la cohérence entre l'état hydraulique, l'état géométrique du conduit et l'état
de transport des particules. L'organisation générale découle directement de la structure du couplage
présentée à la section~\ref{sec:couplage} : à partir de la géométrie courante \(R(x,t)\) et du champ de
concentration \(\phi(x,t)\), le problème hydraulique est résolu sous les conditions de pression imposées ;
la solution obtenue permet d'évaluer la vitesse et la contrainte de cisaillement à la paroi ; ces grandeurs
déterminent le taux d'érosion, qui contrôle simultanément l'évolution du rayon et l'alimentation du champ
de particules en suspension. Le champ de concentration est ensuite transporté par l'écoulement et corrigé
par l'intégration du terme source d'érosion.

L'état du système à l'instant \(t^n\) peut être représenté par l'ensemble
\begin{equation}
\mathcal{S}^n = \left\{R_i^n, \phi_i^n\right\}_{i=1}^{N_x},
\label{eq:etat}
\end{equation}
auquel s'ajoutent les grandeurs hydrauliques déduites de la résolution du problème, notamment
\(Q^n, u_i^n, p_i^n, \tau_{b,i}^n, \rho_i^n\) et \(f_{m,i}^n\). L'objectif du pas de temps est de construire
l'état suivant \(\mathcal{S}^{n+1}\), en respectant simultanément les conditions hydrauliques imposées, la
conservation du transport et la loi locale d'érosion.

\subsection{Décomposition séquentielle et nature du couplage}

Avant de détailler chaque étape, il importe de préciser la nature exacte du couplage temporel, car
plusieurs formulations sont possibles et la distinction a des conséquences sur la vérification.

\begin{itemize}
  \item \textbf{Le sous-problème hydraulique est implicite.} Pour un état \((R^n,\phi^n)\) donné, le débit
        \(Q^n\) est obtenu par résolution d'une équation scalaire non linéaire dont le résidu dépend de
        \(Q\) par l'intermédiaire de \(u\), \(Re\), \(f_D\) et \(\tau_b\). Il n'y a pas de linéarisation
        lagrangienne du débit : la résolution est bien implicite dans ce sous-problème.
  \item \textbf{Le couplage entre les trois sous-systèmes est séquentiel et décalé d'un pas.} L'érosion et
        le transport au pas \(n\) sont évalués à partir de l'état hydraulique déterminé à la
        \emph{fin} du pas précédent, puis la géométrie et la concentration sont avancées. Le schéma
        global est donc du type \og résoudre l'hydraulique, puis avancer les variables lentes \fg{}, sans
        sous-itération de convergence à l'intérieur du pas.
\end{itemize}
Cette distinction est importante pour l'interprétation de la vérification numérique. Si le schéma était
entièrement implicite (avec itérations internes jusqu'à convergence), l'erreur de troncature temporelle
serait dominée par les schémas d'advection et d'érosion. Dans la formulation retenue, il s'y ajoute une
erreur de \emph{découplage}, proportionnelle au rapport entre le pas de temps et le temps de relaxation du
sous-système hydraulique. La séparation d'échelles établie au paragraphe~\ref{sec:echelles}
(\(\tau_h/\ter \sim 10^{-5}\)) garantit que cette erreur reste négligeable sur toute la plage étudiée, mais
c'est cette garantie qui doit être vérifiée par le diagnostic CFL de la section~\ref{sec:cfl-diagnostic},
et non postulée.

\subsection{Initialisation de l'état du système}

La simulation débute à partir d'une géométrie initiale prescrite \(R(x,0)\), d'un champ initial de
concentration
\begin{equation}
\phi(x,0) = \phi_0(x),
\label{eq:ci}
\end{equation}
ainsi que des paramètres hydrauliques et géotechniques du problème. Les pressions à l'entrée et à la
sortie sont imposées par
\begin{equation}
p(0,t) = P_{\mathrm{in}},
\qquad
p(L,t) = P_{\mathrm{out}},
\label{eq:cl}
\end{equation}
avec la possibilité d'introduire une perte de charge locale à l'exutoire au moyen de \(\kout\). La
discrétisation spatiale fournit une grille de \(N_x\) cellules de largeur \(\Delta x = L/N_x\), et
l'évolution temporelle est réalisée avec un pas \(\dt\) déterminé dynamiquement à partir d'une condition
CFL, sous réserve des bornes \(\dt_{\min}\) et \(\dt_{\max}\).

\subsection{Détermination du débit à partir des conditions de pression}
\label{sec:debit-implicite}

La première étape essentielle de chaque cycle consiste à déterminer \(Q^n\). Contrairement à une
formulation à débit imposé, \(Q^n\) est une inconnue implicite : le champ hydraulique est intégré depuis
l'entrée vers la sortie pour une valeur d'essai du débit, et le résidu scalaire
\begin{equation}
F(Q) = p_{\mathrm{calc}}(L; Q) - p_{\mathrm{out}}(Q)
\label{eq:residu}
\end{equation}
est réduit à zéro. La forme précise de \(p_{\mathrm{out}}(Q)\), qui inclut la perte singulière
\(\kout\frac{1}{2}\rho u^2(L)\) lorsque celle-ci est prise en compte, ainsi que la méthode de résolution,
sont détaillées à la section~\ref{sec:hydraulique}. Le débit ainsi obtenu fournit les vitesses moyennes
\begin{equation}
u_i^n = \frac{Q^n}{A_i^n} = \frac{Q^n}{\pi (R_i^n)^2}.
\label{eq:u-de-Q}
\end{equation}
Cette étape constitue le cœur du traitement des conditions de pression imposées : le débit doit être
recalculé à chaque état d'évolution du conduit.

\subsection{Évaluation des grandeurs hydrauliques et du taux d'érosion}

Une fois le débit déterminé, les grandeurs locales nécessaires au calcul de l'érosion sont évaluées :
masse volumique du mélange~\eqref{eq:rho-mix}, nombre de Reynolds~\eqref{eq:re}, facteur de frottement
\(f_D\), facteur de mélange~\eqref{eq:fm} avec sa borne~\eqref{eq:fmmax}, puis contrainte pariétale
\eqref{eq:dw2}. La contrainte est comparée à la contrainte critique et le flux massique érodé est évalué
par~\eqref{eq:erosion}. Le coefficient local intervenant dans l'équation de transport est alors
\begin{equation}
k_i^n = \frac{2\mdot_i^n}{R_i^n\rho_s}.
\label{eq:k-local}
\end{equation}
L'intérêt de cette organisation est que la source d'érosion n'est pas imposée indépendamment du calcul
hydraulique : elle est une quantité dérivée de l'état hydraulique courant.

\subsection{Détermination adaptative du pas de temps}
\label{sec:dt-adaptatif}

Le pas de temps est ensuite calculé à partir de la contrainte CFL
\begin{equation}
\dt = \max\left[\dt_{\min},\;
\min\left(\frac{\mathrm{CFL}_{\mathrm{tgt}}\,\Delta x}{\max_i|a_i|},\;\dt_{\max}\right)\right],
\label{eq:dt}
\end{equation}
où \(a_i\) est la vitesse utilisée pour le contrôle de la propagation advective et
\(\mathrm{CFL}_{\mathrm{tgt}}\) le nombre de Courant visé. Deux vitesses caractéristiques doivent être
distinguées, et leur confusion constitue une source d'erreur fréquente :
\begin{itemize}
  \item la vitesse du \emph{fluide} \(u = Q/A\), qui fixe le transport de la phase fluide ;
  \item la vitesse de \emph{transport des particules} \(u_p = \beta_u\,u\), qui fixe la propagation de la
        concentration \(\phi\) dans l'équation~\eqref{eq:transport}. Le coefficient \(\beta_u\) est le
        rapport entre la vitesse moyenne des particules et celle du fluide ; il est égal à l'unité dans
        l'hypothèse d'un transport purement advectif sans glissement, hypothèse retenue dans la suite, mais
        il permet de représenter un éventuel retard des particules lorsque le modèle est étendu.
\end{itemize}
C'est la vitesse \(u_p\), et non \(u\), qui doit être utilisée dans le calcul du pas de temps
advectif~\eqref{eq:dt}, puisque c'est elle qui gouverne la propagation du champ \(\phi\). La nécessité
d'un recalcul à chaque pas provient du caractère fortement évolutif de l'écoulement : l'agrandissement du
conduit modifie la section et la vitesse, tandis que l'évolution de la concentration modifie la résistance
du mélange.

\subsection{Avancement de l'équation de transport}

La partie advective de l'équation~\eqref{eq:transport} est traitée séparément, par une méthode de volumes
finis. Sur la grille, l'étape d'advection fournit un état intermédiaire
\begin{equation}
\phi_i^* = \phi_i^n - \frac{\dt}{\Delta x}\left(F_{i+1/2}^n - F_{i-1/2}^n\right),
\label{eq:advection}
\end{equation}
où les flux d'interface sont construits par reconstruction MUSCL avec limiteur. Cette étape assure la
conservation discrète de la quantité transportée tout en limitant les oscillations au voisinage des fronts
de concentration raides. Sa construction détaillée, le choix du limiteur et les questions de positivité
sont traités à la section~\ref{sec:transport}.

\subsection{Intégration du terme source d'érosion}

À l'issue de l'étape d'advection, le terme source est traité localement par séparation d'opérateurs. En
considérant \(k\) constant pendant le pas \(\dt\), l'équation locale
\begin{equation}
\frac{\mathrm{d}\phi}{\mathrm{d}t} = k\left(\phisol - \phi\right)
\label{eq:relax}
\end{equation}
possède la solution analytique exacte
\begin{equation}
\phi_i^{n+1} = \phisol - \left(\phisol - \phi_i^*\right)e^{-k_i^n\dt}.
\label{eq:source-exacte}
\end{equation}
Cette opération présente l'avantage de traiter sans approximation temporelle la raideur locale de la
relaxation pour \(k\) fixé, et de garantir que la concentration ne dépasse pas la limite \(\phisol\) dès
lors que l'état issu de l'étape advective respecte lui-même les bornes admissibles. Ce point, qui n'est pas
automatique, est discuté à la section~\ref{sec:source-erosion}.

\subsection{Mise à jour de la géométrie}

Parallèlement à l'évolution du transport, le taux d'érosion gouverne l'évolution du rayon
par~\eqref{eq:dRdt}. La nouvelle géométrie \(R^{n+1}\) devient la géométrie de référence du cycle
temporel suivant. Cette mise à jour est essentielle puisque le rayon intervient directement dans la
section \(A = \pi R^2\), dans la vitesse \(u = Q/A\), dans le nombre de Reynolds et dans la relation liant
la contrainte pariétale au gradient de pression. Il existe donc une dépendance cyclique explicite,
\begin{equation}
\left(R^n,\phi^n\right) \to \left(Q^n,u^n,\tau_b^n\right) \to \mdot^n
\to \left(R^{n+1},\phi^{n+1}\right) \to \left(Q^{n+1},\dots\right),
\label{eq:cycle}
\end{equation}
conformément à la structure de la boucle décrite à la section~\ref{sec:couplage}. Le décalage d'un pas
entre l'état hydraulique utilisé et l'état géométrique produit correspond à l'erreur de découplage
évoquée au paragraphe~\ref{sec:algo}.

\subsection{Critères d'arrêt et contrôle de la validité géométrique}
\label{sec:criteres-arret}

Après chaque mise à jour, l'état courant est comparé aux critères d'arrêt de la simulation. Trois critères
sont distingués, et leur nature doit rester claire dans l'interprétation des résultats :
\begin{enumerate}
  \item \textbf{critère de validité physique} : l'atteinte de la limite géométrique \(R/L = 0.05\)
        introduite à la section~\ref{sec:hypotheses}. Au-delà, l'approximation 1D n'est plus
        quantitativement valide ; l'arrêt signale une sortie de domaine, non un échec de calcul ;
  \item \textbf{critère de sécurité numérique} : le dépassement d'un rayon maximal
        \(R > R_{\mathrm{break}}\), qui interrompt le calcul avant que les quantités manipulées ne
        sortent du domaine des nombres représentables ;
  \item \textbf{critère de temps} : l'atteinte du temps final prescrit, exprimé en multiples de
        \(\ter\).
\end{enumerate}
La distinction entre ces trois critères est indispensable lorsque l'on compare des simulations de
configurations différentes. Un calcul peut en effet se terminer pour trois raisons de nature totalement
distincte, et la valeur du rapport d'agrandissement atteint dépend directement du critère retenu. Le
tableau de résultats doit donc toujours préciser lequel des trois critères a conduit à l'arrêt.

Dans les simulations sans résistance aval, l'arrêt sur le critère de validité est interprété comme la
manifestation du caractère instable de la configuration sous pression imposée, et non comme une divergence
du solveur numérique. Cette lecture est renforcée par la section~\ref{sec:instabilite}, qui établit que
l'amplification géométrique est structurelle et que son taux est prédictible analytiquement.

\subsection{Algorithme général}

La logique de calcul peut être synthétisée par la procédure suivante.

\begin{enumerate}[label=\arabic*.]
  \item Initialisation : \(R^0\), \(\phi^0\), \(t=0\), état \(\mathcal{S}^0\).
  \item Calcul des propriétés locales du mélange (\(\rho\), \(f_m\), \dots) et du coefficient de
        fermeture \(f_D\).
  \item Résolution implicite du débit : \(F(Q)=0\) par encadrement puis dichotomie.
  \item Calcul de \(u\), \(Re\), \(f_D\), \(\tau_b\) sur toutes les cellules.
  \item Calcul du flux érodé \(\mdot\) et du coefficient de relaxation \(k\).
  \item Calcul du pas de temps par la condition CFL, avec bornes \(\dt_{\min}\) et \(\dt_{\max}\).
  \item Étape d'advection : volumes finis MUSCL avec limiteur.
  \item Intégration analytique du terme source de relaxation.
  \item Mise à jour de la géométrie \(R^{n+1}\).
  \item \(t \leftarrow t+\dt\) ; contrôle des critères d'arrêt (validité, sécurité, temps).
  \item Retour à l'étape 2.
\end{enumerate}

Cette organisation traduit la nature séquentielle du couplage. Elle ne repose pas sur une résolution
monolithique de l'ensemble des équations, mais sur une décomposition structurée du problème multiphysique
en sous-problèmes présentant chacun une méthode numérique appropriée : recherche de racine scalaire robuste
pour l'hydraulique, schéma conservateur à haute résolution pour le transport, intégration exacte pour le
terme source, avancement explicite pour la géométrie. Cette stratégie est cohérente avec l'objectif de
construire un outil suffisamment robuste pour les études paramétriques et suffisamment contrôlable pour une
analyse systématique de convergence.

\begin{table}[H]
\centering
\caption{Récapitulatif des sous-problèmes et des méthodes numériques associées.}
\label{tab:methodes}
\small
\begin{tabular}{llll}
\toprule
Sous-problème & Type & Inconnue & Méthode \\
\midrule
Hydraulique global & scalaire non linéaire & \(Q(t)\) & encadrement + dichotomie \\
Champ de pression & quadrature spatiale & \(p(x,t)\) & intégration explicite \\
Advection de \(\phi\) & hyperbolique conservatif & \(\phi(x,t)\) & volumes finis MUSCL--TVD \\
Source d'érosion & relaxation locale raide & \(\phi(x,t)\) & intégration analytique exacte \\
Géométrie & équation d'évolution locale & \(R(x,t)\) & Euler explicite (ordre 1) \\
\bottomrule
\end{tabular}
\end{table}

% TODO auteur : preciser si la mise a jour de R est effectivement effectuee par un schema d'ordre 1
% (Euler explicite) ou par un schema d'ordre superieur, et le mentionner dans l'analyse d'erreur.


%=====================================================================================
\section{Résolution du problème hydraulique sous conditions de pression imposées}
\label{sec:hydraulique}
%=====================================================================================

Dans le modèle considéré, la résolution hydraulique a pour objectif de déterminer l'état d'écoulement
compatible avec les conditions de pression prescrites aux deux extrémités du conduit. La spécificité de la
formulation réside dans le fait que le débit constitue une inconnue du problème et doit être déterminé à
partir des conditions aux limites plutôt que prescrit \emph{a priori}. Cette section expose successivement
la formulation quasi-stationnaire et sa quadrature, la fermeture complète et le caractère non linéaire du
gradient de pression, la formulation implicite du débit, l'algorithme de résolution, puis la justification
de la robustesse et du coût de la méthode retenue.

\subsection{Formulation quasi-stationnaire et quadrature du champ de pression}

Comme établi au paragraphe~\ref{sec:echelles}, le rapport \(\tau_h/\ter \sim 10^{-5}\) et le rapport
acoustique \(\tau_a/\ter \sim 10^{-7}\) sont tous deux très faibles. Il en résulte que, pour une géométrie
et un état de concentration donnés, le champ hydraulique atteint son équilibre beaucoup plus rapidement
que le conduit n'évolue sous l'effet de l'érosion : les termes transitoires associés à l'inertie de
l'écoulement ne sont pas résolus explicitement dans la réponse hydraulique à chaque instant. L'équation
de quantité de mouvement se réduit alors à l'équilibre local
\begin{equation}
\frac{\partial p}{\partial x} = \frac{2\tau_b}{R},
\label{eq:quasi}
\end{equation}
où \(p(x,t)\) désigne la pression moyenne de section, \(R(x,t)\) le rayon effectif et \(\tau_b(x,t)\) la
contrainte de cisaillement à la paroi.

L'intérêt de cette formulation est de transformer le problème hydraulique en une relation différentielle
spatiale intégrable le long du conduit à partir de la pression connue à l'entrée. En imposant
\(p(0,t) = P_{\mathrm{in}}\), on obtient
\begin{equation}
p(x,t) = P_{\mathrm{in}} + \int_0^x \frac{2\tau_b(\xi,t)}{R(\xi,t)}\,\mathrm{d}\xi .
\label{eq:quadrature}
\end{equation}
La distribution de pression résulte donc directement de l'intégration du bilan hydraulique longitudinal,
et la condition imposée à l'extrémité aval intervient comme condition de fermeture du problème, ce qui
conduit à la détermination implicite du débit présentée au paragraphe~\ref{sec:residu}.

Cette formulation présente également un intérêt numérique important. Puisque l'écoulement est traité comme
quasi-stationnaire, il n'est pas nécessaire de résoudre une équation temporelle supplémentaire pour la
vitesse ou la pression entre deux états successifs du système : à chaque instant macroscopique \(t^n\), le
champ hydraulique est recalculé à partir de l'état courant du conduit. Le problème temporel est ainsi
porté par les variables lentes, le rayon \(R\) et la concentration \(\phi\), tandis que la pression et le
débit sont des variables hydrauliques déterminées par l'état instantané du système. Comme souligné à la
section~\ref{sec:effets-non-representes}, cette situation ne saurait être confondue avec la résolution
d'une équation de type Poisson : l'équation de continuité a déjà été intégrée en une contrainte globale,
et le champ de pression est obtenu par la quadrature~\eqref{eq:quadrature}, sans inversion d'opérateur
différentiel spatial.

\subsection{Fermeture complète et non-linéarité du gradient de pression}

En substituant la fermeture~\eqref{eq:dw2} dans l'équation~\eqref{eq:quasi}, le gradient de pression
s'écrit
\begin{equation}
\frac{\partial p}{\partial x}
= -\frac{f_D(Re)\,f_m(\phi)\,\rho(\phi)}{4R}\,u\,|u|,
\label{eq:dp-dx}
\end{equation}
relation dans laquelle la vitesse est elle-même reliée au débit par \(u(x,t) = Q(t)/\left[\pi R^2(x,t)\right]\).
Cette écriture montre explicitement que le gradient de pression dépend simultanément de la géométrie, de
la concentration et de la vitesse,
\begin{equation}
p_x = p_x\left(R,\phi,u\right),
\qquad\text{avec}\qquad
Q \to u \to Re \to f_D \to \tau_b \to \frac{\partial p}{\partial x},
\label{eq:chaine-non-lineaire}
\end{equation}
chaîne à laquelle s'ajoute la dépendance en concentration
\(\phi \to \rho(\phi), f_m(\phi) \to \tau_b \to \partial_x p\). Le débit ne peut donc pas être obtenu par
une relation linéaire entre perte de charge et débit : il doit être recherché comme la valeur qui rend
l'intégration de~\eqref{eq:dp-dx} compatible avec la condition de pression aval.

Il est utile de distinguer les deux sources de non-linéarité, car elles n'ont pas la même nature :
\begin{itemize}
  \item une non-linéarité \emph{cinématique}, géométrique, provenant de \(u \propto Q/R^2\) et de la
        dépendance \(f_D \propto Re^{-1/4}\) : elle existe même en fluide clair et en géométrie fixée ;
  \item une non-linéarité \emph{rhéologique}, provenant de \(\rho(\phi)\) et surtout de \(f_m(\phi)\) qui
        contient le terme \(\lambda(\phi)^{n_\lambda}\) divergente lorsque \(\phi \to \phisol\). Même si le
        plafond~\eqref{eq:fmmax} borne la dépendance, celle-ci reste très raide à proximité de la limite de
        concentration.
\end{itemize}
C'est la seconde qui motive le choix d'un solveur à convergence garantie plutôt qu'à convergence rapide,
argument développé au paragraphe~\ref{sec:robustesse}.

\subsection{Formulation implicite du débit}
\label{sec:residu}

À un instant donné \(t^n\), la géométrie \(R(x,t^n)\) et la concentration \(\phi(x,t^n)\) sont connues. Pour
une valeur d'essai \(Q\), la vitesse est obtenue par la continuité,
\begin{equation}
u(x;Q) = \frac{Q}{\pi R^2(x)},
\label{eq:uQ}
\end{equation}
puis le nombre de Reynolds et le facteur de frottement sont évalués sur chaque cellule. Le débit admissible
est la racine du résidu scalaire
\begin{equation}
F(Q) = p_{\mathrm{calc}}(L;Q) - p_{\mathrm{out}}(Q) = 0,
\label{eq:FQ}
\end{equation}
où la pression calculée à l'exutoire est donnée par la quadrature~\eqref{eq:quadrature} et où la pression
de sortie \emph{imposée} doit tenir compte de la perte singulière lorsque celle-ci est représentée :
\begin{equation}
p_{\mathrm{out}}(Q) = P_{\mathrm{out}} + \kout\,\frac{1}{2}\rho(L)\,u(L;Q)^2 .
\label{eq:pout}
\end{equation}
Deux configurations doivent être distinguées :
\begin{itemize}
  \item \(\kout = 0\) : la condition aval se réduit à \(p(L,t) = P_{\mathrm{out}}\) ; c'est le cas de
        référence utilisé dans la plupart des simulations présentées ;
  \item \(\kout \neq 0\) : la perte locale est intégrée au résidu. Cette configuration représente une
        restriction, un filtre ou un contrôle aval, et joue le rôle de rétroaction négative externe
        décrit à la section~\ref{sec:couplage}.
\end{itemize}
Le résidu peut être interprété directement comme une mesure d'incompatibilité entre la solution hydraulique
produite par un débit d'essai et la condition de pression aval : \(F(Q)<0\) indique une pression calculée
insuffisante, \(F(Q)=0\) une compatibilité, \(F(Q)>0\) une pression calculée excédentaire.

Un point mérite d'être ajouté dans la perspective des effets d'entrée analysés au
paragraphe~\ref{par:entree}. Par symétrie exacte avec~\eqref{eq:pout}, un coefficient de perte singulière
d'entrée \(\kin\) peut être introduit en écrivant
\begin{equation}
p(0,t) = P_{\mathrm{in}} + \kin\,\frac{1}{2}\rho(0)\,u(0;Q)^2 ,
\label{eq:pin}
\end{equation}
ce qui revient à modifier le résidu en
\begin{equation}
F(Q) = p_{\mathrm{calc}}(L;Q) - p_{\mathrm{out}}(Q)
+ \kin\,\frac{1}{2}\rho(0)\,u(0;Q)^2 .
\label{eq:FQ-kin}
\end{equation}
L'introduction de ce terme ne change pas la structure du problème scalaire ni la nature du solveur, mais
elle modifie la relation débit--pression et donc les temps d'érosion. Il est recommandé de l'introduire
lorsque la charge cinétique d'entrée n'est pas négligeable devant la charge motrice, ce qui est
précisément le cas pour les configurations à forte vitesse rapportées au tableau~\ref{tab:sansdim}.

\subsection{Résolution par encadrement et dichotomie}
\label{sec:bracketing}

\paragraph{Principe de l'encadrement.} La première étape consiste à identifier deux valeurs de débit
\(Q_a<Q_b\) telles que \(F(Q_a)F(Q_b)<0\). Sous hypothèse de continuité du résidu, cette condition garantit
l'existence d'au moins une racine \(Q^*\in[Q_a,Q_b]\). En pratique, l'intervalle initial est élargi
progressivement, par exemple par \(Q_b \leftarrow 2Q_b\), jusqu'à obtenir le changement de signe. La
procédure ne dépend donc pas d'une estimation initiale précise du débit, ce qui est essentiel dans un
contexte où \(Q\) varie de plusieurs ordres de grandeur au cours d'un même calcul.

\paragraph{Itération de dichotomie.} Une fois l'intervalle établi, on calcule \(Q_m = (Q_a+Q_b)/2\) et on
évalue \(F(Q_m)\). Si \(F(Q_a)F(Q_m)<0\), la racine reste dans \([Q_a,Q_m]\) ; sinon elle est dans
\([Q_m,Q_b]\). L'opération est répétée jusqu'à satisfaction du critère d'arrêt. La largeur de l'intervalle
décroît géométriquement,
\begin{equation}
\Delta Q^{(k)} = \frac{\Delta Q^{(0)}}{2^{k}},
\label{eq:bisection-decay}
\end{equation}
de sorte que, après \(k\) itérations, l'erreur sur le débit est bornée par
\begin{equation}
\left|Q^{(k)}-Q^*\right| \le \frac{Q_b^{(0)}-Q_a^{(0)}}{2^{k+1}} .
\label{eq:bisection-borne}
\end{equation}
Cette borne est une garantie \emph{a priori}, indépendante de la pente locale du résidu.

\paragraph{Critère d'arrêt fondé sur la pression.} Le critère de résolution est exprimé directement en
termes de pression,
\begin{equation}
\left|F(Q^{(k)})\right| = \left|p_{\mathrm{calc}}(L;Q^{(k)}) - p_{\mathrm{out}}(Q^{(k)})\right|
\le \varepsilon_p ,
\label{eq:tol-p}
\end{equation}
ce qui est cohérent avec la formulation du problème : la grandeur imposée à la frontière est la pression et
non le débit. Le débit est la variable recherchée, la pression calculée à la sortie est la mesure de
qualité de la solution. Ce critère en pression est complété en pratique par une borne sur l'intervalle
\(Q_b-Q_a \le \varepsilon_Q\) et par un nombre maximal d'itérations, qui protègent le calcul contre les
situations où la sensibilité \(\mathrm{d}p(L)/\mathrm{d}Q\) devient très faible en fin d'emballement ; ces
garde-fous doivent être documentés car ils déterminent la précision effective du solveur en régime
fortement érodé.

\paragraph{Déroulement du solveur.} Le fonctionnement peut être résumé comme suit.
\begin{enumerate}[label=\arabic*.]
  \item Choisir un intervalle initial \([Q_a, Q_b]\).
  \item Évaluer \(F(Q_a)\) et \(F(Q_b)\).
  \item Si \(F(Q_a)F(Q_b) \ge 0\), élargir l'intervalle et répéter l'étape 2.
  \item Calculer \(Q_m = (Q_a+Q_b)/2\) et évaluer \(F(Q_m)\).
  \item Si \(|F(Q_m)| \le \varepsilon_p\) ou si les bornes d'intervalle sont atteintes, retourner \(Q_m\).
  \item Sinon, conserver le sous-intervalle contenant la racine et répéter à partir de l'étape 4.
\end{enumerate}

\subsection{Monotonie du résidu et robustesse du solveur}
\label{sec:robustesse}

\paragraph{Monotonie.} Le caractère monotone de \(F(Q)\) sur la plage de paramètres étudiée constitue le
fondement du choix de la méthode. Pour un état géométrique et un état de concentration fixés, une
augmentation du débit accroît la vitesse, donc le nombre de Reynolds, le frottement et la contrainte
pariétale, ce qui accroît la perte de charge totale et diminue la pression calculée à l'exutoire. Cette
dépendance conduit à une variation monotone du résidu, ce qui garantit l'unicité de la racine encadrée.
La propriété est vérifiée numériquement \emph{a posteriori} sur chaque calcul, en contrôlant que
\(F(Q_a)\) et \(F(Q_b)\) restent de signes opposés avec \(F\) décroissante ; elle n'est pas postulée.

\paragraph{Avantage fondamental de l'encadrement.} La propriété essentielle de la dichotomie est qu'elle
conserve explicitement un intervalle contenant la racine, et cela \emph{par construction} et à chaque
itération : l'algorithme ne dépend jamais d'une approximation locale de la fonction autour d'un point
unique. À l'inverse, la méthode de Newton--Raphson repose sur le développement local
\begin{equation}
F(Q) \approx F(Q_n) + F'(Q_n)\left(Q-Q_n\right)
\qquad\Longrightarrow\qquad
Q_{n+1} = Q_n - \frac{F(Q_n)}{F'(Q_n)},
\label{eq:newton}
\end{equation}
dont la convergence est quadratique lorsque les hypothèses locales sont favorables, mais qui n'offre aucune
garantie \emph{a priori} en dehors de ces hypothèses. Il convient donc d'être précis : la dichotomie n'est
pas plus rapide que Newton, et elle ne lui est pas supérieure dans l'absolu ; elle est préférable ici
parce que la convergence par réduction d'intervalle est garantie dès que la racine est encadrée et le
résidu continu, alors que la convergence de Newton dépend du comportement local de la fonction.

\paragraph{Absence de dérivée.} Un avantage complémentaire est l'absence de toute évaluation de \(F'(Q)\).
Dans le problème considéré, chaque évaluation de \(F\) implique déjà une intégration complète du champ
hydraulique sur les \(N_x\) cellules. Une approximation numérique de la dérivée exigerait au moins une
évaluation supplémentaire par itération,
\begin{equation}
F'(Q) \approx \frac{F(Q+\delta Q) - F(Q)}{\delta Q},
\label{eq:derivee-num}
\end{equation}
et le choix de \(\delta Q\) deviendrait critique lorsque le résidu est très raide, c'est-à-dire précisément
dans les régimes où le solveur doit rester fiable. La dichotomie n'introduit pas ce paramètre
supplémentaire.

\paragraph{Robustesse vis-à-vis de l'évolution du couplage.} Le problème hydraulique n'est pas résolu une
seule fois : à chaque pas, le changement d'état \((R^n,\phi^n)\to(R^{n+1},\phi^{n+1})\) produit une nouvelle
fonction de résidu \(F_n(Q) \to F_{n+1}(Q)\). Une méthode exigeant une excellente estimation initiale à
chaque appel serait fragile ; une méthode à encadrement conserve un comportement contrôlé et reproductible
à travers les transitions rapides. Cette propriété a une conséquence méthodologique importante : si le
solveur hydraulique a un comportement contrôlé, alors une évolution brutale de \(R\) ou de \(\phi\) peut
être attribuée aux équations physiques plutôt qu'à une convergence défaillante du calcul du débit. La
distinction entre instabilité physique et échec numérique repose donc en partie sur la robustesse du
solveur.

\paragraph{Coût et compromis.} La contrepartie de la dichotomie est son coût : la convergence est
linéaire, la largeur de l'intervalle étant divisée par deux à chaque itération, contre une convergence
quadratique pour Newton dans des conditions favorables. Ce surcoût est explicitement accepté. Le solveur
est appelé une fois par pas de temps, sur des milliers de pas, et sert également de brique de base aux
études paramétriques ; la priorité est donc la fiabilité globale du calcul couplé plutôt que la
minimisation du nombre d'itérations de chaque résolution individuelle. Un échec du solveur hydraulique
interromprait l'ensemble de la chaîne de calcul, ce qui justifie ce compromis.

\subsection{Un résultat analytique : emballement structurel et temps caractéristique}
\label{sec:instabilite}

La simplicité de l'équation~\eqref{eq:quasi} permet d'obtenir un résultat exact qui éclaire l'ensemble des
résultats numériques et fournit un cas de vérification analytique. Considérons une géométrie uniforme,
\(R(x,t) = R(t)\). L'équation~\eqref{eq:quasi} intégrée entre l'entrée et la sortie donne alors, en valeur
absolue,
\begin{equation}
\frac{\Delta p}{L} = \frac{2\tau_b}{R}
\qquad\Longrightarrow\qquad
\tau_b(t) = \frac{\Delta p}{2L}\,R(t),
\label{eq:tau-analytique}
\end{equation}
résultat remarquable par sa généralité : la contrainte pariétale est fixée par la géométrie et par la
charge imposée, \emph{indépendamment de la loi de frottement}. La fermeture hydraulique ne détermine donc
pas \(\tau_b\) mais le débit \(Q\) (et donc \(u\) et \(Re\)) qui permet de réaliser cette contrainte.
Physiquement, si le conduit s'élargit, la résistance diminue et le débit augmente précisément de manière à
maintenir la contrainte à la valeur imposée par le gradient de pression.

En reportant~\eqref{eq:tau-analytique} dans la loi d'érosion~\eqref{eq:erosion} et
l'équation~\eqref{eq:dRdt}, et en négligeant le facteur géométrique \(\sqrt{1+(\partial_xR)^2}\) pour une
géométrie uniforme, on obtient une équation différentielle ordinaire linéaire en \(R\) :
\begin{equation}
\frac{\mathrm{d}R}{\mathrm{d}t}
= \frac{k_{er}}{\rho_s}\left(\frac{\Delta p}{2L}R - \tau_c\right)
\qquad\Longrightarrow\qquad
R(t) = R_c + \left(R_0 - R_c\right)e^{t/\ter},
\label{eq:R-exponentielle}
\end{equation}
avec
\begin{equation}
R_c = \frac{2L\tau_c}{\Delta p},
\qquad
\ter = \frac{2L\rho_s}{k_{er}\Delta p}.
\label{eq:Rc-ter}
\end{equation}
Trois conséquences méritent d'être soulignées.

\begin{enumerate}
  \item \textbf{L'emballement est structurel.} Le taux de croissance est strictement positif tant que
        \(R>\tau_c 2L/\Delta p\), et \(R\) croît exponentiellement sans borne finie. Il n'existe donc pas
        de point fixe géométrique stable sous pression imposée en l'absence d'une résistance aval : le
        régime d'emballement est la conséquence directe de la condition aux limites, et non de la
        fermeture de frottement ni de la rhéologie du mélange. Ce résultat analytique justifie et
        généralise l'observation numérique d'instabilité systématique.
  \item \textbf{\(\ter\) est un taux de croissance, pas seulement une échelle de normalisation.} Le temps
        de doublement du rayon vaut \(t_2 = \ter\ln 2\), indépendamment de \(R_0\) et de \(\tau_c\) dès
        que \(R_0 \gg R_c\). Pour la configuration de la campagne A (\(\ter = 949.4\) s), on obtient
        \(t_2 = 658\) s, à comparer aux \(665\) s obtenus par le calcul couplé complet : l'accord à
        \(1\%\) valide simultanément l'implémentation du solveur hydraulique et la cohérence du temps
        caractéristique retenu.
  \item \textbf{Le rôle stabilisateur d'une résistance aval se lit directement.} Avec une perte singulière
        \(\kout\), la contrainte pariétale devient \(\tau_b = R\left[\Delta p - \kout\frac{1}{2}\rho u^2\right]/(2L)\),
        de sorte que l'augmentation du débit associée à l'élargissement produit une contre-pression :
        c'est le mécanisme exact par lequel \(\kout>0\) ralentit l'érosion, retardant le temps de
        doublement sans pour autant annuler la croissance à long terme.
\end{enumerate}

Deux confirmations numériques de ces relations méritent d'être mentionnées, car elles valident la
reconstruction du modèle utilisée pour l'interprétation. D'une part, le calcul direct du champ hydraulique
en géométrie uniforme vérifie la relation \eqref{eq:tau-analytique} à mieux de \(10^{-6}\) relatif pour
tous les rayons testés (de \(R_0\) à \(10R_0\)), ce qui confirme l'indépendance de \(\tau_b\) vis-à-vis
de la fermeture. D'autre part, le débit suit exactement la loi de puissance
\(Q \propto R^{19/7}\) (pente logarithmique \(2.714\) à mieux de \(10^{-3}\) pour
\(R/R_0 = 1.5\) à \(10\)), qui quantifie la force de la rétroaction \eqref{eq:retro-positive} : un
doublement du rayon multiplie le débit par \(6.6\) et la contrainte pariétale par \(2\), ce qui explique
la brutalité de l'emballement observé numériquement.

Ce résultat fournit également un critère de vérification indépendant. Pour toute configuration à géométrie
initialement uniforme, le calcul couplé doit reproduire (i) la contrainte analytique~\eqref{eq:tau-analytique}
au premier pas et (ii) la loi exponentielle~\eqref{eq:R-exponentielle} pour l'intervalle pendant lequel la
géométrie reste quasi uniforme. Le tableau~\ref{tab:verif-analytique} présente ce contrôle pour la
configuration de référence.

\begin{table}[H]
\centering
\caption{Vérification analytique de l'état initial et du taux de croissance. Comparaison entre la
prédiction analytique~\eqref{eq:tau-analytique}--\eqref{eq:Rc-ter} et le calcul couplé, pour la
configuration de la campagne A.}
\label{tab:verif-analytique}
\small
\begin{tabular}{llrl}
\toprule
Grandeur & Expression analytique & Valeur analytique & Valeur calculée \\
\midrule
Vitesse initiale & \(u_0 = Q_0/(\pi R_0^2)\) & --- & \(4.5\) m/s \\
Contrainte pariétale initiale & \(\tau_b = R_0\Delta p/(2L)\) & \(62.9\) Pa & \(62.9\) Pa \\
Débit initial & \(Q_0\) & --- & \(1.3\times10^{-4}\ \mathrm{m^3/s}\) \\
Temps de doublement & \(t_2 = \ter\ln 2\) & \(658\) s & \(665\) s \\
Rayon critique & \(R_c = 2L\tau_c/\Delta p\) & \(0.048\) mm & --- \\
\bottomrule
\end{tabular}
\end{table}

% TODO auteur : verifier la colonne "valeur calculee" a partir des sorties effectives (etat initial de
% la simulation de reference) et completer la ligne du temps de doublement si la definition de t_2
% differe (doublement du rayon a la section de sortie et non a l'entree).


%=====================================================================================
\section{Discrétisation numérique du transport advectif des particules}
\label{sec:transport}
%=====================================================================================

La modélisation du transport des particules érodées constitue une composante déterminante du modèle
couplé, car la concentration en solides \(\phi(x,t)\) intervient directement dans la modification des
propriétés hydrauliques du mélange. Le transport est décrit par l'équation d'advection--relaxation
\eqref{eq:transport}, dans laquelle le terme advectif assure la propagation longitudinale des particules
tandis que le terme source traduit leur alimentation par l'érosion des parois. La résolution numérique de
cette équation doit donc simultanément préserver la conservation de la quantité transportée et rester
suffisamment précise au voisinage des gradients élevés susceptibles d'apparaître dans le champ de
concentration.

La difficulté numérique provient principalement du caractère fortement advectif du problème : les
simulations mettent en évidence la formation de fronts de concentration abrupts se propageant vers l'aval
du conduit. Ces fronts présentent des variations spatiales beaucoup plus rapides que celles observées dans
les champs de pression ou de rayon et constituent, par conséquent, la principale contrainte pour la
discrétisation spatiale. Le schéma retenu est une méthode conservative de volumes finis centrée sur les
cellules, associée à une reconstruction MUSCL et à un limiteur de pente. Cette section établit la
formulation conservative, développe la reconstruction MUSCL, précise le rôle du limiteur, puis compare le
schéma retenu aux deux extrêmes que constituent le schéma amont du premier ordre et les reconstructions
non limitées.

\subsection{Formulation conservative et discrétisation par volumes finis}

\paragraph{Structure conservative de l'équation.} La forme retenue pour le terme advectif est la forme
conservative \(\partial(u\phi)/\partial x\), et non l'écriture développée \(u\,\partial\phi/\partial x\).
Cette distinction est fondamentale pour la discrétisation. Pour un volume de contrôle
\(C_i = [x_{i-1/2}, x_{i+1/2}]\), l'intégration de~\eqref{eq:transport} donne en effet
\begin{equation}
\frac{\mathrm{d}}{\mathrm{d}t}\int_{C_i}\phi\,\mathrm{d}x
+ (u\phi)_{i+1/2} - (u\phi)_{i-1/2}
= \int_{C_i} k\left(\phisol-\phi\right)\mathrm{d}x .
\label{eq:bilan-integral}
\end{equation}
Cette relation exprime directement le bilan sur la cellule : la variation de la quantité de solides
contenue dans la cellule est déterminée par la différence entre les flux entrants et sortants, augmentée
de la quantité produite localement par l'érosion. La conservation ne dépend donc pas d'une reconstruction
artificielle d'une dérivée point par point ; elle résulte du bilan intégral, et le flux traversant une
interface apparaît avec des signes opposés dans les bilans des deux cellules adjacentes.

\paragraph{Discrétisation centrée sur les cellules.} Le domaine est subdivisé en \(N_x\) cellules de
largeur \(\Delta x = L/N_x\), et l'inconnue \(\phi_i(t)\) est définie comme la moyenne cellulaire
\begin{equation}
\phi_i(t) = \frac{1}{\Delta x}\int_{x_{i-1/2}}^{x_{i+1/2}}\phi(x,t)\,\mathrm{d}x .
\label{eq:moyenne-cellulaire}
\end{equation}
La forme semi-discrète s'écrit alors
\begin{equation}
\frac{\mathrm{d}\phi_i}{\mathrm{d}t}
= -\frac{F_{i+1/2} - F_{i-1/2}}{\Delta x} + S_i,
\qquad
F_{i+1/2} = (u\phi)_{i+1/2},
\label{eq:semi-discret}
\end{equation}
où \(S_i\) est la moyenne cellulaire du terme source. Toute amélioration du traitement du transport est
introduite au niveau de la définition du flux \(F_{i+1/2}\), sans modifier la structure conservative du
bilan.

\paragraph{Avancement explicite et conservation discrète.} En isolant provisoirement le terme source,
l'étape advective s'écrit
\begin{equation}
\phi_i^* = \phi_i^n - \frac{\dt}{\Delta x}\left(F_{i+1/2}^n - F_{i-1/2}^n\right).
\label{eq:update-fv}
\end{equation}
La somme des contributions de flux internes s'annule lors d'une sommation sur toutes les cellules : les
seules variations globales proviennent des flux traversant les frontières du domaine. Dans l'équation,
cette propriété revêt une importance particulière, car une erreur de conservation du transport se
propagerait vers le calcul hydraulique par l'intermédiaire de \(\rho(\phi)\) et de \(f_m(\phi)\), puis vers
le taux d'érosion par l'intermédiaire de \(\tau_b\). La conservation du transport est donc une condition
nécessaire non seulement pour la qualité du champ \(\phi\), mais également pour la cohérence du couplage
multiphysique dans son ensemble.

\subsection{Reconstruction MUSCL et construction du flux interfacial}

\paragraph{Définition de la reconstruction.} La formulation conservative fixe la structure du bilan mais ne
détermine pas la manière dont le flux interfacial \(F_{i+1/2}\) doit être évalué à partir des moyennes
cellulaires. Pour un schéma de premier ordre, la valeur d'interface est prise égale à la valeur cellulaire :
\(\phi_{i+1/2} \approx \phi_i\), ce qui conduit à une erreur de troncature en \(\mathcal{O}(\Delta x)\).
La reconstruction MUSCL (Monotonic Upstream-centered Scheme for Conservation Laws) consiste à reconstruire
dans chaque cellule un \emph{profil linéaire} au lieu d'un profil constant,
\begin{equation}
\phi(x,t^n) \simeq \phi_i^n + s_i^n\left(x - x_i\right),
\qquad x \in C_i ,
\label{eq:reconstruction-lineaire}
\end{equation}
où \(s_i^n\) est une approximation limitée de la pente de \(\phi\) dans la cellule \(i\). Les états
reconstruits de part et d'autre d'une interface \(x_{i+1/2}\) sont alors
\begin{equation}
\phi_{i+1/2}^{L} = \phi_i + \frac{1}{2}s_i,
\qquad
\phi_{i+1/2}^{R} = \phi_{i+1} - \frac{1}{2}s_{i+1},
\label{eq:etats-interface}
\end{equation}
le premier étant associé à la cellule amont immédiate pour un écoulement dirigé vers les \(x\) croissants.
Le flux numérique est ensuite construit par un flux de Godunov pour l'équation d'advection linéaire, qui
se réduit au décentrage selon le signe de la vitesse :
\begin{equation}
F_{i+1/2} = a_{i+1/2}\,
\begin{cases}
\phi_{i+1/2}^{L}, & a_{i+1/2} \ge 0,\\[1mm]
\phi_{i+1/2}^{R}, & a_{i+1/2} < 0,
\end{cases}
\qquad
a_{i+1/2} = \beta_u\,\frac{u_i + u_{i+1}}{2}.
\label{eq:flux-godunov}
\end{equation}
Cette construction est cohérente : elle reproduit le schéma amont lorsque la pente est nulle et se réduit à
un schéma d'ordre deux lorsque la pente est exacte.

\paragraph{Rôle du limiteur : du rapport de pentes à la pente effective.} La pente \(s_i\) ne peut pas être
choisie librement. Une reconstruction linéaire non contrôlée produit, au voisinage des fronts, des
oscillations non physiques et peut conduire à des valeurs négatives de la concentration. La pente est donc
exprimée sous la forme \(s_i = \Psi(r_i)\,\Delta_{R,i}\), où
\begin{equation}
\Delta_{L,i} = \phi_i - \phi_{i-1},
\qquad
\Delta_{R,i} = \phi_{i+1} - \phi_i,
\qquad
r_i = \frac{\Delta_{L,i}}{\Delta_{R,i}},
\label{eq:rapport-pentes}
\end{equation}
et où \(\Psi\) est une fonction \emph{limiteur}. Toutes les fonctions de limitation usuelles partagent la
propriété
\begin{equation}
\Psi(r) = 0 \quad \text{pour } r \le 0,
\label{eq:psi-nul}
\end{equation}
c'est-à-dire qu'elles annulent la pente aux extrema locaux et aux discontinuités. C'est cette propriété
qui distingue essentiellement le schéma limité d'une reconstruction centrale non limitée, et qui explique
la perte locale d'ordre deux aux extrema : le schéma devient localement d'ordre un précisément là où
l'ordre deux produirait un dépassement.

\paragraph{Condition TVD et fonctions de limitation usuelles.} La condition assurant qu'un schéma
\(\phi_i^{n+1} = H(\phi^{n}_{i-1},\phi^n_i,\phi^n_{i+1})\) est \emph{Total Variation Diminishing}
s'écrit, sous condition CFL adéquate,
\begin{equation}
\mathrm{TV}(\phi^{n+1}) \le \mathrm{TV}(\phi^{n}),
\qquad
\mathrm{TV}(\phi) = \sum_i\left|\phi_{i+1}-\phi_i\right| .
\label{eq:tvd}
\end{equation}
Pour la forme~\eqref{eq:reconstruction-lineaire}--\eqref{eq:flux-godunov}, la condition suffisante de
stabilité est
\begin{equation}
0 \le \Psi(r) \le \min\left(2r,\,2\right),
\label{eq:condition-psi}
\end{equation}
ce qui définit la région admissible des limiteurs. Les trois fonctions les plus utilisées sont
\begin{align}
\text{minmod} &: \quad \Psi(r) = \max\left(0,\min\left(r,1\right)\right),
\label{eq:minmod}\\[1mm]
\text{van Leer} &: \quad \Psi(r) = \frac{r + |r|}{1 + |r|},
\label{eq:vanleer}\\[1mm]
\text{Monotonised Central (MC)} &: \quad \Psi_{\mathrm{MC}}(r)
= \max\left(0,\min\left(2r,\tfrac{1}{2}(1+r),2\right)\right).
\label{eq:mc}
\end{align}
Le limiteur minmod est le plus restrictif : il est le plus diffusif des trois et conserve la propriété TVD
pour tout nombre de Courant. Le limiteur de van Leer occupe une position intermédiaire et présente
l'avantage d'être une fonction régulière (\(\mathcal{C}^1\)) de \(r\), ce qui évite les discontinuités de
pente qui peuvent exciter le couplage avec le facteur de résistance \(f_m(\phi)\). Le limiteur MC est le
moins diffusif des trois tout en satisfaisant~\eqref{eq:condition-psi} : il correspond à la plus grande
pente admissible dans la région TVD, ce qui explique qu'il soit fréquemment retenu pour les fronts raides.

\subsection{Propriété TVD et seuil de stabilité du schéma retenu}

La condition~\eqref{eq:condition-psi} garantit la décroissance de la variation totale, mais elle ne fixe
pas à elle seule la limite de stabilité temporelle du schéma. Pour la forme reconstruire--résoudre--avancer
utilisée ici, c'est-à-dire une reconstruction linéaire limitée suivie d'une avancée explicite d'Euler, la
propriété TVD est établie théoriquement sous la condition plus restrictive
\begin{equation}
\mathrm{CFL} \le \frac{1}{2},
\label{eq:cfl-half}
\end{equation}
et non sous la condition \(\mathrm{CFL} \le 1\) valable pour les schémas de type flux-limiteur \og
Lax--Wendroff limité \fg{} (Sweby, 1984 ; LeVeque, 2002 ; Toro, 2009). Cette distinction est importante en
pratique, car c'est la valeur \eqref{eq:cfl-half} qui conditionne le réglage du pas de temps.

Une vérification numérique directe de ce seuil a été effectuée sur l'équation d'advection linéaire avec un
état initial en marche et une période complète, pour les deux limiteurs retenus et pour les deux variantes
du schéma. Le tableau~\ref{tab:seuil-cfl} en présente les résultats.

\begin{table}[H]
\centering
\caption{Seuil de stabilité observé pour le schéma d'advection (équation d'advection linéaire, état
initial en marche, \(N_x = 120\), 200 pas). \og OK \fg{} indique un schéma stable, sans dépassement
\(\max(\phi) > 1.02\) ni valeur négative \(\min(\phi) < 0\). \og Échec \fg{} indique l'apparition
d'oscillations divergentes.}
\label{tab:seuil-cfl}
\small
\begin{tabular}{llccccccc}
\toprule
Limiteur & Variante & 0.30 & 0.40 & \textbf{0.50} & 0.55 & 0.70 & 0.90 & 0.99 \\
\midrule
MC & reconstruire--résoudre--avancer & OK & OK & \textbf{OK} & Échec & Échec & Échec & Échec \\
MC & MUSCL--Hancock (prédicteur) & OK & OK & \textbf{OK} & OK & OK & OK & OK \\
van Leer & reconstruire--résoudre--avancer & OK & OK & \textbf{OK} & Échec & Échec & Échec & Échec \\
van Leer & MUSCL--Hancock (prédicteur) & OK & OK & \textbf{OK} & OK & OK & OK & OK \\
\bottomrule
\end{tabular}
\begin{tablenotes}
\small
\item Le seuil observé pour la variante reconstruire--résoudre--avancer est
\(\mathrm{CFL} \approx 0.5\), indépendamment du limiteur retenu, en accord avec~\eqref{eq:cfl-half}.
L'ajout d'un prédicteur de type MUSCL--Hancock (avancement des états d'interface d'un demi-pas avant le
calcul du flux) repousse ce seuil à \(\mathrm{CFL} \approx 1\), conformément à la théorie
(Sweby, 1984).
\end{tablenotes}
\end{table}

Deux conséquences pratiques en découlent. Premièrement, le nombre de Courant retenu pour les simulations
présentées dans ce travail, \(\mathrm{CFL}_{\mathrm{tgt}} = 0.35\), se situe à l'intérieur du domaine de
stabilité avec une marge confortable, ce qui justifie \emph{a posteriori} le choix de cette valeur, et non
un simple empirisme. Le réglage n'est donc pas arbitraire : il est à la fois suffisamment proche du seuil
pour ne pas dégrader inutilement la résolution temporelle, et suffisamment éloigné pour absorber les
variations locales du pas de temps. Deuxièmement, si l'on souhaitait augmenter \(\mathrm{CFL}\) et donc
accélérer le calcul, la voie correcte n'est pas d'augmenter le nombre de Courant du schéma actuel, mais
d'ajouter un prédicteur de type Hancock ; cette possibilité est mentionnée comme perspective
d'optimisation et non comme un réglage disponible.

\subsection{Comparaison avec le schéma amont du premier ordre et les reconstructions non limitées}

\paragraph{Schéma amont du premier ordre : diffusion numérique.} Le schéma amont est obtenu en posant
\(s_i \equiv 0\) dans~\eqref{eq:reconstruction-lineaire}, soit \(F_{i+1/2} = a_{i+1/2}\phi_i\) pour
\(a \ge 0\). Il est inconditionnellement positif et TVD pour \(\mathrm{CFL} \le 1\), mais introduit une
diffusion numérique explicite, dont le coefficient est donné par l'équation modifiée du schéma,
\begin{equation}
\Disc = \frac{|a|\,\Delta x}{2}\left(1 - \mathrm{CFL}\right).
\label{eq:diffnum}
\end{equation}
Cette diffusion n'est pas un artefact mineur : elle étale le front proportionnellement à
\(\sqrt{\Disc\,t}\). Pour la configuration de référence de la campagne A (\(u \approx 4.5\) m/s,
\(\Delta x = 1.95\times10^{-4}\) m pour \(N_x = 600\), \(\mathrm{CFL} = 0.35\)), on obtient
\(\Disc \approx 2.9\times10^{-4}\ \mathrm{m^2/s}\). Au bout d'une demi-unité de temps caractéristique
(\(t = 0.5\,\ter \approx 475\) s), l'étalement caractéristique du front atteint
\(\sqrt{\Disc\,t} \approx 0.37\) m, soit environ \emph{trois fois} la longueur du conduit. Autrement dit,
un schéma amont aurait, à l'échelle d'une fraction de \(\ter\), diffusé la totalité de la structure
spatiale que le modèle cherche précisément à décrire.

Le tableau~\ref{tab:fronts} confirme ce diagnostic sur le test d'advection linéaire : le schéma amont
étale le front sur environ le quart du domaine, contre quatre à six cellules pour les schémas limités,
et sa variation totale reste égale à sa valeur initiale (absence d'oscillation, mais diffusion maximale).

\begin{table}[H]
\centering
\caption{Comparaison quantitative des schémas d'advection sur un test de référence : advection
linéaire d'un créneau de largeur \(0.3L\), \(\mathrm{CFL} = 0.5\), \(N_x = 120\) à \(\Delta x\) fixé, après
un déplacement correspondant à un nombre entier de cellules. Les erreurs sont calculées contre la
solution exacte translatée ; la largeur de front est le nombre de cellules telles que
\(0.05 < \phi < 0.95\) (deux fronts, donc un minimum de 4 cellules).}
\label{tab:fronts}
\small
\begin{tabular}{llccccc}
\toprule
Schéma & Limiteur & Largeur de front & \(\min\phi\) & \(\mathrm{TV}/\mathrm{TV}_0\)
& \(\varepsilon_{L_1}\) & \(\varepsilon_{L_\infty}\) \\
\midrule
Amont (ordre 1) & --- & \(24\) & \(\ge 0\) & \(1.00\) & \(5.1\times10^{-2}\) & \(4.5\times10^{-1}\) \\
MUSCL & minmod & \(4\) & \(\ge 0\) & \(1.00\) & \(8.3\times10^{-3}\) & \(2.1\times10^{-1}\) \\
MUSCL & van Leer & \(4\) & \(\ge 0\) & \(1.00\) & \(4.6\times10^{-3}\) & \(1.4\times10^{-1}\) \\
MUSCL & MC & \(4\) & \(\ge 0\) & \(1.00\) & \(4.2\times10^{-3}\) & \(1.3\times10^{-1}\) \\
Reconstruction linéaire non limitée & --- & oscillations & \(<0\) & \(\gg 1\) & --- & --- \\
\bottomrule
\end{tabular}
\begin{tablenotes}
\small
\item À résolution égale, le schéma limité réduit l'erreur \(L_1\) d'un facteur 6 à 12 par rapport au
schéma amont, pour une largeur de front divisée par six. La reconstruction non limitée devient
instable et produit des valeurs négatives, la valeur de \(\mathrm{TV}/\mathrm{TV}_0\) n'ayant alors
plus de sens.
\end{tablenotes}
\end{table}

\paragraph{Reconstructions non limitées : oscillations et perte de positivité.} En l'absence de
limiteur, la pente est estimée par une formule centrée, \(s_i = \frac{1}{2}(\phi_{i+1}-\phi_{i-1})\),
et le schéma devient une variante de type Lax--Wendroff. Les tests du tableau~\ref{tab:seuil-cfl} montrent
que cette variante n'est pas seulement oscillante : elle produit des valeurs \emph{négatives} de la
concentration dès les premiers pas, et ce même pour un nombre de Courant faible
(\(\mathrm{CFL} = 0.25\)). Or une concentration négative n'a pas de sens physique et, dans le modèle
couplé, une valeur négative de \(\phi\) se propagerait vers \(f_m(\phi)\) et \(\lambda(\phi)\), où la
fonction \(\lambda\) n'est définie que pour \(\phi \in\,]0,\phisol[\). La non-linéarité de la fermeture
rhéologique transforme donc immédiatement une oscillation numérique du transport en erreur irrécupérable
sur l'hydraulique. C'est cette raison, et non une préférence esthétique, qui impose le recours à un
limiteur.

\subsection{Contraintes de positivité et de monotonie de la concentration}
\label{sec:positivite}

Le champ \(\phi\) doit satisfaire deux contraintes physiques : la positivité, \(\phi \ge 0\), et la borne
supérieure \(\phi \le \phisol\), qui traduit l'impossibilité pour la suspension de dépasser la fraction
solide du matériau intact.

\paragraph{Borne supérieure.} La contrainte \(\phi \le \phisol\) est assurée par la structure même du
traitement du terme source, puisque la solution exacte~\eqref{eq:source-exacte} est une combinaison
convexe entre l'état post-advectif et \(\phisol\) :
\begin{equation}
\phi_i^{n+1} = \phisol - \left(\phisol - \phi_i^*\right)e^{-k_i^n\dt},
\qquad\Longrightarrow\qquad
\phi_i^{n+1} \le \max\left(\phisol,\, \phi_i^*\right).
\label{eq:borne-sup}
\end{equation}
Cette propriété est une conséquence directe de l'exposant négatif et ne dépend pas de la valeur de
\(k\dt\), ce qui explique l'intérêt de l'intégration analytique pour les régimes raides où
\(k\dt \gg 1\). Elle n'est toutefois valable que si \(\phi_i^*\) respecte lui-même la borne : c'est
l'étape d'advection qui doit la préserver.

\paragraph{Positivité à l'étape d'advection.} La reconstruction et le flux limités préservent la
positivité sous deux conditions : (i) les valeurs reconstruites aux interfaces restent encadrées par les
moyennes cellulaires voisines, propriété satisfaite par construction par les limiteurs de
type~\eqref{eq:condition-psi} puisque \(|s_i| \le 2\min(|\Delta_L|,|\Delta_R|)\) ; et (ii) le nombre de
Courant reste dans le domaine de stabilité \eqref{eq:cfl-half}. Les résultats du
tableau~\ref{tab:seuil-cfl} confirment ce comportement : pour \(\mathrm{CFL} \le 0.5\), le schéma ne
produit aucune valeur négative, quelle que soit la position du front, alors qu'au-delà du seuil
l'oscillation apparaît immédiatement.

\paragraph{Implication pour le réglage de \(\Delta t_{\max}\).} Cette analyse fournit un critère de
cohérence supplémentaire pour le réglage du pas de temps. Si le nombre de Courant effectif venait à
dépasser durablement \(0.5\) — par exemple parce qu'une borne \(\dt_{\max}\) trop élevée prendrait le
relais du contrôle CFL —, la positivité du champ \(\phi\) ne serait plus garantie. Le diagnostic
d'effectivité présenté à la section~\ref{sec:cfl-diagnostic} doit donc être complété par un contrôle du
nombre de Courant effectif maximal, et non de sa seule moyenne.

\paragraph{Contrôle pratique et défaut de masse.} Dans la pratique, une troncature
\(\phi \leftarrow \max(\phi,0)\) est appliquée après l'étape d'advection, afin de prémunir le calcul contre
toute excursion résiduelle. Cette opération doit être documentée, car elle n'est pas conservative :
la masse retirée doit être cumulée et rapportée, et elle doit rester négligeable devant la masse
transitée. Un contrôle systématique consiste à suivre le rapport
\begin{equation}
\delta_m(t) = \frac{\displaystyle\sum_i\max\left(-\phi_i^*,0\right)\Delta x}
{\displaystyle\sum_i\phi_i^n\,\Delta x},
\label{eq:defaut-masse}
\end{equation}
dont la valeur doit rester inférieure à la tolérance de conservation retenue pour la vérification. Ce
contrôle complète le diagnostic CFL et fournit une mesure quantitative de la qualité du couplage
transport--hydraulique.

% TODO auteur : faire apparaitre explicitement (i) le limiteur unique retenu pour les resultats publies
% (van Leer ou MC) et (ii) la valeur du defaut de masse (eq. 4.xx) mesuree sur les simulations de reference.


%=====================================================================================
\section{Intégration du terme d'érosion et couplage des sous-modèles}
\label{sec:source-erosion}
%=====================================================================================

\subsection{Décomposition opératoirelle du problème}

Le terme source de l'équation~\eqref{eq:transport} présente une nature mathématique différente de celle du
terme advectif. Le terme advectif est un opérateur de transport non local, dont la discrétisation exige un
traitement conservateur et une reconstruction haute résolution ; le terme source est un opérateur
\emph{local} de relaxation,
\begin{equation}
S_\phi(\phi) = k\left(\phisol - \phi\right),
\label{eq:opsrc}
\end{equation}
dont la raideur est fixée par la valeur locale de \(k = 2\mdot/(R\rho_s)\). Cette différence de nature
justifie un traitement en deux étapes, selon la décomposition d'opérateurs
\begin{equation}
\frac{\partial\phi}{\partial t} =
\underbrace{-\frac{\partial(u\phi)}{\partial x}}_{\mathcal{L}_{\mathrm{adv}}(\phi)}
+ \underbrace{k\left(\phisol-\phi\right)}_{\mathcal{L}_{\mathrm{src}}(\phi)} .
\label{eq:splitting}
\end{equation}
Dans la stratégie retenue, l'étape d'advection produit un état intermédiaire \(\phi^*\), puis l'étape de
relaxation est appliquée localement, cellule par cellule, sur cet état. Le couplage avec l'hydraulique
intervient en amont : \(k\) est calculé à partir de la contrainte pariétale issue de la résolution
hydraulique du même pas, et non imposé indépendamment.

\subsection{Intégration analytique du terme source de relaxation}

En gelant \(k_i^n\) pendant le pas de temps, l'équation locale à résoudre dans chaque cellule est
\begin{equation}
\frac{\mathrm{d}\phi}{\mathrm{d}t} = k_i^n\left(\phisol - \phi\right),
\qquad \phi(0) = \phi_i^*,
\label{eq:relax-locale}
\end{equation}
équation linéaire du premier ordre dont la solution exacte est
\begin{equation}
\phi_i^{n+1} = \phisol - \left(\phisol - \phi_i^*\right)\exp\left(-k_i^n\,\dt\right).
\label{eq:source-exacte-2}
\end{equation}
Cette intégration analytique présente trois avantages par rapport à un traitement explicite ou à un
schéma de type Runge--Kutta.

\paragraph{Absence d'erreur de troncature temporelle pour la source.} Quelle que soit la valeur de
\(k_i^n\dt\) — et celle-ci peut être très grande dans les épisodes d'érosion intense — la
relation~\eqref{eq:source-exacte-2} reste la solution exacte de l'équation locale. Les schémas explicites
introduiraient au contraire une erreur \emph{et} une contrainte de stabilité du type \(k\dt<1\) qui
limiterait le pas de temps sans rapport avec la physique du transport.

\paragraph{Préservation stricte de la borne supérieure.} Comme \(\exp(-k\dt)\in\,]0,1]\), la solution est
une combinaison convexe entre \(\phi_i^*\) et \(\phisol\) :
\begin{equation}
\phi_i^{n+1} = \phi_i^* + \left(\phisol-\phi_i^*\right)\left[1 - e^{-k\dt}\right],
\qquad
\min\left(\phi_i^*,\phisol\right) \le \phi_i^{n+1} \le \max\left(\phi_i^*,\phisol\right).
\label{eq:convexite}
\end{equation}
La borne \(\phi\le\phisol\) est donc préservée indépendamment de la valeur de \(k\dt\), sans troncature
corrective. Dans le couplage retenu, où \(\lambda(\phi)\) et \(f_m(\phi)\) deviennent singulières lorsque
\(\phi\to\phisol\), cette propriété n'est pas cosmétique : elle empêche l'apparition d'arguments hors
domaine dans les fermetures hydrauliques, et donc l'interruption du calcul.

\paragraph{Saturation naturelle du terme d'érosion.} La source tend vers zéro lorsque
\(\phi\to\phisol\), de sorte que le modèle possède un mécanisme d'auto-limitation : la vitesse
d'injection des solides décroît à mesure que la suspension se concentre. Ce comportement est celui décrit à
la section~\ref{sec:couplage} et contraste avec une formulation à source directe, qui injecterait
\(\mdot\) indépendamment de l'état de concentration atteint et autoriserait \(\phi > \phisol\).

\subsection{Préservation des bornes et conditions de validité}

La préservation de la borne supérieure par~\eqref{eq:source-exacte-2} repose sur l'hypothèse que
\(\phi_i^*\) est lui-même admissible. Cette hypothèse se décompose en deux conditions, correspondant aux
deux phases du pas de temps.
\begin{enumerate}
  \item \textbf{Admissibilité de l'état initial et de la condition d'entrée.} Le champ initial doit
        satisfaire \(0\le\phi_0(x)\le\phisol\), et la condition d'entrée doit respecter la même borne.
        Pour les simulations présentées, l'entrée est en eau claire (\(\phi_{\mathrm{in}} = 0\)), ce qui
        est le cas le plus contraignant en termes de gradient à la limite amont.
  \item \textbf{Admissibilité de l'état post-advectif.} La reconstruction limitée et le flux de
        Godunov~\eqref{eq:flux-godunov} préservent la positivité sous la condition de
        stabilité~\eqref{eq:cfl-half}, comme établi à la section~\ref{sec:positivite}. Il en résulte que
        \(\phi_i^*\) reste compris entre les valeurs extrêmes des cellules voisines à l'instant
        précédent, donc dans l'intervalle \([0,\phisol]\) par récurrence.
\end{enumerate}
Le résultat est donc le suivant : la borne \([0,\phisol]\) est préservée \emph{par construction} pour tout
pas de temps tel que la condition CFL soit respectée, sans avoir recours à une troncature de sécurité.
Cette propriété mérite d'être explicitement identifiée comme un résultat de la stratégie numérique
retenue, car elle conditionne la robustesse du couplage non linéaire. En pratique, une troncature
résiduelle est néanmoins appliquée et son effet est quantifié par le défaut de
masse~\eqref{eq:defaut-masse}, afin que la vérification de conservation reste mesurable.

\subsection{Couplage entre érosion, concentration et évolution géométrique}

La cohérence temporelle du couplage repose sur l'ordre des opérations. Dans l'algorithme retenu, pour un
pas de temps donné, les opérations sont effectuées dans la séquence suivante :
\begin{enumerate}[label=(\roman*)]
  \item résolution du débit \(Q^n\) à partir de l'état \((R^n,\phi^n)\), par le solveur
        d'encadrement--dichotomie ;
  \item calcul des grandeurs locales (\(u^n\), \(Re^n\), \(f_D^n\), \(f_m^n\), \(\tau_b^n\)) et du flux
        érodé \(\mdot^n\) ;
  \item calcul du pas de temps \(\dt\) à partir de la condition CFL ;
  \item mise à jour de la géométrie, \(R^n \to R^{n+1}\), par~\eqref{eq:dRdt} ;
  \item étape d'advection du transport, \(\phi^n \to \phi^*\) ;
  \item intégration analytique de la source, \(\phi^* \to \phi^{n+1}\).
\end{enumerate}
Deux remarques sont nécessaires pour interpréter correctement cette séquence.

\paragraph{Le coefficient de relaxation est évalué sur l'état hydraulique courant.} Le coefficient \(k^n\)
est calculé à partir de \(\mdot^n\), lui-même issu de \(\tau_b^n\), c'est-à-dire de l'état hydraulique
déterminé \emph{avant} la mise à jour de la géométrie. Le transport et l'érosion utilisent donc le même
état de référence à l'intérieur d'un pas ; il n'y a pas d'incohérence interne au pas, mais un décalage
d'un pas entre l'état hydraulique et l'état géométrique produit.

\paragraph{Le rayon utilisé dans \(k = 2\mdot/(R\rho_s)\) et dans l'advection.} Le choix de l'évaluer
avant ou après mise à jour du rayon constitue un choix de discrétisation de premier ordre en \(\dt\) : les
deux variantes sont consistantes, mais elles doivent être appliquées de manière identique sur toute la
campagne de simulation. Ce point, sans conséquence physique, doit être documenté car il influence la
comparaison fine entre configurations.

\subsection{Conséquences numériques sur la stabilité du modèle}

Le traitement retenu a trois conséquences directes sur la stabilité et la robustesse du modèle couplé.

\paragraph{Découplage de la contrainte de stabilité de la source.} L'intégration analytique supprime la
contrainte classique \(k\dt \ll 1\) qui pèserait sur un traitement explicite. Le pas de temps reste donc
gouverné par la seule condition CFL advective, ce qui est indispensable lorsque l'érosion s'accélère : dans
le régime d'emballement, \(k\) croît fortement et un schéma explicite imposerait un pas de temps
décroissant sans limite, alors que la dynamique du transport reste gouvernée par l'advection.

\paragraph{Préservation de l'admissibilité des fermetures.} La garantie \(\phi\le\phisol\) assure que les
arguments de \(\lambda(\phi)\) et \(f_m(\phi)\) restent dans leur domaine de définition. Sans cette
propriété, une seule valeur légèrement supérieure à \(\phisol\) suffirait à produire une valeur
non finie dans la fermeture, puis à contaminer l'ensemble du champ hydraulique et de l'érosion dès le pas
suivant. La stabilité de la chaîne couplée repose donc en partie sur la structure analytique de
l'intégration de la source.

\paragraph{Contrôle explicite des deux mécanismes de dépassement restants.} Deux voies de dépassement
subsistent et doivent être surveillées : (i) la reconstruction spatiale, si la condition CFL n'est pas
respectée, comme établi au tableau~\ref{tab:seuil-cfl} ; (ii) l'erreur de découplage entre l'hydraulique
et la géométrie, contrôlée par le rapport \(\tau_h/\ter\) et par le diagnostic CFL présenté ci-après.
Ces deux contrôles constituent, avec le défaut de masse~\eqref{eq:defaut-masse}, les trois indicateurs de
qualité recommandés pour chaque simulation.

%=====================================================================================
\section{Stratégie temporelle, stabilité et vérification numérique}
\label{sec:verification}
%=====================================================================================

\subsection{Avancement temporel explicite de la partie advective}

Le pas de temps est explicitement borné par la propagation de la concentration. La vitesse utilisée pour
ce contrôle est celle du transport des particules, \(a_i = \beta_u\,u_i\), et non la vitesse du fluide
seule, ce qui autorise à représenter un éventuel glissement lorsque \(\beta_u\neq1\). Le pas admissible
pour l'étape d'advection est donc
\begin{equation}
\dt_{\mathrm{adv}} = \mathrm{CFL}_{\mathrm{tgt}}\,
\frac{\Delta x}{\max_i\left|\beta_u\,u_i\right|}.
\label{eq:dt-adv}
\end{equation}
La contrainte est évaluée à chaque pas, car la vitesse maximale évolue fortement au cours de la
simulation : elle augmente avec l'élargissement local du conduit et diminue lorsque la résistance du
mélange s'accroît.

\subsection{Condition de Courant et adaptation du pas de temps}

Deux situations doivent être distinguées avec soin, car elles ont des conséquences opposées sur la
signification d'une étude de convergence.

\paragraph{Écoulement établi à vitesse bornée.} Lorsque la vitesse maximale varie lentement, le pas de
temps calculé par~\eqref{eq:dt-adv} est proche de sa borne haute \(\dt_{\max}\) et le nombre de Courant
effectif, \(\mathrm{CFL}_{\mathrm{eff}} = \max_i|a_i|\,\dt/\Delta x\), est proche de
\(\mathrm{CFL}_{\mathrm{tgt}}\). La résolution temporelle suit alors fidèlement le réglage demandé.

\paragraph{Régime d'emballement.} Lorsque le rayon croît rapidement en sortie, la vitesse locale augmente
comme \(R^{-2}\) à débit croissant, de sorte que \(\dt_{\mathrm{adv}}\) décroît fortement. Si la borne
\(\dt_{\max}\) prend le relais, le nombre de Courant effectif chute très en dessous de la valeur visée et
la résolution temporelle n'est plus contrôlée par le réglage CFL. C'est précisément ce phénomène que
le diagnostic d'effectivité doit détecter.

\subsection{Diagnostic d'effectivité de la condition CFL}
\label{sec:cfl-diagnostic}

Le diagnostic retenu est fondé sur la mesure de la \emph{fraction de pas bornés par la borne haute} du pas
de temps. On définit
\begin{equation}
\mathcal{C}_{\max} = \frac{1}{N_{\mathrm{steps}}}
\sum_{n} \mathbb{1}\left[\dt_{\mathrm{adv}}^n > \dt_{\max}\right],
\label{eq:clipping}
\end{equation}
c'est-à-dire la proportion de pas de temps pour lesquels la contrainte CFL aurait conduit à un pas
supérieur à la borne imposée. Lorsque \(\mathcal{C}_{\max}\to 1\), la totalité des pas est bornée par
\(\dt_{\max}\) : le pas de temps n'est plus déterminé par la physique du transport, mais par un réglage
numérique. Dans cette situation, modifier le nombre de Courant nominal ne change pas le pas effectivement
utilisé, et toute étude de convergence temporelle devient non discriminante.

Ce diagnostic doit être complété par trois indicateurs quantitatifs qui en font un véritable outil de
vérification, et non un simple voyant :
\begin{itemize}
  \item le nombre de Courant effectif, en moyenne \emph{et} en maximum, ce dernier devant rester inférieur
        au seuil de stabilité~\eqref{eq:cfl-half} ;
  \item la fraction de pas bornés par la borne basse, qui signale un éventuel sur-raffinement inutile ou
        une instabilité naissante ;
  \item le défaut de masse~\eqref{eq:defaut-masse} cumulé, qui mesure l'effet des troncatures appliquées au
        champ de concentration.
\end{itemize}
L'ensemble de ces indicateurs constitue la procédure de vérification recommandée avant toute
interprétation d'une étude de convergence.

\subsection{Vérification de l'influence de \(\dt_{\max}\)}
\label{sec:dtmax}

La borne \(\dt_{\max}\) est un paramètre numérique dont l'influence doit être quantifiée séparément, et non
absorbée dans un réglage global. La procédure recommandée consiste à :
\begin{enumerate}[label=\arabic*.]
  \item exécuter une paire de simulations ne différant que par \(\dt_{\max}\), toutes choses égales par
        ailleurs, et comparer les grandeurs intégrales (temps de doublement, rapport d'agrandissement
        final, position du maximum de concentration) ;
  \item vérifier que la fraction \(\mathcal{C}_{\max}\) de l'équation~\eqref{eq:clipping} est soit proche de
        \(0\) (régime réellement contrôlé par la CFL), soit proche de \(1\) (régime borné par
        \(\dt_{\max}\), dans lequel une étude CFL est sans objet) ;
  \item ne présenter une analyse de convergence en \(\mathrm{CFL}\) que dans le premier cas, ou après
        avoir relevé \(\dt_{\max}\) de manière à rendre la contrainte CFL active.
\end{enumerate}
Cette procédure est directement transposable au protocole expérimental numérique décrit dans la section
suivante : le point essentiel est qu'une campagne \((N_x \times \mathrm{CFL})\) n'a de sens que si elle est
précédée du contrôle~\eqref{eq:clipping}.

\subsection{Étude de convergence spatiale et temporelle}

L'étude de convergence retenue combine le raffinement spatial (\(N_x\)) et le réglage du nombre de Courant
cible, sous réserve que la condition d'effectivité soit satisfaite. Les grandeurs comparées sont le
rayon \(R(x)\), la pression \(p(x)\) et la concentration \(\phi(x)\) à des instants de référence exprimés en
multiples de \(\ter\), l'état le plus raffiné servant de référence. Les erreurs sont évaluées après
interpolation sur la grille de référence, en normes \(L_1\), \(L_2\) et \(L_\infty\).

Les champs \(p\) et \(R\) sont des champs réguliers : leurs erreurs décroissent de manière monotone et à un
taux proche de l'ordre formel du schéma. Le champ \(\phi\), en revanche, présente un front dont la
description numérique obéit à une analyse différente, développée au paragraphe~\ref{sec:fronts-convergence}.

\subsection{Analyse des erreurs en normes \(L_1\), \(L_2\) et \(L_\infty\)}

Le comportement des normes pour une discontinuité capturée par un schéma à limiteur constitue un point
théorique essentiel, qui éclaire directement les résultats numériques observés sur le champ \(\phi\). Un
test de référence permet de l'illustrer quantitativement : une marche de concentration est advectée à
\(\mathrm{CFL} = 0.5\) pendant un nombre entier de traversées de cellules, de sorte que la position exacte
du front soit connue sans ambiguïté ; les erreurs sont alors calculées contre la solution exacte
translatée. Le tableau~\ref{tab:convergence-front} présente l'évolution des normes avec le raffinement.

\begin{table}[H]
\centering
\caption{Comportement des normes avec le raffinement lors de l'advection d'une discontinuité
(\emph{marche} de concentration, \(\mathrm{CFL} = 0.5\), déplacement correspondant à \(N_x/4\) cellules,
erreurs calculées contre la solution exacte translatée). Les produits \(\varepsilon_{L_1}N_x\) et
\(\varepsilon_{L_2}\sqrt{N_x}\) permettent de lire directement l'ordre de convergence.}
\label{tab:convergence-front}
\small
\begin{tabular}{cccccc}
\toprule
\(N_x\) & \(\varepsilon_{L_1}\) & \(\varepsilon_{L_1}\,N_x\) & \(\varepsilon_{L_2}\)
& \(\varepsilon_{L_2}\sqrt{N_x}\) & \(\varepsilon_{L_\infty}\) \\
\midrule
\multicolumn{6}{l}{\emph{Limiteur MC (schéma retenu)}}\\
\(64\)  & \(7.81\times10^{-3}\) & \(0.500\) & \(3.13\times10^{-2}\) & \(0.250\) & \(0.1250\) \\
\(128\) & \(3.91\times10^{-3}\) & \(0.500\) & \(2.21\times10^{-2}\) & \(0.250\) & \(0.1250\) \\
\(256\) & \(1.95\times10^{-3}\) & \(0.500\) & \(1.56\times10^{-2}\) & \(0.250\) & \(0.1250\) \\
\(512\) & \(9.77\times10^{-4}\) & \(0.500\) & \(1.11\times10^{-2}\) & \(0.250\) & \(0.1250\) \\
\(1024\) & \(4.88\times10^{-4}\) & \(0.500\) & \(7.81\times10^{-3}\) & \(0.250\) & \(0.1250\) \\
\(2048\) & \(2.44\times10^{-4}\) & \(0.500\) & \(5.52\times10^{-3}\) & \(0.250\) & \(0.1250\) \\
\midrule
\multicolumn{6}{l}{\emph{Schéma amont du premier ordre (comparaison)}}\\
\(64\)  & \(7.00\times10^{-2}\) & \(4.48\) & \(1.43\times10^{-1}\) & \(1.140\) & \(0.4300\) \\
\(128\) & \(4.97\times10^{-2}\) & \(6.36\) & \(1.20\times10^{-1}\) & \(1.361\) & \(0.4503\) \\
\(256\) & \(3.52\times10^{-2}\) & \(9.01\) & \(1.01\times10^{-1}\) & \(1.623\) & \(0.4648\) \\
\(512\) & \(2.49\times10^{-2}\) & \(12.75\) & \(8.54\times10^{-2}\) & \(1.932\) & \(0.4751\) \\
\(1024\) & \(1.76\times10^{-2}\) & \(18.05\) & \(7.18\times10^{-2}\) & \(2.298\) & \(0.4824\) \\
\(2048\) & \(1.25\times10^{-2}\) & \(25.53\) & \(6.04\times10^{-2}\) & \(2.734\) & \(0.4875\) \\
\bottomrule
\end{tabular}
\begin{tablenotes}
\small
\item Pour le schéma limité, les deux produits \(\varepsilon_{L_1}N_x\) et
\(\varepsilon_{L_2}\sqrt{N_x}\), ainsi que \(\varepsilon_{L_\infty}\), sont rigoureusement constants :
l'erreur \(L_1\) décroît comme \(\Delta x\), l'erreur \(L_2\) comme \(\Delta x^{1/2}\), et l'erreur
\(L_\infty\) \emph{ne décroît pas}. Pour le schéma amont, l'erreur \(L_\infty\) tend vers la moitié de
l'amplitude du saut et l'erreur \(L_1\) décroît beaucoup plus lentement.
\end{tablenotes}
\end{table}

Le résultat est sans ambiguïté, et il est plus riche que la seule décroissance de \(L_1\). Pour une
discontinuité capturée par le schéma limité :
\begin{itemize}
  \item l'erreur \(L_1\) décroît proportionnellement à \(\Delta x\), soit un ordre de convergence égal à
        l'unité, valeur attendue pour un front non résolu ;
  \item l'erreur \(L_2\) décroît comme \(\Delta x^{1/2}\), c'est-à-dire beaucoup plus lentement, ce qui est
        la loi asymptotique pour une discontinuité (et non l'ordre deux du schéma en zone régulière) ;
  \item l'erreur \(L_\infty\) tend vers une \emph{constante} strictement indépendante de la résolution,
        égale ici à \(1/8\) de l'amplitude du saut.
\end{itemize}
Le raffinement ne peut donc pas réduire l'erreur maximale locale d'un front capturé : il peut seulement
réduire l'étendue spatiale de la région où cette erreur est concentrée. Ce point explique pourquoi une
étude de convergence fondée sur la seule norme \(L_\infty\) conduirait à une conclusion erronée de
non-convergence, alors que le schéma converge effectivement. À titre de comparaison, le schéma amont
présente une erreur \(L_1\) qui ne décroît que très lentement (le produit \(\varepsilon_{L_1}N_x\) croît
avec le raffinement) et une erreur \(L_\infty\) qui plafonne à la moitié du saut : le bénéfice de la
reconstruction limitée n'est donc pas seulement une réduction d'un facteur constant, mais un
\emph{changement d'ordre de convergence}.

\subsection{Particularités de convergence des fronts advectifs de concentration}
\label{sec:fronts-convergence}

Dans les simulations couplées, l'erreur \(L_\infty\) mesurée sur \(\phi\) combine deux contributions de
nature différente, dont les lois de décroissance ne sont pas les mêmes :
\begin{itemize}
  \item une contribution \emph{intrinsèque}, liée à la forme du profil capturé ; comme établi au
        tableau~\ref{tab:convergence-front}, elle ne décroît pas avec le raffinement et reste de l'ordre
        d'une fraction de l'amplitude du saut ;
  \item une contribution de \emph{position}, liée au décalage du front entre deux maillages. Elle décroît
        proportionnellement à \(\Delta x\), puisqu'un décalage d'une ou deux cellules produit, sur un front
        raide, une différence locale du même ordre que le saut lui-même.
\end{itemize}
Selon le maillage et l'instant considérés, l'une ou l'autre contribution domine. Cette décomposition
explique les deux observations apparemment contradictoires que l'on rencontre dans les études de
convergence de modèles à front raide : d'une part, la décroissance régulière des erreurs \(L_1\)/\(L_2\)
avec le raffinement, qui traduit la convergence au sens intégral ; d'autre part, la persistance d'une
erreur \(L_\infty\) élevée sur le champ de concentration, alors même que les champs lisses
(\(p\), \(R\)) convergent rapidement et régulièrement.

\begin{figure}[H]
\centering
\includegraphics[width=\linewidth]{figure4.png}
\caption{Sensibilité au maillage et convergence spatiale des champs couplés. Les profils de pression et de
rayon se superposent rapidement lorsque le maillage est raffiné ; l'écart sur la concentration en
suspension reste localisé au niveau du front advectif, ce qui explique la persistance de la sensibilité en
norme \(L_\infty\).}
\label{fig:convergence-maillage}
\end{figure}

La conclusion méthodologique est donc la suivante : pour un modèle couplé comportant un front advectif
raide, la convergence doit être interprétée norme par norme. La convergence \(L_1\)/\(L_2\) atteste la
convergence au sens intégral et donc la validité des grandeurs globales extraites du calcul (temps
caractéristiques, rapports d'agrandissement, indicateurs de régime). L'erreur \(L_\infty\) sur \(\phi\) ne
constitue pas un critère de convergence adapté, et son ampleur doit être expliquée par la nature de la
discontinuité plutôt que corrigée par un raffinement supplémentaire.

\subsection{Vérification de la loi de résistance du mélange et validation croisée}
\label{sec:verif-fm}

La fermeture~\eqref{eq:fm} et son plafond~\eqref{eq:fmmax} étant appliqués au sein du solveur, leur
traitement doit être vérifié sur deux plans : la non-régression des résultats de référence, et l'aptitude
du modèle complet à reproduire les temps publiés pour des configurations indépendantes.

\paragraph{Non-régression de la campagne A.} Trois traitements du plafond ont été comparés : la loi sans
plafond, le minimum abrupt~\eqref{eq:fmmax} et une variante régularisée à raccord \(\mathcal{C}^1\). Pour la
configuration de référence de la campagne A, les trois donnent \emph{exactement} le même temps de
doublement, \(t_2 = 668.5\) s : le facteur de résistance du mélange y reste compris entre \(1.0000\) et
\(1.0001\), c'est-à-dire que la rétroaction rhéologique est inactive à l'échelle du laboratoire et pour un
matériau fin. Les résultats publiés de la campagne A sont donc insensibles au mode de traitement du
plafond, ce qui autorise à retenir le minimum abrupt, conforme à la formulation publiée, sans réserve de
non-régression.

\paragraph{Effet de l'absence de plafond.} Le tableau~\ref{tab:fm-diagnostic} compare les valeurs
maximales du facteur de résistance atteintes au cours des calculs.

\begin{table}[H]
\centering
\caption{Diagnostic du facteur de résistance du mélange. La colonne « sans plafond » correspond à la
loi~\eqref{eq:fm} appliquée sans le minimum~\eqref{eq:fmmax} ; la comparaison met en évidence l'artefact
numérique que le plafond a pour fonction d'éviter.}
\label{tab:fm-diagnostic}
\small
\begin{tabular}{lcccc}
\toprule
Campagne & \(f_{m,\max}\) & \(\max f_m\) sans plafond & \(\max f_m\) avec plafond & symptôme sans plafond \\
\midrule
A (lab., \(d_p = 20\)--\(75\ \si{\micro m}\)) & \(5\) & \(1.0001\) & \(1.0001\)
& aucun (\(f_m \approx 1\)) \\
B (terrain, \(d_p = 1\) mm) & \(2000\) & \(1.46\times10^{12}\) & \(2000\)
& pic local de rayon, arrêt sur \(R_{\mathrm{break}}\) \\
\bottomrule
\end{tabular}
\begin{tablenotes}
\small
\item Pour la campagne B sans plafond, le rayon maximal atteint \(179\,R_0\) alors que le rayon de sortie
n'est que de \(3.5\,R_0\) : l'érosion se concentre en un pic isolé, la comparaison avec une évolution
d'ensemble perd son sens, et le calcul s'interrompt sur le critère de sécurité après une durée sans
rapport avec la dynamique recherchée.
\end{tablenotes}
\end{table}

\paragraph{Validation croisée sur la campagne B.} Une fois le plafond appliqué dans le solveur, la
campagne B peut être confrontée aux valeurs publiées. Le critère retenu est l'atteinte de
\(R_{\max} = R_{\mathrm{break}} = 0.5\) m, soit \(167\,R_0\) ; il correspond, dans le code, à la
condition d'arrêt effective. Le tableau~\ref{tab:validation-campagneB} présente la comparaison pour les
quatre configurations (deux charges hydrauliques, avec et sans résistance aval).

\begin{table}[H]
\centering
\caption{Validation croisée de la campagne B : temps d'échec (atteinte de \(R_{\max} = 0.5\) m) comparés
aux valeurs publiées. La dernière colonne donne l'écart relatif du modèle aux valeurs publiées.}
\label{tab:validation-campagneB}
\small
\begin{tabular}{ccccc}
\toprule
\(\Delta p\) (MPa) & \(\kout\) & \(\ter\) (s) & \(t_{\mathrm{fail}}\) modèle (s) & \(t_{\mathrm{fail}}\) publié (s) \\
\midrule
\(1.50\) & \(0\)   & \(24.7\) & \(128.60\) & \(128.1\) \quad (\(+0.4\%\)) \\
\(1.50\) & \(500\) & \(24.7\) & \(162.13\) & \(161.8\) \quad (\(+0.2\%\)) \\
\(3.33\) & \(0\)   & \(11.1\) & \(55.26\)  & \(55.0\) \quad (\(+0.5\%\)) \\
\(3.33\) & \(500\) & \(11.1\) & \(63.22\)  & \(62.9\) \quad (\(+0.5\%\)) \\
\bottomrule
\end{tabular}
\begin{tablenotes}
\small
\item Les rapports de retard induits par la résistance aval sont reproduits à mieux de \(0.2\%\) :
\(\times1.261\) contre \(\times1.263\) à \(1.5\) MPa, et \(\times1.144\) contre \(\times1.144\) à
\(3.33\) MPa. Le seuil \(f_m = f_{m,\max}\) est atteint dans les quatre configurations, de sorte que
l'accord porte bien sur le régime plafonné.
\end{tablenotes}
\end{table}

L'accord obtenu, meilleur que \(0.5\%\) sur quatre configurations et deux grandeurs indépendantes, est le
résultat de validation le plus solide de ce chapitre : il porte simultanément sur la fermeture
hydraulique~\eqref{eq:dw}, sur la loi d'érosion~\eqref{eq:erosion}, sur le traitement de la perte
singulière aval~\eqref{eq:charge-aval} et sur le plafond~\eqref{eq:fmmax}. Il confirme également que le
temps d'échec est, dans ce régime, fixé par la seule balance de charge et non par la résolution spatiale :
les valeurs du tableau~\ref{tab:validation-campagneB} sont stables lorsque \(N_x\) passe de \(100\) à
\(600\) cellules.

\paragraph{Cohérence de la campagne A.} Pour mémoire, la même procédure appliquée à la campagne A
reproduit les deux valeurs publiées : \(668\) s contre \(666\) s publiés pour la configuration libre
(\(+0.4\%\)) et \(877.8\) s contre \(872\) s pour la configuration avec résistance aval
(\(+0.7\%\)), cette dernière correspondant à un coefficient de perte modéré
\(\kout \approx 0.15\) (section~\ref{sec:kout}).

\subsection{Critères de vérification synthétiques}

En rassemblant les éléments des sections précédentes, la procédure de vérification recommandée pour chaque
simulation comporte sept contrôles :
\begin{enumerate}[label=\arabic*.]
  \item \textbf{Contrôle analytique de l'état initial} : la contrainte pariétale initiale doit vérifier
        \(\tau_b = R_0\Delta p/(2L)\) et le rapport \(\tau_b/\tau_c\) doit être cohérent avec le régime
        attendu (section~\ref{sec:instabilite}) ;
  \item \textbf{Contrôle analytique du taux de croissance} : tant que la géométrie reste quasi uniforme,
        le rayon doit suivre la loi~\eqref{eq:R-exponentielle} avec l'écart attendu lié à la non-uniformité
        et aux effets d'entrée ;
  \item \textbf{Contrôle d'effectivité CFL} : la fraction de pas bornés par \(\dt_{\max}\)
        (\eqref{eq:clipping}) et le nombre de Courant effectif maximal, qui doit rester inférieur au seuil
        de stabilité~\eqref{eq:cfl-half} ;
  \item \textbf{Contrôle de conservation} : le défaut de masse cumulé~\eqref{eq:defaut-masse} doit rester
        négligeable devant la masse transitée ;
  \item \textbf{Contrôle de convergence} : décroissance des erreurs \(L_1\)/\(L_2\) sur les champs
        \(p\), \(R\) et \(\phi\), avec l'interprétation norme par norme établie
        au paragraphe~\ref{sec:fronts-convergence} ;
  \item \textbf{Contrôle de la fermeture de résistance} : la valeur maximale de \(f_m\) atteinte au cours
        du calcul doit être comparée au plafond \(f_{m,\max}\) retenu, et le résultat interprété
        différemment selon que le plafond est atteint ou non (paragraphe~\ref{par:fm-plafond} et
        section~\ref{sec:verif-fm}) ;
  \item \textbf{Contrôle du bilan de charge} : la part de \(\Delta p\) consommée par chaque terme
        (frottement réparti, perte singulière d'entrée, perte singulière aval) doit être extraite des
        sorties, car c'est elle qui détermine si une résistance singulière peut retarder ou arrêter le
        processus (tableau~\ref{tab:bilan-charge}).
\end{enumerate}
Ces sept contrôles constituent la base de la vérification présentée au chapitre suivant. Ils présentent
l'avantage d'être tous calculables à partir des sorties standard du code et de ne requérir aucun
développement supplémentaire, à l'exception du dernier, qui suppose de distinguer explicitement les
contributions de charge dans le résidu hydraulique.


%=====================================================================================
\section{Applications numériques et analyse paramétrique}
\label{sec:resultats}
%=====================================================================================

\subsection{Cas de référence et régime d'instabilité}

Le cas de référence correspond à la campagne A (tableau~\ref{tab:param-campagnes}), c'est-à-dire un conduit
de laboratoire de longueur \(L = 0.117\) m et de rayon initial \(R_0 = 3\) mm, soumis à une perte de charge
imposée équivalente à \(0.5\) m de colonne d'eau, en l'absence de résistance aval (\(\kout = 0\)). L'état
hydraulique initial prédit par le modèle est rassemblé dans le tableau~\ref{tab:etat-initial} et comparé aux
valeurs analytiques~\eqref{eq:tau-analytique}.

\begin{table}[H]
\centering
\caption{État hydraulique initial du cas de référence (campagne A, \(\kout=0\)) : comparaison entre les
valeurs analytiques exactes et le calcul.}
\label{tab:etat-initial}
\small
\begin{tabular}{llcc}
\toprule
Grandeur & Relation & Analytique & Calculé \\
\midrule
Contrainte pariétale & \(\tau_b = R_0\Delta p/(2L)\) & \(62.9\) Pa & \(62.9\) Pa \\
Rapport de supercriticité & \(\tau_b/\tau_c\) & \(62.9\) & --- \\
Vitesse moyenne & \(u_0 = Q_0/(\pi R_0^2)\) & --- & \(4.5\) m/s \\
Débit & \(Q_0\) & --- & \(1.28\times10^{-4}\ \mathrm{m^3/s}\) \\
Nombre de Reynolds & \(Re = 2\rho u_0R_0/\mu_w\) & --- & \(2.7\times10^{4}\) \\
Source volumique & \(S_A = 2\pi R_0\mdot/\rho_s\) & --- & \(5.9\times10^{-8}\ \mathrm{m^2/s}\) \\
Rayon critique & \(R_c = 2L\tau_c/\Delta p\) & \(0.048\) mm & --- \\
\bottomrule
\end{tabular}
\end{table}

Le rapport de supercriticité \(\tau_b/\tau_c \approx 63\) est très supérieur à l'unité : l'érosion est
immédiatement active, sans phase de latence. La conséquence dynamique a été établie analytiquement au
paragraphe~\ref{sec:instabilite} : en l'absence de résistance aval, le rayon suit une loi exponentielle
de taux \(1/\ter\) et il n'existe aucun point fixe géométrique stable. Le calcul couplé reproduit ce
comportement, avec une croissance monotone du rayon jusqu'à l'atteinte du critère d'arrêt.

Trois observations doivent accompagner ce résultat, conformément aux critères d'arrêt
distingués au paragraphe~\ref{sec:criteres-arret}. Premièrement, l'arrêt ne résulte pas d'une divergence
numérique : le solveur hydraulique conserve son encadrement à chaque pas et le résidu en pression reste
inférieur à la tolérance. Deuxièmement, l'horizon de validité quantitative
(\(R/R_0 < 1.95\) pour \(R/L<0.05\)) est franchi bien avant la fin du calcul, de sorte que les rapports
d'agrandissement atteints constituent des indicateurs de régime et non des prédictions. Troisièmement, le
temps caractéristique d'arrêt, exprimé en multiples de \(\ter\), est essentiellement fixé par la définition
du critère de sécurité \(R_{\mathrm{break}}\) ; il ne doit donc pas être présenté comme une grandeur
prédictive sans préciser cette dépendance.

\subsection{Influence de l'érodibilité du matériau}

Deux sols idéalisés ont été considérés afin d'explorer la gamme d'érodibilité des mélanges silto-argileux
compactés : un sol érodible (A, \(k_{er} = 10^{-4}\ \mathrm{s/m}\), \(\tau_c = 1\) Pa) et un sol résistant
(B, \(k_{er} = 10^{-5}\ \mathrm{s/m}\), \(\tau_c = 10\) Pa). Le tableau~\ref{tab:erodibilite} présente les
temps caractéristiques et les temps de doublement.

\begin{table}[H]
\centering
\caption{Influence de l'érodibilité sur les temps caractéristiques (campagne A, \(\kout = 0\)). Les
valeurs analytiques sont issues de~\eqref{eq:ter} et de \(t_2 = \ter\ln 2\).}
\label{tab:erodibilite}
\small
\begin{tabular}{lccccc}
\toprule
Sol & \(k_{er}\) (s/m) & \(\tau_c\) (Pa) & \(\ter\) (s) & \(t_2\) analytique (s) & \(t_2\) calculé (s) \\
\midrule
Érodible (A) & \(1.0\times10^{-4}\) & \(1.0\) & \(949.4\) & \(658\) & \(666\) \\
Résistant (B) & \(1.0\times10^{-5}\) & \(10.0\) & \(9493.6\) & \(6579\) & \(6976\) \\
\bottomrule
\end{tabular}
\begin{tablenotes}
\small
\item Le passage de A à B divise le coefficient d'érosion par dix et multiplie la contrainte critique par
dix. Le rapport des temps de doublement calculés, \(10.5\), est proche du rapport
\(\ter(\mathrm{B})/\ter(\mathrm{A}) = 10\) prédit par~\eqref{eq:ter}, l'écart de quelques pourcents
provenant de la dépendance non triviale au terme \(R_c\) et du caractère non uniforme du profil.
\end{tablenotes}
\end{table}

L'accord entre les temps analytiques et les temps calculés est de \(1\%\) pour le sol A et de \(6\%\) pour
le sol B. Ce résultat valide simultanément la définition du temps caractéristique~\eqref{eq:ter} et la
capacité du modèle à restituer l'effet de premier ordre de l'érodibilité sur les échelles de temps. Il
convient toutefois de rappeler que ces échelles sont d'abord des durées d'émoussement géométrique
(\(t_2\)) et non des temps de rupture : elles se rattachent à la phase initiale de l'érosion, celle qui
reste compatible avec l'hypothèse d'un conduit quasi uniforme et avec le domaine de validité
\eqref{eq:validite-1D}.

\subsection{Influence de la résistance hydraulique aval}
\label{sec:kout}

L'effet d'une résistance aval est gouverné par l'équation de charge
\begin{equation}
\Delta p = \underbrace{\frac{f_D\,\rho\,u^2\,L}{4R}}_{\text{frottement réparti}}
+ \underbrace{\kout\,\frac{1}{2}\rho u^2}_{\text{perte singulière aval}},
\label{eq:charge-aval}
\end{equation}
dans laquelle, \emph{à vitesse donnée}, la perte répartie décroît comme \(R^{-5/4}\) (loi de Blasius)
lorsque le conduit s'élargit, tandis que la perte singulière reste proportionnelle à \(u^2\) et
indépendante de la géométrie. Il en résulte une
propriété structurelle importante : lorsque \(\kout\) est grand, la vitesse de sortie tend vers une valeur
asymptotique indépendante de la géométrie,
\begin{equation}
u \longrightarrow u_{\max} = \sqrt{\frac{2\Delta p}{\kout\,\rho}},
\label{eq:umax}
\end{equation}
qui borne la contrainte pariétale accessible et transforme la croissance exponentielle en une croissance
quasi linéaire. La résistance aval ne module donc pas seulement la vitesse du processus : au-delà d'un
certain seuil, elle \emph{change la nature de la dynamique} et peut conduire à un arrêt de l'érosion
lorsque la contrainte pariétale devient inférieure à la contrainte critique.

Le tableau~\ref{tab:kout} présente la sensibilité du temps de doublement à \(\kout\) pour la configuration
de la campagne A. Les colonnes \(\mathrm{ODE}\) et \emph{simulée} distinguent deux niveaux de modèle :
l'intégration de l'équation d'évolution avec la relation de charge~\eqref{eq:charge-aval} à géométrie
figée, et la simulation couplée complète, où le temps de doublement est déduit du taux de croissance
mesuré au début du calcul. La valeur \(\kout = 0\) reproduit le temps de doublement mesuré (\(666\) s) à
\(0.4\%\) près, ce qui constitue une vérification préalable à l'interprétation des autres lignes.

\begin{table}[H]
\centering
\caption{Sensibilité de la dynamique d'érosion à la résistance aval (campagne A). Les grandeurs sont
évaluées à la géométrie initiale ; \(t_2\) est le temps de doublement calculé par intégration de
l'équation d'évolution avec la relation~\eqref{eq:charge-aval}.}
\label{tab:kout}
\small
\begin{tabular}{ccccccc}
\toprule
\(\kout\) & \(\Pi_{\mathrm{aval}}\) & \(u_0\) (m/s) & \(\tau_{b,0}\) (Pa) & \(\Delta p_{\mathrm{aval}}/\Delta p\)
& \(t_2\) ODE (s) & \(t_2\) simulée (s) \\
\midrule
\(0\)    & \(0\)      & \(4.52\) & \(62.9\) & \(0\)      & \(666\) & \(668\) \\
\(0.1\)  & \(0.21\)   & \(4.07\) & \(52.3\) & \(0.17\)   & \(895\) & \(810\) \\
\(0.5\)  & \(1.04\)   & \(3.09\) & \(32.3\) & \(0.49\)   & \(1725\) & \(1330\) \\
\(1\)    & \(2.08\)   & \(2.51\) & \(22.5\) & \(0.64\)   & \(2689\) & \(1937\) \\
\(10\)   & \(20.8\)   & \(0.96\) & \(4.2\)  & \(0.96\)   & \(2.09\times10^{4}\) & \(1.27\times10^{4}\) \\
\(500\)  & \(1.04\times10^{3}\) & \(0.14\) & \(0.14\) & \(1.00\) & arrêt immédiat & arrêt immédiat \\
\bottomrule
\end{tabular}
\begin{tablenotes}
\small
\item Pour \(\kout = 500\), la contrainte pariétale initiale (\(0.14\) Pa) est inférieure à la contrainte
critique (\(1\) Pa) : la loi~\eqref{eq:erosion} prédit alors une absence totale d'érosion dès l'état
initial. La ligne \(\kout = 10\) correspond à un régime où la perte singulière capte \(96\%\) de la charge
motrice.
\item La colonne simulée est obtenue par une simulation couplée de \(20\) s (\(N_x = 40\) à
\(1200\)), le taux de croissance \(\dot R/R\) étant mesuré sur la fenêtre puis extrapolé en
\(t_2 = \ln 2/\dot R\,R^{-1}\). L'écart systématique avec la colonne ODE, croissant avec \(\kout\),
provient de l'accroissement du débit avec le rayon (\(Q \propto R^{19/7}\), voir section~\ref{sec:instabilite}),
que l'intégration à \(\Delta p\) constant néglige.
\end{tablenotes}
\end{table}

Trois enseignements se dégagent de ce tableau. Premièrement, l'effet d'une résistance aval est fortement
non linéaire : un coefficient de \(0.1\) n'allonge le temps de doublement que de \(20\%\), tandis qu'un
coefficient de \(10\) le multiplie par plus de dix-neuf. Deuxièmement, l'effet n'est pas seulement
quantitatif : à partir d'un seuil, la croissance exponentielle laisse place à une croissance linéaire
bornée, puis à un arrêt complet. Troisièmement, le rôle du \emph{rapport} entre perte répartie et perte
singulière est déterminant, et non la valeur brute de \(\kout\) : pour un conduit plus long ou un rayon
plus petit, une même valeur de \(\kout\) serait négligeable. La comparaison des configurations doit donc
être conduite à rapport de pertes comparable, ce qui plaide pour l'introduction d'un paramètre sans
dimension
\begin{equation}
\Pi_{\mathrm{aval}} = \frac{\kout\,\rho/2}{f_D\,\rho\,L/(4R)}
= \frac{2\kout R}{f_D L},
\label{eq:pi-aval}
\end{equation}
qui mesure la part de la charge consommée par la singularité aval. Pour la configuration de référence, ce paramètre vaut \(0.21\) pour
\(\kout = 0.1\), \(1.04\) pour \(\kout = 0.5\), \(2.1\) pour \(\kout = 1\) et \(20.8\) pour
\(\kout = 10\) : il identifie immédiatement le régime atteint, la fraction de charge consommée par la
singularité aval s'écrivant simplement \(\Pi_{\mathrm{aval}}/(1+\Pi_{\mathrm{aval}})\).

\paragraph{Résolution de la comparaison avec les valeurs publiées.}
Le tableau~\ref{tab:kout} appelle une mise au point, car il porte le message d'ingénierie du chapitre. La
valeur publiée pour la configuration avec résistance aval (\(t_2 = 872\) s) correspond, dans le modèle, à
un coefficient de perte aval \emph{modéré}, \(\kout \approx 0.15\), et non à \(\kout = 10\) : la
configuration correspondante donne \(t_2 = 877.8\) s, soit un écart de \(0.7\%\) seulement, du même ordre
que celui obtenu pour la configuration libre (\(668\) s contre \(666\) s publiés). L'accord est donc
retrouvé dès lors que les deux campagnes ne sont plus confondues : les coefficients élevés
(\(10\), \(500\)) présents dans les fichiers de scénarios du code décrivent la \emph{campagne B}, à
l'échelle du terrain, et non la campagne A. Le rapport de pertes~\eqref{eq:pi-aval} permet de s'en
convaincre sans calcul : à l'échelle du laboratoire, \(\kout = 10\) conduit à
\(\Pi_{\mathrm{aval}} \approx 21\), régime où la singularité aval capte \(96\%\) de la charge et où le temps
de doublement dépasse dix-neuf fois sa valeur de référence. Une telle configuration ne décrit plus un essai
au trou, mais un conduit presque obturé ; elle n'a pas de contrepartie expérimentale dans la campagne A.

\subsection{Sensibilité au gradient hydraulique}

La campagne B permet d'explorer l'effet du chargement hydraulique à l'échelle du terrain
(\(L = 100\) m, \(\Delta p\) de \(1.5\) à \(10.0\) MPa). La loi~\eqref{eq:ter} prévoit une décroissance du
temps caractéristique comme \(1/\Delta p\) : multiplier la charge par un facteur \(k\) doit diviser les
temps par le même facteur. Le tableau~\ref{tab:gradient} confronte cette prédiction aux temps
d'échec calculés pour \(\kout = 0\), en prenant pour critère l'atteinte d'un rayon de \(0.5\) m.

\begin{table}[H]
\centering
\caption{Influence du gradient hydraulique sur le temps d'échec (campagne B, \(\kout = 0\)). Le temps
d'échec est défini par l'atteinte de \(R = 0.5\) m. Les valeurs analytiques sont obtenues par
\(t_{\mathrm{fail}} = \ter\ln\left[(R_{\mathrm{break}}-R_c)/(R_0-R_c)\right]\).}
\label{tab:gradient}
\small
\begin{tabular}{cccccc}
\toprule
\(\Delta p\) (MPa) & Gradient \(\Delta p/L\) & \(\ter\) (s) & \(t_{\mathrm{fail}}\) analytique (s)
& \(t_{\mathrm{fail}}\) intégré (s) & \(t_{\mathrm{fail}}/\ter\) \\
\midrule
\(1.50\) & \(15.0\) kPa/m & \(24.7\) & \(140.8\) & \(140.6\) & \(5.7\) \\
\(3.33\) & \(33.3\) kPa/m & \(11.1\) & \(59.3\) & \(59.3\) & \(5.3\) \\
\(10.0\) & \(100.0\) kPa/m & \(3.7\) & \(19.6\) & \(19.2\) & \(5.2\) \\
\bottomrule
\end{tabular}
\begin{tablenotes}
\small
\item Le rapport \(t_{\mathrm{fail}}/\ter\) est quasi constant (\(5.2\)--\(5.7\)), ce qui confirme
quantitativement que le temps d'échec est piloté par le rapport \(\Delta p/L\) par l'intermédiaire du temps
caractéristique~\eqref{eq:ter} : la dépendance en la charge est correctement capturée par
\(t_{\mathrm{fail}} \propto 1/\Delta p\).
\end{tablenotes}
\end{table}

Ce résultat est important pour l'interprétation géotechnique : la cinétique de l'érosion est dominée par
le gradient hydraulique, et l'effet d'une augmentation de charge est correctement représenté par la seule
réduction du temps caractéristique. Le modèle se comporte donc, sur ce point, comme un outil de
\emph{transposition d'échelle} : les paramètres d'érodibilité du sol fixent \(\ter\), et le chargement
hydraulique fixe le rapport entre le temps d'échec et \(\ter\).

\paragraph{Une résistance aval ne peut que retarder l'instabilité.}
On pourrait s'attendre à ce qu'une forte perte singulière aval finisse par arrêter l'érosion, en captant la
quasi-totalité de la charge motrice lorsque le conduit s'élargit : c'est ce que suggère la vitesse
asymptotique~\eqref{eq:umax}, qui ne dépend plus de la géométrie et devient donc très faible pour un
\(\kout\) élevé. L'intégration couplée montre que ce mécanisme \emph{ne se réalise pas}, parce que la
perte singulière est \emph{auto-limitante} : elle varie comme \(\rho u^2\) et décroît donc en \(R^{-4}\)
lorsque le conduit s'élargit, alors que la perte répartie reste alimentée par la résistance du mélange
\(f_m\), qui sature à son plafond. Le tableau~\ref{tab:bilan-charge} décompose la charge motrice pour la
configuration la plus contraignante (\(\Delta p = 1.5\) MPa, \(\kout = 500\)).

\begin{table}[H]
\centering
\caption{Décomposition de la charge motrice au cours du temps (campagne B, \(\Delta p = 1.5\) MPa,
\(\kout = 500\)). Les deux dernières colonnes donnent la part de \(\Delta p\) consommée par la perte
singulière aval et par le frottement réparti. \(\rho_L\) désigne la masse volumique du mélange en sortie.}
\label{tab:bilan-charge}
\small
\begin{tabular}{ccccccc}
\toprule
\(t/\ter\) & \(t\) (s) & \(R_{\mathrm{out}}/R_0\) & \(Q\) (m\(^3\)/s) & \(u_{\mathrm{out}}\) (m/s)
& perte aval & frottement \\
\midrule
\(1\) & \(24.6\) & \(1.14\) & \(5.4\times10^{-5}\) & \(1.48\) & \(43.3\%\) & \(56.7\%\) \\
\(2\) & \(49.3\) & \(1.99\) & \(8.3\times10^{-5}\) & \(0.74\) & \(14.7\%\) & \(85.3\%\) \\
\(3\) & \(74.0\) & \(6.51\) & \(9.9\times10^{-5}\) & \(0.08\) & \(0.2\%\) & \(99.8\%\) \\
\(4\) & \(98.7\) & \(13.05\) & \(7.0\times10^{-4}\) & \(0.15\) & \(0.7\%\) & \(99.3\%\) \\
\(5\) & \(123.3\) & \(34.11\) & \(1.0\times10^{-2}\) & \(0.32\) & \(3.1\%\) & \(96.9\%\) \\
\(6\) & \(148.0\) & \(95.66\) & \(1.6\times10^{-1}\) & \(0.60\) & \(11.1\%\) & \(88.9\%\) \\
\bottomrule
\end{tabular}
\begin{tablenotes}
\small
\item La décomposition est établie avec la masse volumique du mélange en sortie ; le résidu du solveur
reste inférieur à \(10^{-3}\) Pa, ce qui garantit que la somme des deux parts vaut \(\Delta p\) à la
précision d'affichage. Au-delà de \(t \approx 6\ter\), l'état dépasse l'horizon de validité géométrique de
la réduction unidimensionnelle et n'est plus interprété quantitativement.
\end{tablenotes}
\end{table}

La perte aval domine donc uniquement tant que le conduit est étroit (\(R_{\mathrm{out}}/R_0 \lesssim 2\),
soit \(t \lesssim 2\ter\)). Dès que l'élargissement est engagé, elle devient marginale et le frottement du
mélange capte la quasi-totalité de la charge. La vitesse de sortie mesurée à \(t = 5\ter\) (\(0.32\) m/s)
reste très inférieure à la vitesse asymptotique~\eqref{eq:umax} (\(2.4\) m/s) : l'asymptote n'est jamais
approchée, et l'érosion \emph{ne s'arrête pas} — le calcul atteint \(R_{\mathrm{break}}\) en \(162.1\) s, en
accord avec la valeur publiée (\(161.8\) s). Une résistance aval ne peut donc que retarder l'instabilité,
jamais l'empêcher, tant que le matériau reste érodible ; le retard lui-même sature, puisqu'il vaut
\(\times1.14\) à \(3.33\) MPa et \(\times1.26\) à \(1.5\) MPa, valeurs reproduites à mieux de \(0.2\%\).

\begin{vigilance}
Deux points doivent accompagner l'interprétation de ces résultats. D'une part, le temps d'échec dépend
linéairement de la définition du critère \(R_{\mathrm{break}}\) et n'a de sens qu'accompagné de son seuil
(\(0.5\) m ici, soit \(167\,R_0\)). D'autre part, la valeur du plafond \(f_{m,\max}\) gouverne
qualitativement la réponse de la campagne B : avec la valeur \(2000\) des scénarios du code, le calcul
reproduit les temps publiés à mieux de \(0.5\%\), tandis qu'un plafond de \(5\) conduit à un arrêt de
l'érosion (\(R_{\mathrm{out}}/R_0\) plafonnant à \(5.4\) pour \(\kout = 500\)). Le plafond doit donc être
présenté comme un \emph{paramètre identifié} du modèle, et non comme un simple garde-fou numérique.
\end{vigilance}

\subsection{Influence de la perte de charge à l'entrée}
\label{sec:kin}

Le résidu hydraulique~\eqref{eq:residu} peut être complété par une perte singulière à l'entrée du
conduit, dont la forme est symétrique de celle de la sortie. Cette perte a été introduite dans le bilan de
charge~\eqref{eq:charge-entree} ; elle s'écrit, en notant \(u_0 = Q/(\pi R_0^2)\),
\begin{equation}
p(0) = P_{\mathrm{in}} - \kin\,\frac{\rho_0 u_0^2}{2},
\label{eq:cond-entree}
\end{equation}
Ce terme a une portée différente de celle du coefficient aval, car une partie en est \emph{inévitable}.
Lorsque le modèle est écrit comme un bilan de charge entre un réservoir amont et le débouché, la
conversion de la charge statique en énergie cinétique à l'entrée est un effet de premier ordre : elle
impose \(\kin = 1\) pour l'écoulement de référence, valeur retenue ici comme cas de base.

Le tableau~\ref{tab:kin} présente l'effet d'une perte d'entrée sur la dynamique de la campagne A.

\begin{table}[H]
\centering
\caption{Influence de la perte de charge à l'entrée (campagne A, \(\kout = 0\)). La colonne
\(K_{\mathrm{in,tot}}\) cumule le coefficient singulier et le terme d'énergie cinétique
(\(\mathbf{1}_{\mathrm{cin}} = 1\) si l'entrée cinétique est comptée). \(t_2\) est déduit du taux de
croissance mesuré en début de simulation.}
\label{tab:kin}
\small
\begin{tabular}{ccccc}
\toprule
\(K_{\mathrm{in}}\) & \(K_{\mathrm{in,tot}}\) & \(u_0\) (m/s) & \(\tau_{b,0}\) (Pa) & \(t_2\) (s) \\
\midrule
\(0\)   & \(0\)   & \(4.52\) & \(62.9\)  & \(668\) \quad (\(1.00\)) \\
\(0\)   & \(1\)   & \(2.51\) & \(22.5\)  & \(1937\) \quad (\(2.90\)) \\
\(0.5\) & \(1.5\) & \(2.17\) & \(17.5\)  & \(2522\) \quad (\(3.77\)) \\
\(1\)   & \(2\)   & \(1.95\) & \(14.4\)  & \(3096\) \quad (\(4.63\)) \\
\(2\)   & \(3\)   & \(1.65\) & \(10.8\)  & \(4235\) \quad (\(6.33\)) \\
\bottomrule
\end{tabular}
\begin{tablenotes}
\small
\item Le facteur entre parenthèses est le rapport \(t_2/t_2(K_{\mathrm{in,tot}} = 0)\). La contrainte
pariétale initiale vérifie \(\tau_{b,0} = R_0\left[P_{\mathrm{in}} -
K_{\mathrm{in,tot}}\rho u_0^2/2\right]/(2L)\) à mieux de \(0.3\%\) pour les cinq lignes, ce qui confirme
que le terme d'entrée est correctement pris en compte dans le résidu.
\end{tablenotes}
\end{table}

La perte d'entrée constitue un mécanisme de rétroaction \emph{négative} particulièrement efficace : en
\(K_{\mathrm{in,tot}} = 1\), c'est-à-dire en comptant simplement l'énergie cinétique de l'écoulement
d'entrée, le temps de doublement est déjà multiplié par \(2.9\). L'effet est ensuite sous-linéaire (un
facteur \(3\) sur \(K_{\mathrm{in,tot}}\) n'apporte qu'un facteur \(2.2\) sur \(t_2\)), parce que la
contrainte pariétale décroît elle-même avec la vitesse, ce qui amortit la sensibilité. Cette
auto-limitation distingue nettement l'entrée de la sortie : la perte d'entrée agit sur toute la longueur du
conduit par l'intermédiaire du débit, tandis que la perte de sortie est dominée par les dernières cellules
et devient négligeable dès que le conduit s'élargit en aval (tableau~\ref{tab:bilan-charge}).

Ce résultat a une conséquence pratique : le choix de traiter ou non l'entrée comme une singularité doit
être explicité dans la définition du problème, faute de quoi le temps caractéristique peut varier d'un
facteur \(3\) sans qu'aucun paramètre du sol n'ait changé. Le cas de référence retenu dans ce travail
correspond à \(\kout = \kin = 0\), c'est-à-dire à une pression totale imposée en entrée et à
une sortie libre ; c'est cette convention qui reproduit les valeurs publiées
(tableau~\ref{tab:validation-campagneB}).

\subsection{Profils spatiaux et fronts multiphysiques}

La figure~\ref{fig:profils-spatiaux} présente l'évolution spatio-temporelle des quatre champs du modèle
pour une configuration avec perte aval : rayon normalisé, pression relative, concentration relative et
facteur de résistance du mélange. Deux structures se dégagent. D'une part, un front advectif raide de
concentration se propage vers l'aval ; d'autre part, la localisation du gradient de pression et du
maximum de contrainte pariétale coïncide avec la position de ce front. Cette coïncidence n'est pas fortuite :
elle traduit le fait que la résistance du mélange, qui croît avec la concentration, modifie localement la
perte de charge et donc la contrainte pariétale. La localisation aval de l'élargissement du conduit est donc
une conséquence directe du couplage transport--hydraulique décrit à la section~\ref{sec:couplage}, et non
d'un effet de bord numérique. Dans la configuration de référence, cette localisation reste toutefois
modérée, ce qui est cohérent avec l'analyse de sensibilité de la loi de résistance du mélange
(section~\ref{sec:fermeture-hydraulique}) : plus le matériau est fin, plus la rétroaction rhéologique est
faible et plus le profil de rayon tend vers une forme cylindrique.

\begin{figure}[H]
\centering
\includegraphics[width=\linewidth]{figure3.png}
\caption{Profils spatiaux adimensionnels le long du conduit pour la configuration avec perte aval. De haut
en bas : rayon normalisé \(R/R_0\), rapport de pression \(p/P_{\mathrm{in}}\), concentration relative
\(\phi/\phisol\) et facteur de résistance du mélange \(f_m\). La ligne tiretée du premier panneau
indique la limite de validité géométrique \(R/L = 0.05\) de l'approximation 1D.}
\label{fig:profils-spatiaux}
\end{figure}

\subsection{Analyse des temps caractéristiques d'érosion}

Les deux campagnes convergent vers la même lecture des temps caractéristiques :
\begin{itemize}
  \item \(\ter = 2L\rho_s/(k_{er}\Delta p)\) fixe le taux de croissance de l'emballement et le temps de
        doublement \(t_2 = \ter\ln 2\) tant que la géométrie reste quasi uniforme ;
  \item le temps d'échec, rapporté à \(\ter\), est presque indépendant du chargement hydraulique
        (\(t_{\mathrm{fail}}/\ter \approx 5.3\) dans la campagne B) et dépend donc principalement du
        rapport entre le rayon d'arrêt et la géométrie initiale ;
  \item la résistance aval agit en premier lieu sur \(\ter\), par l'intermédiaire de la charge réellement
        disponible pour le frottement réparti.
\end{itemize}
À l'échelle du laboratoire, les temps de doublement obtenus (\(666\) s pour le sol érodible, \(6976\) s
pour le sol résistant) sont du même ordre que ceux rapportés dans les campagnes d'essais au trou (HET) sur
des mélanges silto-argileux compactés, ce qui confirme que le modèle restitue l'échelle temporelle de
l'érosion de conduit. Cette comparaison doit toutefois rester qualitative : elle porte sur des durées et
non sur des morphologies, puisque la géométrie du conduit échappe rapidement au domaine de validité de
l'approche 1D.

\subsection{Domaine de validité et limites de l'approche unidimensionnelle}
\label{sec:validite}

Les résultats précédents doivent être lus à la lumière du tableau~\ref{tab:sansdim}, dont trois lignes
signalent des critères non satisfaits dans la configuration de référence. Le tableau~\ref{tab:validite}
récapitule les limites et leur traduction pratique.

\begin{table}[H]
\centering
\caption{Limites de l'approche unidimensionnelle et conséquences pratiques sur l'interprétation.}
\label{tab:validite}
\small
\begin{tabular}{p{3.4cm}p{4.6cm}p{5.6cm}}
\toprule
Limite & Critère & Conséquence sur l'interprétation \\
\midrule
Élancement du conduit & \(R/L < 0.05\), soit \(R/R_0 < 1.95\) pour la campagne A &
L'horizon quantitatif est atteint dès le doublement du rayon. Au-delà, les rapports d'agrandissement et
les profils spatiaux sont des indicateurs de régime. \\[1mm]
Longueur d'établissement & \(L_e/D \sim 10\)--\(40\) pour \(L/D = 19.5\) &
La fermeture de conduite établie est appliquée sur un domaine du même ordre que son propre domaine de
validité. Une sous-estimation systématique de la perte de charge réelle est possible. \\[1mm]
Charge cinétique d'entrée & \(\frac{1}{2}\rho u^2/\Delta p = 2.1\) &
La condition \(p(0) = P_{\mathrm{in}}\) suppose le fluide déjà accéléré. Une perte singulière d'entrée
\(\kin\) de \(0.5\) réduit la vitesse de \(30\%\) et la contrainte pariétale de \(51\%\), et double le
temps de doublement (valeurs mesurées, section~\ref{sec:kin}). \\[1mm]
Inertie convective & \(8|\partial_xR|/f_D \ll 1\) &
La relation quasi-statique~\eqref{eq:quasi} cesse d'être exacte dès que le profil de rayon se creuse
significativement. \\[1mm]
Rétroaction rhéologique & \(f_m\) dépend de \((d_p/l_m)^2\) &
Pour un silt fin (\(d_p \approx 20\ \si{\micro m}\)), la résistance du mélange reste proche de l'unité et
la dynamique est gouvernée par la seule rétroaction géométrique. \\
\bottomrule
\end{tabular}
\end{table}

Il en résulte une recommandation d'interprétation qui structure l'ensemble des résultats présentés : les
\emph{temps caractéristiques} extraits du modèle (temps de doublement, temps d'échec rapporté à \(\ter\))
sont des grandeurs robustes, car ils reposent sur la relation analytique~\eqref{eq:tau-analytique} et sur
l'échelle~\eqref{eq:ter} ; les \emph{morphologies} (profils de rayon, rapports d'agrandissement atteints,
localisation du front) sont en revanche des résultats qualitatifs dès que \(R/L\) dépasse \(0.05\). Cette
distinction, énoncée dès la section~\ref{sec:hypotheses}, trouve ici sa traduction quantitative.

%=====================================================================================
\section{Synthèse}
\label{sec:synthese-chap}
%=====================================================================================

Ce chapitre a établi le cadre de modélisation et la stratégie de résolution numérique du modèle couplé.
Quatre éléments en constituent la structure.

\paragraph{Une chaîne de réductions explicite.} Le modèle unidimensionnel n'est pas posé \emph{a priori} :
il résulte d'une suite de réductions clairement identifiées, depuis la conservation de la masse
tridimensionnelle jusqu'à la relation quasi-statique~\eqref{eq:quasi}. Chaque hypothèse est associée à un
critère quantitatif (tableau~\ref{tab:sansdim}) qui définit son domaine de validité. Cette façon de
procéder a permis d'identifier trois points sur lesquels les hypothèses sont sollicitées au-delà de leur
domaine de confort : l'élancement du conduit dans la phase d'emballement, la longueur d'établissement, et
la représentation des effets d'entrée.

\paragraph{Un solveur hydraulique robuste et vérifiable.} La détermination implicite du débit
(\eqref{eq:FQ}) par encadrement et dichotomie fournit une garantie de convergence indépendante du
comportement local du résidu, ce qui est indispensable dans un problème où la fermeture rhéologique
devient raide. Le coût supplémentaire par rapport à une méthode de Newton est compensé par la fiabilité
du sous-problème hydraulique, qui conditionne l'ensemble de la chaîne de calcul.

\paragraph{Deux résultats analytiques servant de référence.} L'indépendance de la contrainte pariétale
vis-à-vis de la fermeture de frottement, \(\tau_b = R\Delta p/(2L)\), conduit à la loi exponentielle
\eqref{eq:R-exponentielle} : elle établit que l'emballement est une propriété structurelle de la condition
de pression imposée, et fournit un temps de doublement analytique reproduit par le calcul à mieux de
\(1\%\). Cette vérification est bien plus informative qu'une comparaison à une solution numérique de
référence, car elle ne dépend d'aucune discrétisation.

\paragraph{Un protocole de vérification norme par norme.} L'étude de la discrétisation du transport a
établi trois résultats quantitatifs : la condition de stabilité du schéma reconstruire--résoudre--avancer
(\(\mathrm{CFL}\le1/2\)), qui justifie le réglage retenu et fixe une limite supérieure à ne pas franchir ;
les lois asymptotiques des erreurs pour une discontinuité capturée
(\(L_1 \sim \Delta x\), \(L_2 \sim \Delta x^{1/2}\), \(L_\infty = \) constante) ; et la préservation par
construction de la borne \(0\le\phi\le\phisol\) grâce à l'intégration analytique de la source. Ces trois
éléments transforment l'étude de convergence en un diagnostic interprétable, et non en une simple
comparaison de tableaux.

Le chapitre suivant exploite ces résultats pour l'analyse des régimes d'érosion et l'étude paramétrique,
en appliquant systématiquement les cinq contrôles de vérification énumérés au
paragraphe~\ref{sec:fronts-convergence}.


%=====================================================================================
\appendix
\chapter{Rappel : du bilan de quantité de mouvement à la loi de Darcy}
\label{ann:ns}
%=====================================================================================

Cette annexe rassemble, sous une forme condensée, la chaîne d'établissement des équations utilisées comme
point de départ de la section~\ref{sec:reduction-ns}. Elle est destinée à rendre le chapitre autonome et à
expliciter les hypothèses exactes sous lesquelles les résultats de la réduction sont valables.

\section{Conservation de la masse et incompressibilité}

En écrivant que la variation de la masse d'un volume matériel est nulle, et en utilisant le théorème de
la divergence, on obtient l'équation de continuité
\begin{equation}
\frac{\partial\rho}{\partial t} + \divg\left(\rho\vec{v}\right) = 0 .
\label{eq:cont-masse}
\end{equation}
Pour un fluide dont la masse volumique ne varie pas sous l'effet des contraintes de surface (fluide
\emph{incompressible} au sens de la mécanique des milieux continus), l'équation~\eqref{eq:cont-masse} se
réduit à
\begin{equation}
\divg\vec{v} = \frac{\partial v_x}{\partial x} + \frac{\partial v_y}{\partial y}
+ \frac{\partial v_z}{\partial z} = 0 ,
\label{eq:ann-div}
\end{equation}
qui est l'équation~\eqref{eq:div-v} utilisée à la section~\ref{sec:div-v}. Il convient de noter que
l'incompressibilité n'interdit pas une variation de masse volumique d'origine thermique ou
compositionnelle ; elle interdit seulement la contribution de pression.

\section{Bilan de quantité de mouvement et équation de Cauchy}

En appliquant le principe fondamental de la dynamique à un élément fluide, les forces appliquées se
décomposent en une force de volume (le poids) et en des forces de surface, qui se décomposent elles-mêmes
en une contribution normale isotrope (la pression) et une contribution de cisaillement décrite par le
tenseur des contraintes visqueuses \(\tau_{ij}\). L'identification des forces de surface sur les faces
opposées, par développement au premier ordre, conduit à l'équation de Cauchy
\begin{equation}
\rho\frac{D\vec{v}}{Dt} = -\rho g\vec{k} - \nabla P + \nabla\cdot\tau_{ij},
\qquad
\frac{D}{Dt} = \frac{\partial}{\partial t} + \vec{v}\cdot\nabla .
\label{eq:cauchy}
\end{equation}
Le membre de gauche est la dérivée matérielle de la vitesse, qui n'est autre que l'accélération dans le
cadre eulérien : elle comprend la variation locale et la variation convective, terme responsable de la
non-linéarité intrinsèque des équations du mouvement des fluides.

\section{Fermeture constitutive newtonienne}

L'équation~\eqref{eq:cauchy} n'est pas soluble en l'état, car les contraintes n'y sont pas exprimées en
fonction des vitesses. La fermeture constitutive d'un fluide newtonien relie le tenseur des contraintes au
tenseur des taux de déformation,
\begin{equation}
\tau_{ij} = 2\mu\dot{\epsilon}_{ij},
\qquad
\dot{\epsilon}_{ij} = \frac{1}{2}\left(\frac{\partial v_i}{\partial x_j}
+ \frac{\partial v_j}{\partial x_i}\right).
\label{eq:constitutif}
\end{equation}
En reportant~\eqref{eq:constitutif} dans~\eqref{eq:cauchy} et en utilisant
\(\divg\vec{v} = 0\), les termes de dilatation disparaissent et l'on obtient l'équation de
Navier--Stokes pour un fluide incompressible newtonien,
\begin{equation}
\rho\frac{D\vec{v}}{Dt} = -\rho g\vec{k} - \nabla P + \mu\nabla^2\vec{v}.
\label{eq:ns-ann}
\end{equation}
Le dernier terme, dit visqueux, est la divergence du tenseur des contraintes, dont le caractère
diffusif apparaît explicitement sous la forme du laplacien de la vitesse.

\section{Pression réduite : élimination de la gravité}

En retranchant l'équation de l'hydrostatique, \(\nabla P_s = -\rho g\vec{k}\), de
l'équation~\eqref{eq:ns-ann}, la gravité disparaît au profit de la \emph{pression réduite}
\(\Chat = P - P_s = P + \rho g z\) :
\begin{equation}
\rho\frac{D\vec{v}}{Dt} = -\nabla\Chat + \mu\nabla^2\vec{v},
\qquad
\Chat = P + \rho g z .
\label{eq:ns-reduite}
\end{equation}
Physiquement, seule la pression en excès de l'hydrostatique met le fluide en mouvement : la pression
hydrostatique assure, à elle seule, la compensation exacte du poids. Cette réduction est celle utilisée au
paragraphe~\ref{sec:degenerescence} pour éliminer explicitement le terme de gravité du modèle 1D.

\section{Dégénérescence vers les solutions laminaires : Poiseuille et Darcy}

\subsection{Écoulement rampant et terminaison du bilan de forces}

Le rapport entre les forces d'inertie et les forces visqueuses définit le nombre de Reynolds. Pour
\(Re \ll 1\) — régime rampant de Stokes — le terme inertiel \(\rho D\vec{v}/Dt\) devient négligeable devant
le terme visqueux, et l'équation~\eqref{eq:ns-reduite} se réduit à l'équilibre
\begin{equation}
0 = -\nabla\Chat + \mu\nabla^2\vec{v}.
\label{eq:stokes}
\end{equation}
Cette équation est encore une équation aux dérivées partielles, mais linéaire : elle constitue le point de
départ des solutions analytiques ci-dessous, et le prototype de toutes les fermetures locales utilisées
dans les milieux poreux.

\subsection{Écoulement dans une fracture plane}

Pour un écoulement établi entre deux plans parallèles distants de \(W\), la vitesse ne dépend que de la
coordonnée transverse et l'équation~\eqref{eq:stokes} s'intègre en
\begin{equation}
\vx(z) = \frac{1}{2\mu}\frac{\mathrm{d}\Chat}{\mathrm{d}x}\left(z^2 - Wz\right),
\qquad
q = -\frac{W^2}{12\mu}\frac{\mathrm{d}\Chat}{\mathrm{d}x}.
\label{eq:fracture}
\end{equation}

\subsection{Écoulement dans un tube circulaire}

Pour un écoulement établi dans un tube de rayon \(R\), avec la condition de non-glissement à la paroi, la
solution parabolique (dite de Poiseuille) s'écrit
\begin{equation}
\vx(r) = \frac{1}{4\mu}\frac{\mathrm{d}\Chat}{\mathrm{d}x}\left(r^2-R^2\right),
\qquad
q = -\frac{R^2}{8\mu}\frac{\mathrm{d}\Chat}{\mathrm{d}x} .
\label{eq:poiseuille-ann}
\end{equation}
En introduisant la charge \(h = \Chat/(\rho g)\), cette expression prend exactement la forme de la loi de
Darcy \(q = -K\,\mathrm{d}h/\mathrm{d}s\) avec \(K = \rho g R^2/(8\mu)\).

\subsection{Portée et limites de l'analogie avec le modèle 1D}

La comparaison des trois expressions (fracture, tube, Darcy) met en évidence une structure commune :
\begin{equation}
\underbrace{q = -\frac{w_c^2}{c\,\mu}\frac{\mathrm{d}\Chat}{\mathrm{d}x}}_{\text{solution locale établie}}
\qquad\longleftrightarrow\qquad
\underbrace{\tau_b = \frac{f_D}{8}\rho\,f_m\,\frac{u^2}{2}}_{\text{fermeture turbulente retenue}}
\label{eq:analogie}
\end{equation}
Dans les deux cas, la description réduite d'un écoulement confiné repose entièrement sur une
\emph{fermeture locale} reliant un flux à une contrainte. Ce qui change entre les deux membres est la
nature de cette fermeture, linéaire dans le cas rampant et quadratique dans le cas turbulent. L'analogie
avec les écoulements souterrains est donc structurelle et pédagogique ; elle n'est pas quantitative, et le
modèle de piping ne saurait être interprété comme un modèle d'écoulement dans un milieu poreux.
Le tableau~\ref{tab:sansdim} quantifie cet écart : le conduit de renard hydraulique est en régime
turbulent établi, hors du domaine des solutions de Poiseuille.

%=====================================================================================
\section*{Liste des notations}
\addcontentsline{toc}{section}{Liste des notations}
%=====================================================================================

\begin{tabbing}
\hspace{3.6cm}\=\kill
\(A\) \> section transversale du conduit \\
\(\Chat\) \> pression réduite (ou motrice), \(\Chat = P + \rho g z\) \\
\(C_B\) \> coefficient dispersif de la loi de résistance du mélange \\
\(C_r\) \> coefficient de longueur de mélange, \(l_m = C_r R\) \\
\(d_p\) \> diamètre caractéristique des particules \\
\(f_D\) \> facteur de frottement de Darcy de l'écoulement de référence \\
\(f_m\) \> facteur multiplicatif de résistance dû aux particules \\
\(F(Q)\) \> résidu scalaire du problème hydraulique, \(F(Q)=p_{\mathrm{calc}}(L;Q)-p_{\mathrm{out}}(Q)\) \\
\(F_{i+1/2}\) \> flux numérique de concentration à l'interface \(i+1/2\) \\
\(k\) \> taux de relaxation local de la concentration, \(k = 2\mdot/(R\rho_s)\) \\
\(k_{er}\) \> coefficient d'érosion \\
\(\kin,\ \kout\) \> coefficients de perte singulière à l'entrée et à la sortie \\
\(L\) \> longueur du conduit \\
\(L_e\) \> longueur d'établissement de l'écoulement \\
\(l_m\) \> longueur de mélange \\
\(\mdot\) \> flux massique érodé par unité de surface pariétale \\
\(N_x\) \> nombre de cellules du maillage \\
\(n_\lambda\) \> exposant de concentration de la loi de résistance \\
\(P_{\mathrm{in}},\ P_{\mathrm{out}}\) \> pressions imposées à l'entrée et à la sortie \\
\(p\) \> pression moyenne de section \\
\(Q\) \> débit volumique (constant le long du conduit) \\
\(R\) \> rayon effectif du conduit \\
\(R_c\) \> rayon critique de la loi d'évolution, \(R_c = 2L\tau_c/\Delta p\) \\
\(R_0\) \> rayon initial \\
\(Re\) \> nombre de Reynolds \\
\(r_i\) \> rapport de pentes au nœud \(i\), \(r_i = \Delta_{L,i}/\Delta_{R,i}\) \\
\(s_i\) \> pente reconstruite dans la cellule \(i\) \\
\(\mathcal{S}^n\) \> état du système au pas \(n\), \(\mathcal{S}^n=\{R_i^n,\phi_i^n\}\) \\
\(S_A\) \> source volumique d'érosion par unité de longueur \\
\(t\) \> temps \\
\(\ter\) \> temps caractéristique d'érosion, \(\ter = 2L\rho_s/(k_{er}\Delta p)\) \\
\(t_2\) \> temps de doublement du rayon \\
\(\dt\) \> pas de temps \\
\(u\) \> vitesse moyenne de section \\
\(u_p = \beta_u u\) \> vitesse de transport des particules \\
\(\vec{v}\) \> champ de vitesse tridimensionnel \\
\(\beta_u\) \> rapport entre vitesse des particules et vitesse du fluide \\
\(\gamma\) \> taux de croissance du rayon en géométrie uniforme \\
\(\Delta p\) \> perte de charge imposée, \(\Delta p = P_{\mathrm{in}}-P_{\mathrm{out}}\) \\
\(\dot{\epsilon}_{ij}\) \> tenseur des taux de déformation \\
\(\lambda(\phi)\) \> fonction de concentration de la loi de résistance du mélange \\
\(\mu_w\) \> viscosité dynamique de l'eau \\
\(\Pi_{\mathrm{aval}}\) \> rapport entre perte singulière aval et perte répartie \\
\(\rho,\rho_w,\rho_p\) \> masses volumiques du mélange, de l'eau et des particules \\
\(\rho_s\) \> masse volumique apparente saturée du sol intact \\
\(\tau_b\) \> contrainte de cisaillement à la paroi \\
\(\tau_c\) \> contrainte critique d'érosion \\
\(\tau_{ij}\) \> tenseur des contraintes visqueuses \\
\(\tau_h,\ \tau_a\) \> échelles de temps convective et acoustique \\
\(\phi\) \> fraction volumique de solides en suspension \\
\(\phisol\) \> fraction solide du sol intact \\
\(\Psi(r)\) \> fonction limiteur de pente \\
\(\Disc\) \> coefficient de diffusion numérique \\
\end{tabbing}

%=====================================================================================
\begin{thebibliography}{99}
\addcontentsline{toc}{section}{Références}
%=====================================================================================

\bibitem{Barenblatt2003}
Barenblatt, G.I. (2003) \emph{Scaling}. Cambridge: Cambridge University Press.

\bibitem{Bonelli2006}
Bonelli, S., Brivois, O., Borghi, R. and Benahmed, N. (2006) `On the modelling of piping erosion',
\emph{Comptes Rendus Mécanique}, 334(8--9), pp. 555--559.

\bibitem{Colebrook1939}
Colebrook, C.F. (1939) `Turbulent flow in pipes, with particular reference to the transition region
between the smooth and rough pipe laws', \emph{Journal of the Institution of Civil Engineers}, 11,
pp. 133--156.

\bibitem{Fell2007}
Fell, R. and Fry, J.-J. (eds.) (2007) \emph{Internal erosion of dams and their foundations}. Boca Raton,
FL: Taylor \& Francis.

\bibitem{Julien2012}
Julien, P.Y. (2012) \emph{Erosion and sedimentation}. 2nd edn. Cambridge: Cambridge University Press.

\bibitem{Lachouette2008}
Lachouette, D., Golay, F. and Bonelli, S. (2008) `One-dimensional modeling of piping flow erosion',
\emph{Comptes Rendus Mécanique}, 336(9), pp. 731--738.

\bibitem{Lachouette2009}
Lachouette, D. (2009) \emph{Érosion par renard hydraulique : modélisation et comparaison aux essais
expérimentaux}. Thèse de doctorat. Université de Grenoble.

\bibitem{LeVeque2002}
LeVeque, R.J. (2002) \emph{Finite volume methods for hyperbolic problems}. Cambridge: Cambridge
University Press.

\bibitem{Moody1944}
Moody, L.F. (1944) `Friction factors for pipe flow', \emph{Transactions of the ASME}, 66, pp. 671--684.

\bibitem{Partheniades1965}
Partheniades, E. (1965) `Erosion and deposition of cohesive soils', \emph{Journal of the Hydraulics
Division}, 91(1), pp. 105--139.

\bibitem{Sweby1984}
Sweby, P.K. (1984) `High resolution schemes using flux limiters for hyperbolic conservation laws',
\emph{SIAM Journal on Numerical Analysis}, 21(5), pp. 995--1011.

\bibitem{Toro2009}
Toro, E.F. (2009) \emph{Riemann solvers and numerical methods for fluid dynamics}. 3rd edn. Berlin:
Springer.

\bibitem{Wahl2014}
Wahl, T.L. (2014) \emph{Evaluation of erodibility-based embankment dam breach equations}. Report
No.\ HL-2014-02. Denver, CO: U.S. Department of the Interior, Bureau of Reclamation.

\bibitem{WanFell2002}
Wan, C.F. and Fell, R. (2002) \emph{Investigation of erosion rate of base soils in embankment dams}.
Report R-412. Sydney: University of New South Wales.

\bibitem{WanFell2004}
Wan, C.F. and Fell, R. (2004) `Investigation of rate of erosion of soils in embankment dams',
\emph{Journal of Geotechnical and Geoenvironmental Engineering}, 130(4), pp. 373--380.

\end{thebibliography}

\end{document}
